\documentclass[11pt]{article}
\usepackage[letterpaper,portrait,margin=1in]{geometry}
\usepackage[T1]{fontenc}
\usepackage{lmodern}
\usepackage{amsmath,amssymb,amsthm}
\usepackage{microtype}
\usepackage{xcolor}
\usepackage{tikz}
\usetikzlibrary{arrows.meta,calc,fit,positioning}
\usepackage{aliascnt}
\usepackage[colorlinks=true,linkcolor=teal,citecolor=teal,urlcolor=teal]{hyperref}
\usepackage[nameinlink,capitalize]{cleveref}

\allowdisplaybreaks

\newtheorem{theorem}{Theorem}[section]
\newaliascnt{lemma}{theorem}
\newtheorem{lemma}[lemma]{Lemma}
\aliascntresetthe{lemma}
\newaliascnt{proposition}{theorem}
\newtheorem{proposition}[proposition]{Proposition}
\aliascntresetthe{proposition}
\newaliascnt{corollary}{theorem}
\newtheorem{corollary}[corollary]{Corollary}
\aliascntresetthe{corollary}
\newtheorem{fact}{Fact}[section]
\renewcommand{\thefact}{\thesection.\Roman{fact}}

\crefname{theorem}{Theorem}{Theorems}
\Crefname{theorem}{Theorem}{Theorems}
\crefname{lemma}{Lemma}{Lemmas}
\Crefname{lemma}{Lemma}{Lemmas}
\crefname{proposition}{Proposition}{Propositions}
\Crefname{proposition}{Proposition}{Propositions}
\crefname{corollary}{Corollary}{Corollaries}
\Crefname{corollary}{Corollary}{Corollaries}
\crefname{fact}{Fact}{Facts}
\Crefname{fact}{Fact}{Facts}
\crefname{equation}{equation}{equations}
\Crefname{equation}{Equation}{Equations}

\newcommand{\C}{\mathbb C}
\newcommand{\R}{\mathbb R}
\newcommand{\cA}{\mathcal A}
\newcommand{\cC}{\mathcal C}
\newcommand{\cE}{\mathcal E}
\newcommand{\cF}{\mathcal F}
\newcommand{\cH}{\mathcal H}
\newcommand{\cM}{\mathcal M}
\newcommand{\cS}{\mathcal S}
\newcommand{\cT}{\mathcal T}
\newcommand{\U}{\mathrm U}
\newcommand{\dd}{\mathrm d}
\newcommand{\diag}{\operatorname{diag}}
\newcommand{\Tr}{\operatorname{Tr}}
\newcommand{\eps}{\varepsilon}
\newcommand{\ket}[1]{|#1\rangle}
\newcommand{\proj}[1]{|#1\rangle\!\langle #1|}
\newcommand{\ot}{\otimes}
\newcommand{\Forb}{F_{\mathrm{orb}}}
\newcommand{\Fcl}{F_{\mathrm{cl}}}
\newcommand{\Funiv}{F_{\mathrm{univ}}}
\newcommand{\doi}[2]{\href{https://doi.org/#1}{#2}}
\newcommand{\arxiv}[1]{\href{https://arxiv.org/abs/#1}{\texttt{arXiv:#1}}}

\title{Asymptotically Optimal Mixed-State Cloning}
\author{Jiani Fei\(^{*}\)\qquad
  Jinzhao Wang\(^{*,\dagger}\)}
\date{}
\begin{document}
\maketitle
\begingroup
\renewcommand{\thefootnote}{\fnsymbol{footnote}}
\footnotetext[1]{\textit{Leinweber Institute for Theoretical Physics,
  Stanford University, Stanford, CA 94305, USA}}
\footnotetext[2]{\textit{Zyphra Technology, San Francisco, CA 94105, USA.}
  \enspace Email:
  \href{mailto:wang.jinzhao226@gmail.com}{\texttt{wang.jinzhao226@gmail.com}}}
\endgroup

\begin{abstract}
Unknown quantum states can only be cloned approximately.  Although the
optimal pure-state cloner has long been known, the mixed-state problem remains
challenging.  We give an asymptotically optimal mixed-state cloning protocol,
evaluated by global fidelity.  Let \(d\ge2\) and
\(m_n/n\to\gamma>1\).  We determine the asymptotically
optimal global \(n\to m_n\) cloning root fidelity for full-rank
\(d\)-dimensional states with a non-degenerate spectrum.  For a known spectrum
\(p_1>\cdots>p_d>0\), the minimax fidelity on its unitary orbit is
\[
  \Forb(\gamma,p)
  =\prod_{i<j}
  \frac{\sqrt{\gamma}+\sqrt{q_{ij}(\gamma-1+q_{ij})}}
       {\gamma+q_{ij}},
  \qquad q_{ij}=\frac{p_j}{p_i}.
\]
If the spectrum is unknown, a single spectrum-independent cloner guarantees
the pointwise benchmark
\[
  \Funiv(\gamma,p)
  =\left(\frac{2\sqrt\gamma}{1+\gamma}\right)^{(d-1)/2}
   \Forb(\gamma,p).
\]
On every regular compact spectral set \(K\), the global minimax value is
\(\inf_{p\in K}\Funiv(\gamma,p)\).
The cloner is constructed by combining Schur--Weyl label processing with a Cartan-sector covariant channel.  For the converse, two-way local
asymptotic normality transfers any cloner to classical Gaussian dilation and
quantum Gaussian amplification problems, whose optimal bounds give the
matching converses.
\end{abstract}

\tableofcontents

\section{Introduction}\label{sec:introduction-results}

The no-cloning theorem rules out perfect copying of an arbitrary unknown
quantum state \cite{WoottersZurek1982,Dieks1982}.  Approximate cloning asks
instead how well a quantum channel can transform \(n\) identical inputs into
an output resembling \(m\) ideal copies.  For pure states, early universal
qubit machines were followed by exact solutions for the standard global and
single-clone criteria
\cite{BuzekHillery1996,GisinMassar1997,
BrussDiVincenzoEkertFuchsMacchiavelloSmolin1998}: Werner solved the global
universal problem in arbitrary dimension \cite{Werner1998}, while
Keyl--Werner established optimality when individual clones are tested
\cite{KeylWerner1999}.  State-dependent global cloning was developed in
\cite{BrussDiVincenzoEkertFuchsMacchiavelloSmolin1998,
CheflesBarnett1999}; see also the reviews
\cite{ScaraniIblisdirGisinAcin2005,FanWangJingYueShiZhangMu2014}.

The choice of figure of merit is essential.  At the level of individual
clones, the many-output limit is closely tied to state estimation
\cite{MassarPopescu1995,BrussEkertMacchiavello1998,BaeAcin2006,
ChiribellaDAriano2006}.  When all outputs are examined jointly, the relation
is subtler: global asymptotic cloning has been analyzed for broad pure-state
families, and an optimal measure-and-prepare output need not consist of
independent copies of the estimated state
\cite{ChiribellaYang2014Optimal}.  The dependence of the optimum on a joint,
marginal, or process-level criterion is already visible for coherent-state
and phase-covariant families
\cite{CerfKrugerNavezWernerWolf2005,KimChitambar2022}.  Here we use root
fidelity with the entire product target \(\rho^{\ot m}\).  This global
criterion is stronger than reproducing one-site marginals and is therefore
distinct from broadcasting, which permits correlations among the outputs
\cite{BarnumCavesFuchsJozsaSchumacher1996}.

Mixed states do not
evade the exact no-cloning obstruction: the general
no-broadcasting theorem rules out a single channel that clones every member
of a noncommuting family of density matrices, with an entropic proof given by
Kalev--Hen
\cite{BarnumCavesFuchsJozsaSchumacher1996,KalevHen2008}.  Approximate
mixed-state cloning poses an additional structural difficulty: both the
eigenvalues and the eigenbasis may be unknown, and the two kinds of
information behave differently under repeated sampling.  In contrast to the
pure-state case, earlier mixed-state results either did
not identify an optimal cloner with a matching value or applied only to
special state families or different criteria.  Thus no sharp general solution
was known under the global product-fidelity criterion.  Work directly on
mixed-state cloning derived global-fidelity, relative-error, and
state-dependent bounds for prescribed finite ensembles
\cite{Rastegin2003Upper,Rastegin2002Relative,
Rastegin2003StateDependent}; constructed cloners for symmetric-subspace inputs
and identical mixed qubits under one-site shrinking-factor criteria
\cite{Fan2003,FanLiuShi2007}; studied a restricted mixed-qubit construction
under an averaged Hilbert--Schmidt criterion \cite{Kazakov2007}; and established
general entropic limitations on simultaneous approximate cloning and
broadcasting of mixed states \cite{LemmWilde2017}.  Closely related work solved
optimal mixed-qubit purification and a local-fidelity cloning/estimation
problem, rather than the joint product-fidelity task considered here
\cite{CiracEkertMacchiavello1999}.  Multi-input mixed-state work has also
focused on broadcasting and superbroadcasting, which optimize marginals
rather than the joint product target
\cite{DArianoMacchiavelloPerinotti2005,
BuscemiDArianoMacchiavelloPerinotti2005Optimal,
BuscemiDArianoMacchiavelloPerinotti2006,DangFan2007}.  More recently, the
random-purification channel of Tang--Wright--Zhandry was combined with
pure-state cloning to give a general all-output guarantee, rather than an
exact optimum, for rank-constrained mixed states
\cite{TangWrightZhandry2025,LiTheilHarrowChuang2026}.

In general, mix-state cloning problem admits two models.  In the known-spectrum problem, the
eigenvalues are supplied and the state ranges over one unitary orbit; cloning
must reproduce only the unknown eigenbasis.  In the unknown-spectrum problem,
a single spectrum-independent sequence of channels must handle both the
eigenvalues and the eigenbasis, uniformly over compact subsets of the strict
spectral chamber.
The second problem is not merely the first preceded by spectrum estimation:
the measured spectrum retains fluctuations of order \(n^{-1/2}\), which must
themselves be converted from the \(n\)-copy to the \(m_n\)-copy scale.

Here we resolve the mix-state cloning problem in the proportional-output asymptotic
regime.  For every fixed dimension and \(m_n/n\to\gamma>1\), we determine the
optimal global cloning fidelity and construct asymptotically optimal cloners both on each
full-rank, nondegenerate unitary orbit and, when the spectrum is unknown,
uniformly over regular compact spectral sets.

Two complementary decompositions drive the solution.  Schur--Weyl duality
separates the spectral information, carried by the Young label, from the
eigenbasis information, carried by the conditional irreducible state
\cite{KeylWerner2001Spectrum}.  Our cloners process the label classically and
apply a multiplicity-one Cartan-sector channel to the irreducible state
\cite{ParthasarathyRangaRaoVaradarajan1967,FultonHarris1991}.
Two-way local asymptotic normality (LAN) gives the corresponding asymptotic
separation into a classical Gaussian shift and displaced thermal oscillator
modes
\cite{GutaJencova2007LAN,GutaKahn2006LANQubits,
KahnGuta2009LANFiniteDimensional}.
It transfers any putative improvement to a Gaussian dilation or amplification
problem.  Quantum amplifier noise limits
\cite{Caves1982LinearAmplifiers} and the least-noise mixed-Gaussian cloner
\cite{GutaMatsumoto2006MixedGaussianCloning} identify the quantum benchmark;
a unified flat-prior fidelity argument gives the classical and quantum
converses used here.

\subsection{Results}\label{sec:main-results}

Fix positive integers \(m_n\) such that
\[
  \gamma_n:=\frac{m_n}{n}\longrightarrow\gamma>1,
  \qquad r_d:=\binom d2.
\]
In particular, \(m_n>n\) for all sufficiently large \(n\).
For the finitely many indices with \(m_n\le n\), the two cloners constructed
below are both defined by the marginal channel
\[
  \cC_n^p(X)=\cC_n^{\mathrm{univ}}(X)
  :=\Tr_{m_n+1,\ldots,n}(X),
\]
with the empty partial trace interpreted as the identity when \(m_n=n\).
This channel is spectrum independent, \(\U(d)\)-covariant, and sends
\(\rho^{\ot n}\) exactly to \(\rho^{\ot m_n}\).  Thus all nontrivial
construction formulas and asymptotic proofs below concern \(m_n>n\).
Write \(\mathsf A_d:=\{x\in\R^d:\sum_i x_i=1\}\).  Let
\[
  \Delta_d^\circ
  :=\left\{p\in\mathsf A_d:p_1>\cdots>p_d>0\right\},\qquad
  \rho_p=\diag(p_1,\ldots,p_d),\qquad
  \rho_{p,g}=g\rho_p g^*.
\]
For positive trace-class operators \(A,C\), not necessarily normalized, we
write
\[
  F(A,C)=\|\sqrt A\sqrt C\|_1.
\]
For states this is the unsquared Uhlmann fidelity.  For nonnegative functions
or measures, the same notation denotes their Hellinger affinity; in
particular, \(F(P,Q)=\sum_x\sqrt{P(x)Q(x)}\) on a countable space.  The target
in every statement below is the entire product state
\(\rho_{p,g}^{\ot m_n}\).  For
\(p\in\Delta_d^\circ\), define
\[
  q_{ij}=\frac{p_j}{p_i},
  \qquad
  f_\gamma(q)
  =\frac{\sqrt\gamma+\sqrt{q(\gamma-1+q)}}{\gamma+q},
\]
and
\begin{equation}\label{eq:main-values}
  \Forb(\gamma,p)=\prod_{i<j}f_\gamma(q_{ij}),
  \qquad
  \Fcl(\gamma,d)
  =\left(\frac{2\sqrt\gamma}{1+\gamma}\right)^{(d-1)/2},
  \qquad
  \Funiv(\gamma,p)=\Fcl(\gamma,d)\Forb(\gamma,p).
\end{equation}
The factor \(\Forb\) measures the cost of amplifying the unknown eigenbasis,
whereas \(\Fcl\) measures the additional cost of dilating the unknown
spectral fluctuation.

\paragraph{Known spectrum}
Fixing \(p\) isolates the genuinely noncommutative part of the problem.  The
target distribution of Young labels is then known and can be sampled exactly
at size \(m_n\); only the conditional representation sectors, which encode
the eigenbasis, need to be cloned.

\begin{theorem}[Known-spectrum optimum]\label{thm:known-optimum}
For every \(p\in\Delta_d^\circ\),
\begin{equation}\label{eq:known-optimum}
  \lim_{n\to\infty}
  \sup_{\cM_n}\inf_{g\in\U(d)}
  F\!\left(\cM_n(\rho_{p,g}^{\ot n}),\rho_{p,g}^{\ot m_n}\right)
  =\Forb(\gamma,p),
\end{equation}
where the supremum is over all quantum channels from \(n\) to \(m_n\)
qudits.  The explicit sequence of channels \(\cC_n^p\) defined in
\cref{sec:known-label-kernel} asymptotically attains the displayed value.
\end{theorem}

Thus every pair \(i<j\) contributes one displaced-thermal amplification
factor \(f_\gamma(p_j/p_i)\).  The Schur--Cartan cloner attains their product,
while fixed-spectrum LAN shows that any better finite-dimensional cloner
would violate the Gaussian amplification optimum proved here.

\paragraph{Unknown spectrum}
If \(p\) is not supplied, the target Young-label distribution can no longer
be sampled directly.  The measured input label determines the spectrum only
at the \(n^{-1/2}\) scale, and converting \(n\) copies into
\(m_n\sim\gamma n\) copies requires dilating this classical fluctuation.
This produces the extra, spectrum-independent factor \(\Fcl(\gamma,d)\).

\begin{theorem}[Unknown-spectrum achievability and compact-set optimum]
\label{thm:unknown-optimum}
There exists a sequence of spectrum-independent, \(\U(d)\)-covariant channels
\(\cC_n^{\mathrm{univ}}\) with the following properties.
\begin{enumerate}
\item[(a)] For every compact \(K\Subset\Delta_d^\circ\),
\begin{equation}\label{eq:unknown-uniform-achievability}
  \liminf_{n\to\infty}
  \inf_{\substack{p\in K\\g\in\U(d)}}
  \left[
  F\!\left(
    \cC_n^{\mathrm{univ}}(\rho_{p,g}^{\ot n}),
    \rho_{p,g}^{\ot m_n}
  \right)-\Funiv(\gamma,p)
  \right]\ge0.
\end{equation}

\item[(b)] Let \(K\Subset\Delta_d^\circ\) be nonempty and
  satisfy
  \(K=\overline{\operatorname{int}_{\mathsf A_d}K}^{\,\mathsf A_d}\).
Then
\begin{equation}\label{eq:compact-global-value}
  \lim_{n\to\infty}
  \sup_{\cM_n}
  \inf_{\substack{p\in K\\g\in\U(d)}}
  F\!\left(\cM_n(\rho_{p,g}^{\ot n}),\rho_{p,g}^{\ot m_n}\right)
  =\inf_{p\in K}\Funiv(\gamma,p),
\end{equation}
where the supremum is over all quantum channels from \(n\) to \(m_n\)
qudits.  The same sequence \(\cC_n^{\mathrm{univ}}\), defined in
\cref{sec:unknown-label-kernel}, asymptotically attains the displayed value.
\end{enumerate}
\end{theorem}

The common architecture of \(\cC_n^p\) and \(\cC_n^{\mathrm{univ}}\) is
summarized in \cref{fig:schur-cartan-cloner}.  After the Schur transform, the
Young label \(\mu\) carries the spectral information, while the conditional
irrep state carries the unknown eigenbasis.  A classical kernel chooses an
output label \(\nu\), and the common covariant sector channel maps the conditional
state into the \(\nu\)-sector before the \(m_n\)-system output is reconstructed.
For known \(p\), the label kernel samples the exact target law \(P_{m_n,p}\)
independently of \(\mu\); the spectrum-independent kernel instead randomly
rounds \(\gamma_n\mu\).  Their label factors are therefore \(1\) and \(\Fcl\),
while the common sector factor is \(\Forb\), giving the totals \(\Forb\) and
\(\Funiv\).  The details are developed in \cref{sec:label-kernels}.

\begin{figure}[!tbp]
\centering
\resizebox{\linewidth}{!}{%
\begin{tikzpicture}[
  flow/.style={-{Latex[length=2mm]},thick},
  dependency/.style={-{Latex[length=1.8mm]},thin,gray},
  guide/.style={thin,gray},
  box/.style={draw,rounded corners=2pt,align=center,inner sep=4pt,
    minimum height=9mm,font=\small},
  labelbox/.style={box,fill=blue!6},
  sectorbox/.style={box,fill=orange!9},
  note/.style={draw,rounded corners=2pt,align=left,inner sep=4pt,
    minimum height=12mm,fill=gray!6,font=\scriptsize}
]
  \node[box,fill=gray!8,text width=17mm] (input) at (-0.5,0)
    {input\\[-1pt]\(\rho_{p,g}^{\ot n}\)};

  \node[labelbox,text width=24mm] (labelin) at (3.5,1.15)
    {Young label\\\(\mu\sim P_{n,p}\)};
  \node[sectorbox,text width=24mm] (sectorin) at (3.5,-1.15)
    {irrep state\\\(\rho_{p,g,\mu}\in\cH_\mu\)};

  \node[labelbox,text width=27mm] (labelout) at (7.1,1.15)
    {label kernel \(K_n\)\\\(\nu\in\mathsf Y_{m_n}\)};
  \node[sectorbox,text width=27mm] (sectorout) at (7.1,-1.15)
    {sector channel\\\(\cE_{\mu,\nu}(\rho_{p,g,\mu})\)};

  \node[box,fill=green!8,text width=28mm] (blockout) at (10.8,0)
    {assemble the \(\nu\)-block\\[-1pt]
     \(\cE_{\mu,\nu}(\rho_{p,g,\mu})
       \ot I_{\cS_\nu}/\dim\cS_\nu\)};
  \node[box,fill=gray!8,text width=22mm,font=\scriptsize]
    (output) at (15.55,0)
    {output state\\[-1pt]
     \(\mathcal Q_n[K_n](\rho_{p,g}^{\ot n})\)};

  \coordinate (schursplit) at (2.0,0);
  \draw[flow] (schursplit) |- (labelin.west);
  \draw[flow] (schursplit) |- (sectorin.west);
  \draw[flow] (input.east) --
    node[below=2pt,pos=.45,font=\scriptsize,align=center]
      {Schur\\transform} (schursplit);
  \draw[flow] (labelin.east) -- (labelout.west);
  \draw[flow] (sectorin.east) -- (sectorout.west);
  \draw[dependency] (labelout.south) --
    node[right,font=\scriptsize,text=gray] {selects \(\nu\)} (sectorout.north);
  \coordinate (assemblemerge) at (8.9,0);
  \draw[thick] (labelout.east) -| (assemblemerge);
  \draw[thick] (sectorout.east) -| (assemblemerge);
  \draw[flow] (assemblemerge) -- (blockout.west);
  \draw[flow] (blockout.east) --
    node[below=2pt,font=\scriptsize,align=center]
      {inverse Schur\\transform} (output.west);

  \node[note,text width=47mm] (known) at (4.45,3.0)
    {\textbf{Known \(p\):}
     \(K_n^p(\nu\mid\mu)=P_{m_n,p}(\nu)\).};
  \node[note,text width=58mm] (unknown) at (10.25,3.0)
    {\textbf{Unknown \(p\):}
     \(K_n^{\mathrm{rnd}}\) randomly rounds \(\gamma_n\mu\).};
  \coordinate (kernelchoice) at ([yshift=4mm]labelout.north);
  \draw[guide] (known.south) -- (kernelchoice);
  \draw[guide] (unknown.south) -- (kernelchoice);
  \draw[dependency] (kernelchoice) -- (labelout.north);

\end{tikzpicture}%
}
\caption{\textbf{Cloner Construction}.
The Schur transform separates the classical Young label from the conditional
irrep state.  A Markov kernel selects the output label, while the generic
sector channel \(\cE_{\mu,\nu}\) maps the conditional state into that sector;
a maximally mixed Specht factor is then restored. Here
\(\mathcal Q_n[K_n]\) denotes the lifted channel associated with the selected
kernel. For known spectrum the
target label is sampled independently of \(\mu\), whereas for unknown
spectrum a randomized dilation is used; these choices produce our cloners
for \(m_n>n\), namely
\(\cC_n^p=\mathcal Q_n[K_n^p]\) and
\(\cC_n^{\mathrm{univ}}=\mathcal Q_n[K_n^\mathrm{rnd}]\), respectively.
On the label pairs
carrying asymptotically all probability, the sector fidelity tends to
\(\Forb(\gamma,p)\), giving the displayed benchmarks.  }
\label{fig:schur-cartan-cloner}
\end{figure}


The minimax meaning of \eqref{eq:compact-global-value} is as follows.
The set \(K\) may be known when the channel is designed, but the actual
spectrum \(p\in K\) and eigenbasis \(g\) are not.  The order of optimization
is important: the supremum first allows a channel sequence tailored to
\(K\), and the inner infimum then tests that single sequence against every
\(p\in K\) and \(g\in\U(d)\).  Since the unknown-spectrum benchmark
\(\Funiv(\gamma,p)\) varies with \(p\), the resulting worst-case scalar is
governed by its least favorable value, hence the infimum over \(p\) on the
right-hand side.  No further infimum over \(g\) appears there because
unitary covariance makes the benchmark constant along the orbit of each
fixed spectrum.  Thus, even with knowledge of \(K\), no \(K\)-adapted
sequence of channels can asymptotically beat
\(\inf_{p\in K}\Funiv(\gamma,p)\).  The single construction
\(\cC_n^{\mathrm{univ}}\), which does not depend on \(K\), guarantees
\(\Funiv(\gamma,p)\) pointwise, uniformly for \(p\in K\), and hence
asymptotically attains this worst-case value.

The condition
\(K=\overline{\operatorname{int}_{\mathsf A_d}K}^{\,\mathsf A_d}\)
is a clean sufficient regularity assumption for full-dimensional spectral
uncertainty.  Its interior contains small neighborhoods in all \(d-1\)
spectral directions, and the closure condition allows the infimum over
interior points in the converse to be replaced, by continuity of \(\Funiv\),
with the infimum over all of \(K\).  We do not claim that this condition is
necessary, but some thickness assumption cannot be removed from a
geometry-independent statement of this form.  For a singleton \(K=\{p\}\),
the channel may be tailored to \(p\), so the problem has the known-spectrum
value \(\Forb(\gamma,p)\), not \(\Funiv(\gamma,p)\).  More generally,
lower-dimensional spectral classes encode partial spectral knowledge and
may have different local Gaussian penalties.  Thus
\eqref{eq:compact-global-value} is the broadest geometry-independent
full-dimensional converse proved here; no statement with the same
right-hand side can extend to all compact \(K\).

For example, a natural admissible class keeps the smallest eigenvalue and every adjacent
spectral gap uniformly positive.  For \(0<\kappa<2/[d(d+1)]\),\footnote{At \(\kappa=2/[d(d+1)]\), the identity
\(
  1=dp_d+\sum_{i=1}^{d-1}i(p_i-p_{i+1})
  \ge\frac{\kappa d(d+1)}2
\)
forces a unique spectrum, so the class is a singleton and its optimum is
\(\Forb\), not \(\Funiv\).  The compact-set theorem is not asserted uniformly
as the spectrum approaches an eigenvalue collision or the boundary of the
probability simplex.} set
\[
  \Delta_{d,\kappa}
  =\left\{p\in\mathsf A_d:p_d\ge\kappa,\
    p_i-p_{i+1}\ge\kappa\ (i<d)\right\}.
\]
This is a nonempty compact subset of \(\Delta_d^\circ\) and satisfies
\(\Delta_{d,\kappa}
=\overline{\operatorname{int}_{\mathsf A_d}\Delta_{d,\kappa}}
^{\,\mathsf A_d}\),
so \eqref{eq:compact-global-value} applies with
\(K=\Delta_{d,\kappa}\).  Indeed, choose
\(\kappa'\in(\kappa,2/[d(d+1)])\) and set
\(p_i=1/d+((d+1)/2-i)\kappa'\).  Then
\(p_i-p_{i+1}=\kappa'>\kappa\) and
\(p_d>2/[d(d+1)]>\kappa\), so this is a relative-interior point.
The set is closed in the probability simplex, hence compact; being convex
with nonempty relative interior gives the stated closure property.


\subsection{Remarks}
\label{sec:remarks}\label{sec:discussion}

\paragraph{Comparison with purify-and-clone.}
The random-purification channel of Tang--Wright--Zhandry maps \(n\) copies of
any rank-\(r\) state to the uniform mixture of \(n\) identical purifications
on \(\C^d\ot\C^r\) \cite{TangWrightZhandry2025}.  The purify-and-clone
strategy of Li--Theil--Harrow--Chuang then applies Werner's optimal pure-state
cloner on the \(dr\)-dimensional purification space
\cite{Werner1998,LiTheilHarrowChuang2026}.  Their Eq.~(1a) and
Theorem~S3.1 give the all-output root-fidelity guarantee
\[
  F_{\mathrm{RP}}^{\mathrm{root}}
  \ge
  \left(
  \frac{\binom{n+dr-1}{dr-1}}
       {\binom{m+dr-1}{dr-1}}
  \right)^{1/2}
  \longrightarrow\gamma^{-(dr-1)/2}.
\]
The square root of the binomial ratio is the exact root fidelity of the
pure-state cloning step, but after the purifying registers are discarded it
is only the displayed lower bound on the mixed-state fidelity.  For
full-rank inputs, \(r=d\), our exact values satisfy
\[
  \Forb(\gamma,p)>\gamma^{-d(d-1)/4}
  >\gamma^{-(d^2-1)/2},
  \qquad
  \Funiv(\gamma,p)>\gamma^{-(d^2-1)/4}
  >\gamma^{-(d^2-1)/2},
\]
so, at each fixed full-rank nondegenerate spectrum in our regime, both give
strictly stronger proved guarantees.

The comparison can be strengthened from guarantees to the actual
purify-and-clone fidelity. We explicitly show this for qubits in
Appendix~\ref{app:qubit-rp-comparison}:  \Cref{prop:qubit-rp-comparison} proves that the limsup of its root
fidelity is strictly below \(\Forb(\gamma,p)\) for every full-rank
nondegenerate qubit spectrum.  At \(\gamma=2\) and \(p=(2/3,1/3)\), this
limsup is at most \(0.8604\), whereas the root fidelities of our
spectrum-independent and known-spectrum cloners have limiting lower bounds,
respectively,
\[
  \Funiv(2,p)\approx0.8856,
  \qquad
  \Forb(2,p)\approx0.9121.
\]
Thus both of our cloners are strictly superior to the actual
purify-and-clone protocol in this example.

The proved comparison with the actual purify-and-clone protocol is therefore
the qubit result just stated.  Appendix~\ref{app:qubit-rp-comparison} also
gives a modewise sketch suggesting that the same argument extends to
arbitrary dimension.


\paragraph{Dependence on spectral mixedness}
The product formula makes precise one sense in which a more mixed state is
easier to clone.  For fixed \(\gamma>1\),
\[
  f_\gamma'(q)
  =\frac{(\gamma-1)^2}
  {2\sqrt{q(\gamma-1+q)}
   \bigl(\sqrt{\gamma(\gamma-1+q)}+\sqrt q\bigr)^2}>0,
  \qquad 0<q<1.
\]
Consequently, if \(p,p'\in\Delta_d^\circ\) satisfy
\begin{equation}\label{eq:ratio-wise-flattening}
  \frac{p'_j}{p'_i}\ge \frac{p_j}{p_i}
  \qquad\text{for every }i<j,
\end{equation}
then
\[
  \Forb(\gamma,p')\ge \Forb(\gamma,p),
  \qquad
  \Funiv(\gamma,p')\ge \Funiv(\gamma,p),
\]
with strict inequality if any ratio is strictly larger.  Equivalently, every
logarithmic spectral gap has contracted.  This ratio-wise flattening implies
that \(p'\) is majorized by \(p\), but is stronger than majorization.  For
\(d\ge3\), majorization alone need not increase the cloning fidelity; for
\(d=2\), the single eigenvalue ratio does give majorization monotonicity.

\paragraph{Limitations.}
Our results concern the proportional-output asymptotic regime and assume a
full-rank, nondegenerate spectrum.  We do not treat degenerate spectra or
rank-deficient states, for which the LAN model and stabilizer structure
change.  We also do not determine the exact optimal fidelity for finite
\(n\).  It remains unclear whether the cloners constructed here are already
optimal for the corresponding finite-copy problems; we conjecture that they
are.

\paragraph{Conjecture: degenerate and rank-deficient spectra}
We conjecture the following formula for the optimal cloning fidelity when
the spectrum can be degenerate or rank deficient.  Suppose that \(p\) has
distinct eigenvalues
\(\lambda_1>\cdots>\lambda_s\ge0\), with respective multiplicities
\(r_1,\ldots,r_s\), so that
\(\sum_a r_a=d\) and \(\sum_a r_a\lambda_a=1\).  Set
\[
  c_\gamma=\frac{2\sqrt\gamma}{1+\gamma},
  \qquad
  \Forb(\gamma,p)
  =\prod_{a<b}
    f_\gamma\!\left(\frac{\lambda_b}{\lambda_a}\right)^{r_a r_b}.
\]
We conjecture that this \(\Forb(\gamma,p)\) is the optimal
known-spectrum orbital cloning fidelity.  Only pairs of distinct eigenspaces
contribute: rotations within a degenerate eigenspace belong to the stabilizer
of \(\rho_p\) and hence give no oscillator mode.  If \(\lambda_s=0\), the
zero eigenspace is included as the last block, and each positive--zero pair
contributes \(f_\gamma(0)=\gamma^{-1/2}\).

For the unknown-spectrum problem restricted to a fixed-rank stratum, suppose
that \(\lambda_s=0\), and write \(r_0=r_s\) and \(r=d-r_0\).  We conjecture
that the local minimax fidelity within the rank-\(r\) manifold is
\[
  \Funiv^{(r)}(\gamma,p)
  =c_\gamma^{(\sum_{a=1}^{s-1}r_a^2-1)/2}
   \Forb(\gamma,p).
\]
Only the positive-eigenvalue blocks contribute to the classical exponent;
the support--kernel directions are pure Gaussian modes and are already
included in \(\Forb\) through the factors \(f_\gamma(0)\).  This decomposition
is suggested by fixed-rank low-rank LAN \cite{LahiryNussbaum2024}.  In
particular, if the positive eigenvalues are nondegenerate, then
\[
  \Funiv^{(r)}(\gamma,p)
  =c_\gamma^{(r-1)/2}\gamma^{-r(d-r)/2}
   \prod_{1\le i<j\le r}
    f_\gamma\!\left(\frac{\lambda_j}{\lambda_i}\right).
\]

For the full unknown-state local problem, assume instead that
\(\lambda_s>0\).  Allowing arbitrary \(n^{-1/2}\)-perturbations that may split
the degeneracies, we conjecture the value
\[
  \Funiv(\gamma,p)
  =c_\gamma^{(\sum_a r_a^2-1)/2}
   \Forb(\gamma,p).
\]
Here \(\sum_a r_a^2-1\) is the real dimension of the traceless Hermitian
commutant of \(\rho_p\), whose local coordinates should become classical
Gaussian variables, while the cross-eigenspace directions give the quantum
factors above.  In particular, for the maximally mixed
known spectrum the product is empty and the orbital fidelity is one.

As a sanity check, the fixed-rank conjecture has the correct pure-state
specialization.  A rank-one spectrum has one eigenvalue \(1\) and a zero
eigenspace of multiplicity \(d-1\), so
\[
  \Funiv^{(1)}(\gamma,p)=\Forb(\gamma,p)=f_\gamma(0)^{d-1}
  =\gamma^{-(d-1)/2}.
\]
This is exactly the \(m_n/n\to\gamma\) limit of Werner's optimal pure-state
cloning root fidelity \cite{Werner1998}:
\[
  \left(
    \frac{\binom{n+d-1}{d-1}}{\binom{m_n+d-1}{d-1}}
  \right)^{1/2}
  \longrightarrow\gamma^{-(d-1)/2}.
\]

\subsection{Earlier and related work}\label{sec:related-work}

\paragraph{Pure-state cloning and state estimation.}
Universal qubit cloners were constructed in
\cite{BuzekHillery1996,GisinMassar1997,
BrussDiVincenzoEkertFuchsMacchiavelloSmolin1998}; the arbitrary-dimensional
global and single-clone optima were obtained in
\cite{Werner1998,KeylWerner1999}, while links to finite-copy estimation were
developed in \cite{MassarPopescu1995,BrussEkertMacchiavello1998}.
At the single-clone level, cloning with infinitely many outputs reduces to
estimation \cite{BaeAcin2006,ChiribellaDAriano2006}; global all-output
problems have also been studied for broad pure-state families
\cite{CheflesBarnett1999,ChiribellaYang2014Optimal,KimChitambar2022}.
These settings differ from our proportional-output mixed-state problem; see
\cite{ScaraniIblisdirGisinAcin2005,FanWangJingYueShiZhangMu2014} for reviews.

\paragraph{Finite-dimensional mixed-state cloning.}
The earlier finite-dimensional literature summarized above provides bounds
or constructions for restricted state families and criteria, rather than the
exact proportional-output minimax value under joint product fidelity
\cite{CiracEkertMacchiavello1999,Rastegin2003Upper,
Rastegin2002Relative,Rastegin2003StateDependent,Fan2003,FanLiuShi2007,
Kazakov2007,LemmWilde2017}.
More recently, random purification followed by pure-state cloning gave a
universal all-output guarantee for rank-constrained mixed states, but not an
exact optimum \cite{TangWrightZhandry2025,LiTheilHarrowChuang2026}; companion
work instead amplifies a selected eigenstate
\cite{LiTheilHarrowChuangQPA2026}.  See \Cref{sec:discussion} for the
quantitative comparison.

\paragraph{Broadcasting and superbroadcasting.}
Single-input families are broadcastable precisely when they commute
\cite{BarnumCavesFuchsJozsaSchumacher1996}; an input-independent-fidelity
criterion likewise obstructs universal broadcasting for arbitrary mixed
qubits \cite{ChenChen2007MixedBroadcast}.  Multiple-input superbroadcasting
can nevertheless purify individual mixed-qubit marginals; universal,
phase-covariant, and optimal maps are known
\cite{DArianoMacchiavelloPerinotti2005,
BuscemiDArianoMacchiavelloPerinotti2005Optimal,
BuscemiDArianoMacchiavelloPerinotti2006,DangFan2007}, as is a
continuous-variable version for displaced thermal states
\cite{DArianoPerinottiSacchi2006Continuous}.  These marginal criteria allow
correlated outputs, so they do not determine full tensor-power fidelity.
Virtual broadcasting is a related nonphysical extension linking universal
cloning and optimal physical approximation, without changing this distinction
\cite{ParzygnatFullwoodBuscemiChiribella2024}.

\paragraph{Gaussian cloning and LAN.}
Continuous-variable cloning grew from coherent-state amplifier--beam-splitter
schemes \cite{CerfIpeRottenberg2000,
BraunsteinCerfIblisdirVanLoockMassar2001}; under joint rather than single-clone
fidelity, a non-Gaussian cloner can be optimal
\cite{CerfKrugerNavezWernerWolf2005}.  Caves established the
phase-insensitive amplifier noise limit \cite{Caves1982LinearAmplifiers}, and
Gu{\c t}{\u a}--Matsumoto reduced displaced mixed-Gaussian cloning to
least-noise amplification \cite{GutaMatsumoto2006MixedGaussianCloning}.
Quantum LAN was formulated statistically in \cite{GutaJencova2007LAN} and
developed through likelihood-ratio and explicit-channel formulations in
\cite{YamagataFujiwaraGill2013LAN,GutaKahn2006LANQubits,
KahnGuta2009LANFiniteDimensional}, yielding sharp purification, dilution, and
reversal results \cite{BowlesGutaAdesso2011Purification,
BowlesGutaAdesso2012Reversal}; related Gaussian tasks include amplifying and
purifying noisy coherent states \cite{ZhaoChiribella2017}.  This
LAN--amplification lineage underlies our converse.  Low-rank qudit estimation
has different Gaussian asymptotics \cite{LahiryNussbaum2024}, so boundary
models lie outside our full-rank LAN statement.

\paragraph{Other cloning and replication models.}
Further variants include probabilistic exact cloning \cite{DuanGuo1998};
economical, phase-covariant, asymmetric, and telecloning schemes
\cite{NiuGriffiths1998,BrussCinchettiDArianoMacchiavello2000,
Cerf2000Asymmetric,MuraoJonathanPlenioVedral1999}; and cloning of
transformations and superreplication
\cite{ChiribellaDArianoPerinotti2008,ChiribellaYangYao2013,
GendraCalsamigliaMunozTapiaBaganChiribella2014,
DurSekatskiSkotiniotis2015,ChiribellaYangHuang2015,
ChiribellaYang2016}.  A higher-order framework for general channel cloning
includes state cloning \cite{SekatskiGuryanovaKothakondaSkotiniotis2025};
no-signalling yields general cloning bounds
\cite{SekatskiSkotiniotisDur2015NoSignaling,HuTomamichel2024}, while optimal
pure-state cloning and transposition can be complementary channels
\cite{BrzicGrinkoStudzinskiQuintino2026}.  These models differ from ours in
the success rule, family, target object, or replication rate.

\subsection{Organization}

\Cref{sec:label-kernels} develops the common Schur--Cartan lifting, proves a
fidelity bound reducing the cloning problem to a classical label affinity and
a conditional sector fidelity, and then defines the exact target-label kernel
for known spectrum and the randomized dilation kernel for unknown spectrum.
Together with two uniform estimates, these constructions prove the
achievability half of \cref{thm:known-optimum} and part~\textup{(a)} of
\cref{thm:unknown-optimum}.  \Cref{sec:gaussian-LAN} solves the corresponding
Gaussian amplification problems, states the two-way LAN approximation, and
transfers the Gaussian bounds back to finite dimension.  It proves the
converse half of \cref{thm:known-optimum} and the compact-set converse in
part~\textup{(b)} of \cref{thm:unknown-optimum}.  The quantitative comparison
with random purification is summarized in \cref{sec:remarks}; the strict
comparison is developed in Appendix~\ref{app:qubit-rp-comparison}.  The other
appendices supply the fixed-height PBW calculation of the sector fidelity,
the Young local limit and randomized-rounding estimate, and the unified
flat-prior Gaussian converse.

\subsection*{Acknowledgments}

We would like to thank Patrick Hayden and @@@ for discussions and early collaborations. JF acknowledges support from @@@.

\subsection*{Generative AI usage}

The original idea for the cloner construction was due to the authors. We then make use of OpenAI's ChatGPT Sol 5.6 and Anthropic's Fable 5 to assist with the analysis.  The evaluation of the cloning fidelity and the converse argument were benefited from suggestions by AI. The AI tools was conducted to assist with language editing, manuscript
organization, and figure preparation.  All AI-assisted material was reviewed
and revised by the authors, who take full responsibility for the accuracy and
content of the manuscript.

\section*{Glossary of notation}

We collect the notation that recurs across sections.  Symbols local to a
single proof are defined where they first appear.

\begin{description}
\item[Spaces and maps.]
  \(\C\) and \(\R\) denote the complex and real numbers,
  \(\U(d)\) the unitary group, and \(\cT_1(\cH)\) the trace-class
  operators on a Hilbert space \(\cH\).  We write \(A^*\) for the adjoint,
  \(\Tr A\) for the trace, \(\|A\|_1\) for the trace norm, and \(\ot\) for
  tensor products.  A channel is a completely positive trace-preserving map.

\item[Asymptotic scale.]
  The local dimension \(d\) is fixed, \(n\) is the number of inputs,
  \(m_n\) the number of outputs, and
  \(\gamma_n=m_n/n\to\gamma>1\).  The number of pairs \(i<j\) is
  \(r_d=\binom d2\).  A capital \(N\) denotes a generic sample size.

\item[Spectra and states.]
  \(\mathsf A_d=\{x\in\R^d:\sum_i x_i=1\}\), and
  \(\Delta_d^\circ\) is its strictly ordered, positive spectral chamber.
  For \(p\in\Delta_d^\circ\),
  \(\rho_p=\diag(p)\) and \(\rho_{p,g}=g\rho_pg^*\).
  The notation \(K\Subset\Delta_d^\circ\) means that \(K\) is a compact
  subset of this open chamber; relative interiors and closures are taken in
  \(\mathsf A_d\).

\item[Fidelity and optimal values.]
  The unadorned \(F(A,C)\) denotes root (unsquared) fidelity; for classical
  measures it denotes Hellinger affinity.  With \(q_{ij}=p_j/p_i\), the
  one-mode factor is \(f_\gamma(q_{ij})\).  The quantities \(\Forb\),
  \(\Fcl\), and \(\Funiv=\Fcl\Forb\) are respectively the orbital,
  classical, and universal root fidelities in \eqref{eq:main-values}.

\item[Schur--Weyl notation.]
  \(\mathsf Y_N\) is the set of partitions \(\lambda\vdash N\) with at
  most \(d\) rows.  The corresponding unitary and symmetric-group spaces are
  \(\cH_\lambda\) and \(\cS_\lambda\).  The Young-label law under
  \(\rho_p^{\ot N}\) is \(P_{N,p}\), its conditional irreducible state is
  \(\rho_{p,\lambda}\), and \(\mathsf T_N(p)\) is the typical-label set.

\item[Cloners and label kernels.]
  \(\cM_n\) denotes a generic competing channel.  A Markov kernel
  \(K_n(\nu\mid\mu)\) on Young labels is lifted to the covariant channel
  \(\mathcal Q_n[K_n]\).  For compatible labels, \(\cC_{\mu,\nu}\) is the
  Cartan-sector channel; \(\cE_{\mu,\nu}\) extends it to all label pairs by
  a fallback channel.  The constructed known- and unknown-spectrum cloners
  are \(\cC_n^p\) and \(\cC_n^{\mathrm{univ}}\), respectively.  Thus \(K\)
  denotes a spectral set, while \(K_n\) denotes a label kernel.

\item[Local coordinates.]
  \(\mathsf H=\{h\in\R^d:\sum_i h_i=0\}\) is the spectral tangent
  hyperplane.  The variables \(h\in\mathsf H\) and \(z\in\C^{r_d}\)
  describe local spectral and orbital perturbations, respectively;
  \(Y(z)\) is the skew-Hermitian generator of the orbital perturbation,
  \(g_{n,z}=e^{Y(z)/\sqrt n}\), \(\Omega_L\) is a bounded local box, and
  \(v_{n,L}(p)\) is its minimax value.

\item[Gaussian notation.]
  The calligraphic \(\cF\) denotes one-mode Fock space and
  \(\cF_{\mathrm{orb}}=\cF^{\ot r_d}\), distinct from the roman \(F\) for
 fidelity.  The state \(\tau_p\) is the thermal product in the
  input and ideal target, while \(\tau_p^{(\gamma)}\) is the noisy amplifier
  output at zero displacement.  The states \(\Phi_z^{(p)}\) and
  \(\Psi_z^{(p,\gamma)}\) are the displaced input and ideal target built from
  \(\tau_p\).  The symbols \(\Theta^{\mathrm{in}}\) and
  \(\Theta^{\mathrm{tar}}\) include the classical Gaussian factor.  The
  number operator is \(\widehat N\), distinct from the sample size \(N\).

\item[Limits and uniformity.]
  Unless stated otherwise, asymptotics are as \(n\to\infty\) with \(d\) and
  \(\gamma\) fixed.  An \(o(1)\) term is uniform only over the parameters
 specified in the accompanying statement.
\end{description}

\section{Cloner construction and achievability}\label{sec:label-kernels}

Both cloning protocols use the same finite-dimensional
architecture.  We first recall the Schur--Weyl decomposition in
\cref{sec:schur-weyl-decomposition} and lift an arbitrary Markov kernel on
Young labels to a covariant quantum channel in
\cref{sec:cartan-sector-lifting}.  We then choose the exact target-law kernel
for known spectrum and the randomized dilation kernel for unknown spectrum in
\cref{sec:known-label-kernel,sec:unknown-label-kernel}, respectively, and
combine their label estimates with the common sector fidelity in
\cref{sec:uniform-estimates-achievability}.  The two achievability proofs are
completed in \cref{sec:achievability-completion}.  We now develop the
components summarized in
\cref{fig:schur-cartan-cloner}.

\subsection{Schur--Weyl decomposition}
\label{sec:schur-weyl-decomposition}

Let
\[
  \mathsf Y_N=\{\lambda\vdash N:\ell(\lambda)\le d\}.
\]
Schur--Weyl duality gives
\[
  (\C^d)^{\ot N}
  \cong
  \bigoplus_{\lambda\in\mathsf Y_N}\cH_\lambda\ot\cS_\lambda,
\]
where \(\cH_\lambda\) is the irreducible \(\U(d)\)-module of highest
weight \(\lambda\) and \(\cS_\lambda\) is the corresponding Specht module.
For \(\rho_p=\diag(p_1,\ldots,p_d)\),
\begin{equation}\label{eq:schur-state}
  \rho_p^{\ot N}
  =\bigoplus_{\lambda\in\mathsf Y_N}
  P_{N,p}(\lambda)\,
  \rho_{p,\lambda}\ot\frac{I_{\cS_\lambda}}{\dim\cS_\lambda},
  \qquad
  P_{N,p}(\lambda)=\dim\cS_\lambda\,s_\lambda(p),
  \quad
  \rho_{p,\lambda}=\frac{\pi_\lambda(\rho_p)}{s_\lambda(p)}.
\end{equation}
For \(\rho_{p,g}=g\rho_p g^*\), replace \(\rho_{p,\lambda}\) by
\(\rho_{p,g,\lambda}=U_\lambda(g)\rho_{p,\lambda}U_\lambda(g)^*\).  The
Young law \(P_{N,p}\) depends on \(p\), but not on \(g\).

Define
\[
  \mathsf T_N(p)
  =\left\{\lambda\in\mathsf Y_N:
  \left\|\frac\lambda N-p\right\|_\infty\le N^{-1/3}\right\}.
\]

\begin{lemma}[Uniform Young concentration]\label{lem:young-concentration}
There is a constant \(c_d>0\), depending only on \(d\), such that
\[
  \sup_{p\in\Delta_d^\circ}P_{N,p}(\mathsf T_N(p)^c)
  \le e^{-c_d N^{1/3}}
\]
for all sufficiently large \(N\).
\end{lemma}

\begin{proof}
Write \(q=\lambda/N\).  The Weyl dimension formula gives
\(\dim\cH_\lambda\le(N+1)^{d(d-1)/2}\).  Every weight \(\xi\) of
\(\cH_\lambda\) has \(\lambda-\xi\) equal to a nonnegative integral
combination of positive roots.  Since \(p_j/p_i<1\) for \(i<j\), this gives
\(p^\xi\le p^\lambda\).  Hence
\(s_\lambda(p)\le \dim\cH_\lambda\,p^\lambda\).
The standard bounds
\(\dim\cS_\lambda\le\binom{N}{\lambda_1,\ldots,\lambda_d}\) and
\(\binom{N}{\lambda_1,\ldots,\lambda_d}p^\lambda
\le \exp[-ND(q\|p)]\) give
\begin{equation}\label{eq:finite-Young-large-deviation}
  P_{N,p}(\lambda)
  \le (N+1)^{d(d-1)/2}\exp[-ND(q\|p)],
\end{equation}
which is the finite-sample estimate underlying the Keyl--Werner spectrum
large-deviation theorem \cite{KeylWerner2001Spectrum}.

There are at most \((N+1)^{d-1}\) partitions of \(N\) with at most \(d\)
rows.  If \(\lambda\notin\mathsf T_N(p)\), then
\(\|q-p\|_1>N^{-1/3}\), and Pinsker's inequality gives
\(D(q\|p)\ge\tfrac12N^{-2/3}\).  Therefore
\[
  P_{N,p}(\mathsf T_N(p)^c)
  \le (N+1)^{d(d-1)/2+d-1}e^{-N^{1/3}/2}.
\]
For all sufficiently large \(N\), the polynomial factor is at most
\(e^{N^{1/3}/4}\).  Hence the asserted bound holds with \(c_d=1/4\).
\end{proof}

\subsection{Cartan-sector channel and lifting of a label kernel}
\label{sec:cartan-sector-lifting}

Let \(\lambda,\theta\) be dominant weights, with unit highest-weight vectors
\(v_\lambda,v_\theta\).  The Cartan component \(\cH_{\lambda+\theta}\)
occurs with multiplicity one in \(\cH_\lambda\ot\cH_\theta\).  Let
\begin{equation}\label{eq:cartan-inclusion}
  V_{\lambda,\theta}:\cH_{\lambda+\theta}
  \longrightarrow\cH_\lambda\ot\cH_\theta,
  \qquad
  V_{\lambda,\theta}v_{\lambda+\theta}=v_\lambda\ot v_\theta,
\end{equation}
be the normalized intertwining isometry.  For \(\mu\in\mathsf Y_n\) and
\(\nu\in\mathsf Y_m\), write \(\nu-\mu\vdash m-n\) when the coordinatewise
difference \(\nu-\mu\) is a partition of \(m-n\).  If this holds, put
\(\omega=\nu-\mu\) and define
\begin{equation}\label{eq:sector-channel-formula}
  \cC_{\mu,\nu}(X)
  =\frac{\dim\cH_\mu}{\dim\cH_\nu}
  V_{\mu,\omega}^*(X\ot I_{\cH_\omega})V_{\mu,\omega}.
\end{equation}
The partial trace
\(\Tr_{\cH_\omega}(V_{\mu,\omega}V_{\mu,\omega}^*)\) commutes with the
irreducible \(\U(d)\)-action on \(\cH_\mu\); hence Schur's lemma and a trace
comparison give
\(
  \Tr_{\cH_\omega}(V_{\mu,\omega}V_{\mu,\omega}^*)
  =\frac{\dim\cH_\nu}{\dim\cH_\mu}I_{\cH_\mu}
\).
Thus \eqref{eq:sector-channel-formula} is completely positive by
tensoring with \(I_{\cH_\omega}\), applying
\(Y\mapsto V_{\mu,\omega}^*YV_{\mu,\omega}\), and positive rescaling;
the identity above makes it trace
preserving, while equivariance of \(V_{\mu,\omega}\) gives
\(\U(d)\)-covariance.

For every pair \((\mu,\nu)\in\mathsf Y_n\times\mathsf Y_m\), define a
covariant channel \(\cE_{\mu,\nu}:\cT_1(\cH_\mu)\to\cT_1(\cH_\nu)\) by
\[
  \cE_{\mu,\nu}(X)=
  \begin{cases}
    \cC_{\mu,\nu}(X),&\nu-\mu\vdash m-n,\\[1mm]
    \Tr(X)I_{\cH_\nu}/\dim\cH_\nu,&\text{otherwise}.
  \end{cases}
\]
Let \(K_n:\mathsf Y_n\rightsquigarrow\mathsf Y_{m_n}\), with transition
probabilities \(K_n(\nu\mid\mu)\), be any Markov kernel.  Its Schur--Cartan
lifting is the following channel,
denoted by \(\mathcal Q_n[K_n]\): apply the Schur transform, measure
\(\mu\), discard the input Specht factor, sample \(\nu\) from
\(K_n(\cdot\mid\mu)\), apply \(\cE_{\mu,\nu}\), adjoin
\(I_{\cS_\nu}/\dim\cS_\nu\), and apply the inverse Schur transform.  Each
step is completely positive and trace preserving, so their composition is a
channel.  Because the label measurement and sampling are invariant and every
map \(\cE_{\mu,\nu}\) is covariant, \(\mathcal Q_n[K_n]\) is
\(\U(d)\)-covariant.

\begin{proposition}[Fidelity of a lifted label kernel]
\label{prop:lifted-kernel-bound}
Fix \(p\in\Delta_d^\circ\).  Let
\[
  \Pi_{n,p}(\mu,\nu)=P_{n,p}(\mu)K_n(\nu\mid\mu),
  \qquad
  Q_{m_n,p}(\nu)=\sum_\mu\Pi_{n,p}(\mu,\nu).
\]
Suppose \(G\subset\mathsf Y_n\times\mathsf Y_{m_n}\) satisfies
\(\nu-\mu\vdash m_n-n\) for every \((\mu,\nu)\in G\), and
\[
  F\!\left(\cC_{\mu,\nu}(\rho_{p,\mu}),\rho_{p,\nu}\right)\ge b
  \qquad((\mu,\nu)\in G)
\]
for some \(b\in[0,1]\).  Then, for every \(g\in\U(d)\),
\begin{equation}\label{eq:lifted-kernel-bound}
  F\!\left(
    \mathcal Q_n[K_n](\rho_{p,g}^{\ot n}),
    \rho_{p,g}^{\ot m_n}
  \right)
  \ge b\left(F(Q_{m_n,p},P_{m_n,p})-\sqrt{\Pi_{n,p}(G^c)}\right).
\end{equation}
\end{proposition}

\begin{proof}
Covariance makes every sector fidelity independent of \(g\).  In the output
Schur decomposition, the actual state in the \(\nu\)-block is the mixture of
\(\cE_{\mu,\nu}(\rho_{p,g,\mu})\) with weights
\(\Pi_{n,p}(\mu,\nu)\), whereas the ideal block is
\(P_{m_n,p}(\nu)\rho_{p,g,\nu}\).  The Specht factors are identical
and therefore cancel by multiplicativity of the trace norm under tensor
products.

For \(Q_{m_n,p}(\nu)>0\), set
\[
  r(\nu)=\sum_{\mu:(\mu,\nu)\notin G}
  \frac{\Pi_{n,p}(\mu,\nu)}{Q_{m_n,p}(\nu)},
\]
and set \(r(\nu)=0\) otherwise.  The direct-sum identity for fidelity,
its scaling under positive scalar multiples, and joint concavity give
\[
\begin{aligned}
  &F\!\left(
    \mathcal Q_n[K_n](\rho_{p,g}^{\ot n}),
    \rho_{p,g}^{\ot m_n}
  \right)\\
  &\qquad\ge
  b\left(
    F(Q_{m_n,p},P_{m_n,p})
    -\sum_\nu
      \sqrt{Q_{m_n,p}(\nu)P_{m_n,p}(\nu)}\,r(\nu)
  \right).
\end{aligned}
\]
Since \(0\le r\le1\) and
\(\sum_\nu Q_{m_n,p}(\nu)r(\nu)=\Pi_{n,p}(G^c)\),
Cauchy--Schwarz yields
\[
  \sum_\nu\sqrt{Q_{m_n,p}(\nu)P_{m_n,p}(\nu)}\,r(\nu)
  \le
  \sqrt{\Pi_{n,p}(G^c)}
  \left(\sum_\nu P_{m_n,p}(\nu)r(\nu)\right)^{1/2}
  \le\sqrt{\Pi_{n,p}(G^c)}.
\]
This proves \eqref{eq:lifted-kernel-bound}.
\end{proof}

We now specialize the lifted channel by choosing its output-label kernel.
The two protocols differ only in this choice: the known-spectrum construction
samples the exact target law, whereas the unknown-spectrum construction uses
a randomized dilation of the measured input label.

\subsection{Known spectrum: exact target-label sampling}
\label{sec:known-label-kernel}

Fix \(p\in\Delta_d^\circ\).  For every \(n\) such that \(m_n>n\), define
\[
  K_n^p(\nu\mid\mu)=P_{m_n,p}(\nu).
\]
Thus \(\nu\) is independent of \(\mu\) and has the exact target law.  Define
the amplification branch by
\[
  \cC_n^p:=\mathcal Q_n[K_n^p].
\]
Pairs with
\(\nu-\mu\not\vdash m_n-n\) use
\(\cE_{\mu,\nu}(X)=\Tr(X)I_{\cH_\nu}/\dim\cH_\nu\) without changing the
sampled label \(\nu\).  Together with the common convention in
\cref{sec:main-results}, this defines the full \(\U(d)\)-covariant sequence
\((\cC_n^p)_n\), whose output Young law on \(\rho_{p,g}^{\ot n}\) is
\(P_{m_n,p}\) for every \(n\) and \(g\).

\subsection{Unknown spectrum: randomized rounding}
\label{sec:unknown-label-kernel}

For \(N\in\mathbb N\), let
\[
  \mathsf L_N=\left\{x\in\mathbb Z^d:\sum_i x_i=N\right\}.
\]
For every \(n\) such that \(m_n>n\), set
\(\nu_n^{\mathrm{fb}}=(m_n,0,\ldots,0)\).
Given \(\mu\in\mathsf Y_n\), draw independent variables
\(U_1,\ldots,U_{d-1}\), each uniform on \([-1/2,1/2)\), and set
\begin{equation}\label{eq:randomized-rounding}
  \widetilde\nu_i
  =\left\lfloor\gamma_n(\mu_i+U_i)+\frac12\right\rfloor
  \quad(1\le i<d),
  \qquad
  \widetilde\nu_d=m_n-\sum_{i<d}\widetilde\nu_i.
\end{equation}
This is a spectrum-independent lattice version of the dilation
\(\mu\mapsto\gamma_n\mu\): it sends typical labels near \(np\) to labels near
\(m_np\), with \(\widetilde\nu-\mu\) typically a partition.  The uniform
shifts randomize the rounding error, avoiding lattice bias in the Gaussian
limit.
Then \(\widetilde\nu\in\mathsf L_{m_n}\).  Let
\(D_n(\widetilde\nu\mid\mu)\) denote its distribution.  Define
\(K_n^{\mathrm{rnd}}:\mathsf Y_n\rightsquigarrow\mathsf Y_{m_n}\) by
\begin{equation}\label{eq:rounding-kernel}
  K_n^{\mathrm{rnd}}(\nu\mid\mu)=
  \begin{cases}
    D_n(\nu\mid\mu),
      &\nu\ne\nu_n^{\mathrm{fb}}
       \text{ and }\nu-\mu\vdash m_n-n,\\[1mm]
    1-\displaystyle\sum_{\substack{\lambda\in\mathsf Y_{m_n},\\
            \lambda\ne\nu_n^{\mathrm{fb}}\\
            \lambda-\mu\vdash m_n-n}}
       D_n(\lambda\mid\mu),
      &\nu=\nu_n^{\mathrm{fb}},\\[4mm]
    0,&\text{otherwise}.
  \end{cases}
\end{equation}
Under \eqref{eq:rounding-kernel}, every rounded label outside
\(\mathsf Y_{m_n}\), and every
\(\widetilde\nu\in\mathsf Y_{m_n}\) for which
\(\widetilde\nu-\mu\not\vdash m_n-n\), is replaced by
\(\nu_n^{\mathrm{fb}}\), which also retains any mass rounded onto it.
Define the amplification branch by
\[
  \cC_n^{\mathrm{univ}}:=\mathcal Q_n[K_n^{\mathrm{rnd}}].
\]
The rounding variables, \eqref{eq:rounding-kernel}, and the maps
\(\cE_{\mu,\nu}\) are spectrum independent, while
\(\mathcal Q_n[K_n^{\mathrm{rnd}}]\) is \(\U(d)\)-covariant.  Together
with the common convention in \cref{sec:main-results}, this defines the
full spectrum-independent, \(\U(d)\)-covariant sequence
\((\cC_n^{\mathrm{univ}})_n\).

\subsection{Sector and label asymptotics}
\label{sec:uniform-estimates-achievability}

The following proposition is the only representation-theoretic asymptotic
estimate used in the cloning protocols.

\begin{proposition}[Uniform Cartan-sector fidelity]
\label{prop:sector-fidelity}
Let \(K\Subset\Delta_d^\circ\) and let \(\delta_n\ge0\) satisfy
\(\delta_n\to0\).  For all sufficiently large \(n\), every
\(p\in K\), \(\mu\in\mathsf Y_n\), and \(\nu\in\mathsf Y_{m_n}\) satisfying
\[
  \left\|\frac\mu n-p\right\|_\infty\le\delta_n,
  \qquad
  \left\|\frac\nu{m_n}-p\right\|_\infty\le\delta_n
\]
also satisfies \(\nu-\mu\vdash m_n-n\).  Moreover,
\[
  F\!\left(\cC_{\mu,\nu}(\rho_{p,\mu}),\rho_{p,\nu}\right)
  =\Forb(\gamma,p)+o(1)
\]
uniformly over all such \(p,\mu,\nu\).
\end{proposition}

The mechanism is already visible on states obtained by a fixed number of
lowering operations.  Writing \(\omega=\nu-\mu\), the Cartan intertwiner
satisfies
\[
  V_{\mu,\omega}F_\alpha^{(\nu)}
  =\bigl(F_\alpha^{(\mu)}\ot I+I\ot F_\alpha^{(\omega)}\bigr)
    V_{\mu,\omega},
\]
where \(F_\alpha^{(\lambda)}\) is the lowering operator for
\(\alpha=e_i-e_j\) on \(\cH_\lambda\).  After normalization, each excitation
splits binomially between the two factors, with fraction
\[
  s_{\alpha,n}
  =\frac{\mu_i-\mu_j}{\nu_i-\nu_j}
  =\frac1\gamma+o(1)
\]
entering the \(\mu\)-factor.  Thus the conditional thermal parameter
\(q_{ij}=p_j/p_i\) is transformed into
\[
  1-s_{\alpha,n}+s_{\alpha,n}q_{ij}
  \longrightarrow
  1-\frac{1-q_{ij}}{\gamma}
  =:q_{ij}^{(\gamma)}.
\]
The low-weight matrices of the sector output and ideal sector therefore
approach those of \(\tau_p^{(\gamma)}\) and \(\tau_p\), whose fidelity is
\(\Forb(\gamma,p)\).  This is a statement about conditional states for
typical label pairs, not convergence of the varying channels
\(\cC_{\mu,\nu}\) to the Gaussian amplifier.  The exact PBW calculation and
uniform thermal-tail argument that complete the proof are given in
Appendix~\ref{app:sector-fidelity}.

The second estimate concerns only \(K_n^{\mathrm{rnd}}\).

\begin{proposition}[Randomized Young-label dilation]
\label{prop:young-rounding}
There is a constant \(c_{\mathrm{rnd}}<\infty\), depending only on \(d\) and
the fixed sequence \((\gamma_n)\), such that, for every compact
\(K\Subset\Delta_d^\circ\) and all sufficiently large \(n\), every
\(p\in K\) and \(\mu\in\mathsf T_n(p)\) satisfy
\[
  \operatorname{supp}K_n^{\mathrm{rnd}}(\,\cdot\mid\mu)
  \subseteq
  \left\{\nu\in\mathsf Y_{m_n}:
    \nu-\mu\vdash m_n-n,\quad
    \|\nu-\gamma_n\mu\|_\infty\le c_{\mathrm{rnd}}
  \right\}.
\]

For \(m_n>n\), let
\[
  Q_{m_n,p}(\nu)
  :=\sum_{\mu\in\mathsf Y_n}
    P_{n,p}(\mu)K_n^{\mathrm{rnd}}(\nu\mid\mu),
  \qquad \nu\in\mathsf Y_{m_n},
\]
be the induced output-label law.  Then, for every compact
\(K\Subset\Delta_d^\circ\),
\begin{equation}\label{eq:rounding-label-fidelity}
  F(Q_{m_n,p},P_{m_n,p})
  =\Fcl(\gamma,d)+o(1)
  \qquad\text{uniformly for }p\in K.
\end{equation}
\end{proposition}

The uniform dither in \eqref{eq:randomized-rounding} makes the label kernel a
discrete realization of continuous dilation.  After centering at \(Np\) and
scaling by \(\sqrt N\), embed a label law as a piecewise-constant density on
the fundamental cells of the label lattice.  Before fallback replacement,
the intersection-volume identity \eqref{eq:cell-kernel} identifies the
embedded rounded law with the target-cell average of the
\(\sqrt{\gamma_n}\)-dilation of the embedded input law.  Write
\[
  \Sigma_p
  :=\left.(\diag(p)-pp^{\mathsf T})\right|_{
    \{x\in\R^d:\,\sum_i x_i=0\}}.
\]
The Young local limit is \(\mathsf N_{0,\Sigma_p}\), so randomized rounding changes the limit
to \(\mathsf N_{0,\gamma\Sigma_p}\), whereas the ideal \(m_n\)-label law
still converges to \(\mathsf N_{0,\Sigma_p}\).  Their Hellinger affinity is
\[
  F\!\left(\mathsf N_{0,\gamma\Sigma_p},
           \mathsf N_{0,\Sigma_p}\right)
  =\left(\frac{2\sqrt\gamma}{1+\gamma}\right)^{(d-1)/2}
  =\Fcl(\gamma,d).
\]
The support check and fallback rule affect only vanishing probability on
typical labels.

The proofs of \cref{prop:sector-fidelity,prop:young-rounding} are given in
Appendices~\ref{app:sector-fidelity} and~\ref{app:young-rounding},
respectively.

\subsection{Completion of the achievability proofs}
\label{sec:achievability-completion}

\begin{proof}[Known-spectrum achievability in \cref{thm:known-optimum}]
By the covariance established in
\cref{sec:known-label-kernel}, take \(g=I\).  For \(K_n^p\),
\[
  \Pi_{n,p}(\mu,\nu)=P_{n,p}(\mu)P_{m_n,p}(\nu),
  \qquad Q_{m_n,p}=P_{m_n,p}.
\]
Let \(G_{n,p}=\mathsf T_n(p)\times\mathsf T_{m_n}(p)\).  Young
concentration, \cref{lem:young-concentration}, gives
\(\Pi_{n,p}(G_{n,p}^c)=o(1)\).  Apply \cref{prop:sector-fidelity} with
\(K=\{p\}\) and
\[
  \delta_n=n^{-1/3}+m_n^{-1/3}.
\]
It follows that every pair in \(G_{n,p}\) is compatible for all sufficiently
large \(n\), and there is an \(\eta_n\to0\) such that, uniformly on \(G_{n,p}\),
\[
  F\!\left(\cC_{\mu,\nu}(\rho_{p,\mu}),\rho_{p,\nu}\right)
  \ge\Forb(\gamma,p)-\eta_n.
\]
Applying \cref{prop:lifted-kernel-bound} and using
\(F(Q_{m_n,p},P_{m_n,p})=1\) yields
\[
  F\!\left(\cC_n^p(\rho_p^{\ot n}),\rho_p^{\ot m_n}\right)
  \ge
  \bigl(\Forb(\gamma,p)-\eta_n\bigr)
  \bigl(1-\sqrt{\Pi_{n,p}(G_{n,p}^c)}\bigr)
  =\Forb(\gamma,p)-o(1).
\]
\end{proof}

\begin{proof}[Proof of part~\textup{(a)} of \cref{thm:unknown-optimum}]
By the covariance established in
\cref{sec:unknown-label-kernel}, take \(g=I\), and fix
\(K\Subset\Delta_d^\circ\).
Uniformly for \(p\in K\), Young concentration,
\cref{lem:young-concentration}, gives
\(P_{n,p}(\mathsf T_n(p)^c)=o(1)\).  By \cref{prop:young-rounding}, for all
sufficiently large \(n\), every pair \((\mu,\nu)\) with
\(\mu\in\mathsf T_n(p)\) and \(K_n^{\mathrm{rnd}}(\nu\mid\mu)>0\) satisfies
both \(\nu-\mu\vdash m_n-n\) and
\[
  \left\|\frac\nu{m_n}-p\right\|_\infty
  \le
  \frac{\|\nu-\gamma_n\mu\|_\infty}{m_n}
  +\left\|\frac\mu n-p\right\|_\infty
  \le \frac{c_{\mathrm{rnd}}}{m_n}+n^{-1/3}.
\]
Apply \cref{prop:sector-fidelity} with
\[
  \delta_n=n^{-1/3}+\frac{c_{\mathrm{rnd}}}{m_n}.
\]
Then, for some \(\eta_n\to0\), uniformly in
\(p,\mu,\nu\) as above,
\[
  F\!\left(\cC_{\mu,\nu}(\rho_{p,\mu}),\rho_{p,\nu}\right)
  \ge\Forb(\gamma,p)-\eta_n.
\]
The joint label law has bad-pair probability
\(P_{n,p}(\mathsf T_n(p)^c)\), and
\eqref{eq:rounding-label-fidelity} gives
\(F(Q_{m_n,p},P_{m_n,p})=\Fcl(\gamma,d)+o(1)\) uniformly on \(K\).
Applying \cref{prop:lifted-kernel-bound} therefore gives
\[
\begin{aligned}
  F\!\left(\cC_n^{\mathrm{univ}}(\rho_p^{\ot n}),
            \rho_p^{\ot m_n}\right)
  &\ge
  \bigl(\Forb(\gamma,p)-\eta_n\bigr)
  \left(
    F(Q_{m_n,p},P_{m_n,p})
    -\sqrt{P_{n,p}(\mathsf T_n(p)^c)}
  \right)\\
  &=\Funiv(\gamma,p)-o(1),
\end{aligned}
\]
uniformly for \(p\in K\).
\end{proof}

\section{Optimality from LAN and Gaussian converses}\label{sec:gaussian-LAN}

The constructions above establish achievability for the minimax fidelities.
We now prove the matching upper bounds by restricting arbitrary cloners to
local neighborhoods in which cloning becomes Gaussian dilation and
amplification.  We
first introduce the local coordinates used by the limiting models, then solve
the Gaussian problems
in \cref{sec:gaussian-amplification}, state the two-way LAN approximation in
\cref{sec:local-asymptotic-normality}, and finally transfer the resulting
converse inequalities back to the finite-dimensional models in
\cref{sec:minimax-completion}.

The LAN transport common to these converse arguments is summarized in
\cref{fig:LAN-transfer}.  The upper path is the induced Gaussian-model
competitor; the lower path is the ideal target against which its fidelity is
measured.

\begin{figure}[!htbp]
\centering
\begin{tikzpicture}[
  flow/.style={-{Latex[length=2mm]},thick},
  lan/.style={-{Latex[length=2mm]},thick,dashed,teal!70!black},
  compare/.style={{Latex[length=1.8mm]}-{Latex[length=1.8mm]},densely dashed},
  box/.style={draw,rounded corners=2pt,align=center,inner sep=5pt,
    minimum height=10mm,font=\small},
  finite/.style={box,fill=orange!9},
  gaussian/.style={box,fill=blue!6}
]
  \node[gaussian,text width=22mm] (gin) at (0,1.25)
    {Gaussian input\\\(\Phi_\theta\)};
  \node[finite,text width=25mm] (fin) at (3.35,1.25)
    {finite input\\\(\rho_{n,\theta}\)};
  \node[finite,text width=29mm] (factual) at (7.15,1.25)
    {competitor output\\\(\cM_n(\rho_{n,\theta})\)};
  \node[gaussian,text width=31mm] (gactual) at (11.35,1.25)
    {induced output\\\(\Gamma_n(\Phi_\theta)\)};

  \node[finite,text width=29mm] (ftarget) at (7.15,-1.25)
    {finite ideal target\\\(\sigma_{n,\theta}\)};
  \node[gaussian,text width=31mm] (gtarget) at (11.35,-1.25)
    {Gaussian ideal target\\\(\Psi_\theta\)};

  \draw[lan] (gin) -- node[above,font=\scriptsize,text=black] {\(S_n\)} (fin);
  \draw[flow] (fin) -- node[above,font=\scriptsize] {\(\cM_n\)} (factual);
  \draw[flow] (factual) -- node[above,font=\scriptsize] {\(T_n\)} (gactual);
  \draw[lan] (ftarget) -- node[above,font=\scriptsize,text=black] {\(T_n\)}
    (gtarget);
  \draw[compare] (factual) -- node[left,font=\scriptsize] {fidelity}
    (ftarget);
  \draw[compare] (gactual) -- node[right,font=\scriptsize] {fidelity}
    (gtarget);

  \node[draw,rounded corners=2pt,fill=gray!6,align=center,font=\small,
    inner sep=5pt,text width=92mm] at (5.7,-3.0)
    {\(\Gamma_n=T_n\circ\cM_n\circ S_n\).  Hence a Gaussian minimax
     upper bound also bounds the finite-dimensional cloning fidelity.\\[-1pt]
     For cloning, \(\theta=(h,z)\) and the target dilation is
     \((h,z)\mapsto\sqrt\gamma\,(h,z)\); for known spectrum, set \(h=0\).};
\end{tikzpicture}
\caption{Transporting an arbitrary finite-dimensional cloner through
two-way local asymptotic normality.  A dashed arrow identifies its displayed
endpoint up to uniform \(o(1)\) trace-norm error on compact local parameter
sets; the solid \(T_n\) arrow defines the induced output exactly.  The LAN maps
are not exact inverses, and the induced channels \(\Gamma_n\) need not be
Gaussian or converge.  The converse follows from the optimal Gaussian
dilation and amplification bounds.}
\label{fig:LAN-transfer}
\end{figure}

\subsection{Local coordinates}
\label{sec:local-coordinates}

Let
\[
  \mathsf H=\left\{h\in\R^d:\sum_i h_i=0\right\},
  \qquad
  Y(z)=2\sum_{j<k}\left(\bar z_{jk}E_{jk}-z_{jk}E_{kj}\right),
  \quad z\in\C^{r_d},
\]
where \(E_{jk}\) are the standard matrix units, and define
\[
  g_{n,z}=\exp\!\left(\frac{Y(z)}{\sqrt n}\right).
\]
In \(\rho_{p+h/\sqrt n,g_{n,z}}\), the coordinates \(h\) and \(z\)
describe local changes of eigenvalues and eigenbasis, respectively.
Since \(Y(z)^*=-Y(z)\), one has
\(g_{n,z}\in\U(d)\).  For \(L>0\), set
\[
  \Omega_L
  =\left\{(h,z)\in\mathsf H\times\C^{r_d}:
  \|h\|_\infty\le L,\
  |\Re z_{ij}|,|\Im z_{ij}|\le L\ \text{for all }i<j\right\}.
\]
For every fixed \(p\in\Delta_d^\circ\) and \(L>0\), the vector
\(p+h/\sqrt n\) belongs to \(\Delta_d^\circ\) for every
\((h,z)\in\Omega_L\) once \(n\) is sufficiently large.  Under LAN, \(h\)
becomes a \((d-1)\)-dimensional classical Gaussian shift, while \(z\) becomes
a collection of displaced oscillator coordinates.  The Gaussian theorem
below optimizes the two corresponding fidelity losses jointly.

\subsection{Gaussian amplification}
\label{sec:gaussian-amplification}

Let \(\cF\) be the one-mode Fock space with number basis
\(\{\ket{k}:k\ge0\}\), creation and annihilation operators \(a^*,a\), and
number operator \(\widehat N=a^*a\).  We use the Weyl convention
\begin{equation}\label{eq:Weyl-convention}
  D(z)=\exp(z a^*-\bar z a).
\end{equation}
For \(0<q<1\), set
\[
  \tau_q=(1-q)\sum_{k\ge0}q^k\proj{k}.
\]
Two thermal states commute, and hence
\begin{equation}\label{eq:thermal-root-fidelity}
  F(\tau_q,\tau_{q'})
  =\frac{\sqrt{(1-q)(1-q')}}{1-\sqrt{qq'}}.
\end{equation}
For \(G\ge1\), put \(q^{(G)}=1-(1-q)/G\).  The quantum-limited
phase-insensitive amplifier \(\cA_G\) satisfies
\begin{equation}\label{eq:amplifier-action}
  \cA_G\!\left(D(z)\tau_qD(z)^*\right)
  =D(\sqrt G\,z)\tau_{q^{(G)}}D(\sqrt G\,z)^*.
\end{equation}
Equivalently,
\(\cA_G(X)=\Tr_b[S_r(X\ot\proj{0}_b)S_r^*]\), where \(S_r\) is a
two-mode squeezing unitary with \(\cosh^2 r=G\).  Squeezing supplies the
gain \(\sqrt G\), while tracing out the vacuum idler produces the minimum
added noise encoded by \(q^{(G)}\).

For \(p\in\Delta_d^\circ\), define on
\(\cF_{\mathrm{orb}}=\cF^{\ot r_d}\)
\[
  \tau_p=\bigotimes_{i<j}\tau_{q_{ij}},
  \qquad
  \tau_p^{(\gamma)}
  =\bigotimes_{i<j}\tau_{q_{ij}^{(\gamma)}}.
\]
With \(D(z)=\bigotimes_{i<j}D(z_{ij})\), the three relevant
states are
\begin{align*}
  \text{input:}\quad
  &\Phi_z^{(p)}=D(z)\tau_pD(z)^*,\\
  \text{ideal gain-\(\sqrt\gamma\) target:}\quad
  &\Psi_z^{(p,\gamma)}=D(\sqrt\gamma z)\tau_pD(\sqrt\gamma z)^*,\\
  \text{amplifier output:}\quad
  &\cA_\gamma^{\ot r_d}(\Phi_z^{(p)})
    =D(\sqrt\gamma z)\tau_p^{(\gamma)}D(\sqrt\gamma z)^*.
\end{align*}
Thus the ideal target retains the thermal state \(\tau_p\); only the actual
amplifier output contains the additional noise \(\tau_p^{(\gamma)}\).  By
\eqref{eq:thermal-root-fidelity},
\begin{equation}\label{eq:one-mode-f}
  F(\tau_{q^{(\gamma)}},\tau_q)=f_\gamma(q).
\end{equation}
Thus, for every \(z\in\C^{r_d}\), the product-amplifier output and the
corresponding ideal state satisfy
\[
  F\!\left(\cA_\gamma^{\ot r_d}(\Phi_z^{(p)}),
           \Psi_z^{(p,\gamma)}\right)
  =\Forb(\gamma,p).
\]

Equip \(\mathsf H\) with its induced Euclidean measure and recall
\[
  \Sigma_p=\left.(\diag(p)-pp^{\mathsf T})\right|_{\mathsf H}>0.
\]
Write \(\mathsf N_{h,\Sigma_p}\) for the Gaussian law on \(\mathsf H\)
with mean \(h\) and covariance \(\Sigma_p\).  A
classical--quantum state has the form
\[
  R(x)=r(x)\rho_x,\qquad
  r(x)\ge0,\quad \int_{\mathsf H}r(x)\,\dd x=1,
\]
where \(x\) and \(r\) are classical, while \(\rho_x\) is a quantum state on
\(\cF_{\mathrm{orb}}\).  A channel \(\Lambda\) is completely positive and
preserves \(\int\Tr R(x)\,\dd x\).  Here complete positivity
means positivity of \(\operatorname{id}_{M_k}\otimes\Lambda\) for every
\(k\ge1\), with the matrix factor acting on the quantum trace-class fibers;
the classical algebra is abelian.  If
\(S(x)=s(x)\sigma_x\), then
\[
  F(R,S)=\int_{\mathsf H}F(R(x),S(x))\,\dd x
  =\int_{\mathsf H}\sqrt{r(x)s(x)}\,F(\rho_x,\sigma_x)\,\dd x.
\]
With \(\|R\|_1:=\int_{\mathsf H}\|R(x)\|_1\,\dd x\), the usual monotonicity
and joint-concavity properties remain valid, and the first equality above is
the continuous direct-sum identity for fidelity.

\par\indent
Set
\[
  \Theta_{h,z}^{\mathrm{in}}
  =\mathsf N_{h,\Sigma_p}\ot\Phi_z^{(p)},
  \qquad
  \Theta_{h,z}^{\mathrm{tar}}
  =\mathsf N_{\sqrt\gamma h,\Sigma_p}\ot\Psi_z^{(p,\gamma)}.
\]
Let
\[
  \mathsf Z_L
  =\{z\in\C^{r_d}:|\Re z_{ij}|,|\Im z_{ij}|\le L\ \text{for all }i<j\}.
\]

\begin{theorem}[Gaussian amplification optimum]
\label{thm:gaussian-amplification}
For every \(p\in\Delta_d^\circ\),
\begin{align}
  \lim_{L\to\infty}\sup_\Lambda
  \frac1{|\mathsf Z_L|}\int_{\mathsf Z_L}
  F\!\left(\Lambda(\Phi_z^{(p)}),\Psi_z^{(p,\gamma)}\right)\,\dd z
  &=\Forb(\gamma,p),
  \label{eq:gaussian-orbital-optimum}\\
  \lim_{L\to\infty}\sup_\Lambda
  \frac1{|\Omega_L|}\int_{\Omega_L}
  F\!\left(\Lambda(\Theta_{h,z}^{\mathrm{in}}),
            \Theta_{h,z}^{\mathrm{tar}}\right)\,\dd h\,\dd z
  &=\Funiv(\gamma,p).
  \label{eq:gaussian-full-optimum}
\end{align}
The first supremum is over all quantum channels on
\(\cF_{\mathrm{orb}}\); the second is over all classical--quantum channels.
No Gaussianity, covariance, or product-structure assumption is imposed on
the competing channels in either supremum.
\end{theorem}

The achievability mechanism is explicit.  The product amplifier
\(\cA_\gamma^{\ot r_d}\) attains \(\Forb(\gamma,p)\) in the orbital model.
In the full model, tensor it with the classical pushforward
\(x\mapsto\sqrt\gamma x\).  The resulting Gaussian affinity is
\(\Fcl(\gamma,d)\), so multiplicativity gives
\(\Fcl(\gamma,d)\Forb(\gamma,p)=\Funiv(\gamma,p)\).

For the converse, average a near-optimal channel over translates of the
expanding flat-prior boxes.  The vanishing boundary-to-volume error yields a
displacement-covariant limiting map.  A weighted-fidelity witness then turns
the orbital fidelity into a noise moment, which the least-noise inequality
bounds by \(\Forb(\gamma,p)\).  In the full model, covariance forces the
classical trace kernel to have convolution form, and the corresponding
Gaussian witness contributes at most \(\Fcl(\gamma,d)\).  The two bounds
therefore multiply to \(\Funiv(\gamma,p)\).  The compactness, possible
lost-mass, and witness arguments are given in
\cref{prop:flat-prior-converse} and Appendix~\ref{app:gaussian-converse}.

\subsection{Local asymptotic normality}
\label{sec:local-asymptotic-normality}

Here LAN means that, uniformly for \((h,z)\) in compact sets,
the state families \(\rho_{p+h/\sqrt n,g_{n,z}}^{\ot n}\) and
\(\mathsf N_{h,\Sigma_p}\ot\Phi_z^{(p)}\) are asymptotically interconvertible
by channels with vanishing trace-norm error.  The Gaussian and oscillator
factors represent spectral and eigenbasis fluctuations, respectively; cloning
scales \((h,z)\) by \(\sqrt\gamma\).

The rotation coordinate in \cite{KahnGuta2009LANFiniteDimensional} differs
from the oscillator displacement coordinate used here.  For
\(\zeta\in\C^{r_d}\), set
\[
  X_p(\zeta)
  =\sum_{j<k}
    \frac{\zeta_{jk}E_{kj}-\bar\zeta_{jk}E_{jk}}
         {\sqrt{p_j-p_k}},
  \qquad
  (L_p\zeta)_{jk}=-\frac{\zeta_{jk}}{2\sqrt{p_j-p_k}}.
\]
Then \(L_p\) is invertible and
\begin{equation}\label{eq:LAN-coordinate-map}
  X_p(L_p^{-1}z)=Y(z).
\end{equation}
In the convention of
\cite[Eqs.~(4.4) and~(4.19)]{KahnGuta2009LANFiniteDimensional}, the physical
unitary generated by \(X_p(\zeta)\) produces the oscillator state
\(W(w)^*\tau W(w)\), where \(W\) denotes their Weyl operator and
\(w_{jk}=\zeta_{jk}/(2\sqrt{p_j-p_k})\).  Since our convention is
\(D(z)\tau D(z)^*\) with \(D=W\), the corresponding coordinate is
\(z=-w=L_p\zeta\).  Linearity of \(L_p\) implies that multiplication of the
physical coordinate by \(\sqrt\gamma\) multiplies \(z\) by the same factor.

\begin{theorem}[Finite-dimensional LAN and target-sample scaling]
\label{thm:LAN-full}
Fix \(p\in\Delta_d^\circ\).  There are channels
\[
  T_N^{\mathrm{full}}:
  \cT_1((\C^d)^{\ot N})
  \longrightarrow L^1(\mathsf H;\cT_1(\cF_{\mathrm{orb}})),
  \qquad
  S_N^{\mathrm{full}}:
  L^1(\mathsf H;\cT_1(\cF_{\mathrm{orb}}))
  \longrightarrow\cT_1((\C^d)^{\ot N})
\]
such that, for every compact
\(C\subset\mathsf H\times\C^{r_d}\),
\begin{align}
  \sup_{(h,z)\in C}
  \left\|
    T_n^{\mathrm{full}}\!\left(\rho_{p+h/\sqrt n,g_{n,z}}^{\ot n}\right)
    -\mathsf N_{h,\Sigma_p}\ot\Phi_z^{(p)}
  \right\|_1&\longrightarrow0,
  \label{eq:LAN-full-forward}\\
  \sup_{(h,z)\in C}
  \left\|
    S_n^{\mathrm{full}}\!\left(\mathsf N_{h,\Sigma_p}\ot\Phi_z^{(p)}\right)
    -\rho_{p+h/\sqrt n,g_{n,z}}^{\ot n}
  \right\|_1&\longrightarrow0,
  \label{eq:LAN-full-reverse}\\
  \sup_{(h,z)\in C}
  \left\|
    T_{m_n}^{\mathrm{full}}\!\left(
      \rho_{p+h/\sqrt n,g_{n,z}}^{\ot m_n}\right)
    -\mathsf N_{\sqrt\gamma h,\Sigma_p}\ot\Psi_z^{(p,\gamma)}
  \right\|_1&\longrightarrow0.
  \label{eq:LAN-full-target}
\end{align}
\end{theorem}

\begin{proof}
The forward and reverse assertions
\eqref{eq:LAN-full-forward}--\eqref{eq:LAN-full-reverse} are
\cite[Theorem~4.3]{KahnGuta2009LANFiniteDimensional}, after identifying their
\((d-1)\)-dimensional spectral coordinate with \(\mathsf H\) and applying
the coordinate change \eqref{eq:LAN-coordinate-map}.  For the classical
factor, let
\[
  \iota:\R^{d-1}\longrightarrow\mathsf H,
  \qquad
  \iota(u_1,\ldots,u_{d-1})
  =\left(u_1,\ldots,u_{d-1},-\sum_{i<d}u_i\right),
\]
and let \(A\) be the matrix of \(\iota\).  In the coordinates of
\cite[Eq.~(3.3)]{KahnGuta2009LANFiniteDimensional}, the classical covariance
is the \((d-1)\times(d-1)\) matrix
\[
  V(p)_{ij}=\delta_{ij}p_i-p_ip_j,
  \qquad i,j<d.
\]
A direct calculation gives
\[
  A V(p)A^{\mathsf T}=\diag(p)-pp^{\mathsf T}.
\]
Hence, when \(h=\iota(u)\), the pushforward of their classical Gaussian is
\(\mathsf N_{h,\Sigma_p}\).  The coordinate change
\eqref{eq:LAN-coordinate-map} identifies their oscillator state with
\(\Phi_z^{(p)}\).  Their Schr\"odinger-picture maps are completely positive
and trace preserving on the trace-class and classical--quantum \(L^1\)
spaces shown above.

The target-sample assertion \eqref{eq:LAN-full-target} is a consequence
of the first assertion.  Indeed, the identity
\[
  \rho_{p+h/\sqrt n,g_{n,z}}
  =\rho_{p+\sqrt{\gamma_n}h/\sqrt{m_n},
          g_{m_n,\sqrt{\gamma_n}z}}
\]
follows from \(\sqrt{\gamma_n}/\sqrt{m_n}=1/\sqrt n\) and the real linearity
of \(Y\).  Apply \eqref{eq:LAN-full-forward} at sample size \(m_n\) on one
compact set containing all scaled parameters.  Since \(\gamma_n\to\gamma\),
translation of Gaussian densities and Weyl conjugation are uniformly
trace-norm continuous there, which gives \eqref{eq:LAN-full-target}.
\end{proof}

\begin{corollary}[Fixed-spectrum LAN]\label{cor:LAN-orb}
Fix \(p\in\Delta_d^\circ\).  There are channels
\[
  T_N^{\mathrm{orb}}:\cT_1((\C^d)^{\ot N})\to\cT_1(\cF_{\mathrm{orb}}),
  \qquad
  S_N^{\mathrm{orb}}:\cT_1(\cF_{\mathrm{orb}})\to\cT_1((\C^d)^{\ot N})
\]
such that, uniformly for \(z\) in compact sets,
\begin{align}
  \|T_n^{\mathrm{orb}}(\rho_{p,g_{n,z}}^{\ot n})-\Phi_z^{(p)}\|_1
  &\longrightarrow0,
  \label{eq:LAN-known-forward}\\
  \|S_n^{\mathrm{orb}}(\Phi_z^{(p)})-\rho_{p,g_{n,z}}^{\ot n}\|_1
  &\longrightarrow0,
  \label{eq:LAN-known-reverse}\\
  \|T_{m_n}^{\mathrm{orb}}(\rho_{p,g_{n,z}}^{\ot m_n})
       -\Psi_z^{(p,\gamma)}\|_1
  &\longrightarrow0.
  \label{eq:LAN-known-target}
\end{align}
\end{corollary}

\begin{proof}
Restrict \cref{thm:LAN-full} to \(h=0\).  The classical Gaussian factor
\(\mathsf N_{0,\Sigma_p}\) is independent of \(z\).
Set
\(T_N^{\mathrm{orb}}(X):=\int_{\mathsf H}
T_N^{\mathrm{full}}(X)(x)\,\dd x\) and
\(S_N^{\mathrm{orb}}(\sigma):=S_N^{\mathrm{full}}
(\mathsf N_{0,\Sigma_p}\ot\sigma)\).
The bounds
\(\|\int_{\mathsf H}R(x)\,\dd x\|_1
\le\int_{\mathsf H}\|R(x)\|_1\,\dd x\) and
\(\|\mathsf N_{0,\Sigma_p}\ot X\|_1=\|X\|_1\) give the three assertions.
\end{proof}

We will use the following continuity estimate here and in the appendices.  If
\(A,C,A',C'\) are positive trace-class operators with trace at most one, then
\begin{equation}\label{eq:fidelity-continuity}
  |F(A,C)-F(A',C')|
  \le \|A-A'\|_1^{1/2}+\|C-C'\|_1^{1/2}.
\end{equation}
Indeed, the Powers--St\o rmer inequality bounds
\(\|\sqrt A-\sqrt{A'}\|_2\) by \(\|A-A'\|_1^{1/2}\), and similarly for
\(C,C'\).  Adding and subtracting \(\sqrt{A'}\sqrt C\), the reverse triangle
inequality and Schatten H\"older give \eqref{eq:fidelity-continuity}, since
\(\|\sqrt C\|_2,\|\sqrt{A'}\|_2\le1\).

The following lemma formalizes this LAN fidelity transfer.

\begin{lemma}[LAN fidelity transfer]\label{lem:LAN-fidelity-transfer}
Let \(\{\rho_{n,\theta}\}_{\theta\in C}\) and
\(\{\sigma_{n,\theta}\}_{\theta\in C}\) be input and target models on a
compact parameter set \(C\).  Let \(\{\Phi_\theta\}\) and
\(\{\Psi_\theta\}\) be limiting models, and suppose channels \(S_n,T_n\)
satisfy
\[
  \sup_{\theta\in C}\|S_n(\Phi_\theta)-\rho_{n,\theta}\|_1\to0,
  \qquad
  \sup_{\theta\in C}\|T_n(\sigma_{n,\theta})-\Psi_\theta\|_1\to0.
\]
For every channel \(\cM_n\), put \(\Gamma_n=T_n\circ\cM_n\circ S_n\).
Then
\begin{equation}\label{eq:LAN-transfer}
  \inf_{\theta\in C}F(\Gamma_n(\Phi_\theta),\Psi_\theta)
  \ge
  \inf_{\theta\in C}F(\cM_n(\rho_{n,\theta}),\sigma_{n,\theta})-o(1).
\end{equation}
\end{lemma}

\begin{proof}
Let the two uniform approximation errors be \(\delta_n\) and \(\eta_n\).
Since \(T_n\circ\cM_n\) is a channel, uniformly in
\(\theta\),
\(\|\Gamma_n(\Phi_\theta)-T_n\cM_n(\rho_{n,\theta})\|_1
\le \|S_n(\Phi_\theta)-\rho_{n,\theta}\|_1\le\delta_n\).
Also
\(\|T_n(\sigma_{n,\theta})-\Psi_\theta\|_1\le\eta_n\).  Together with
\eqref{eq:fidelity-continuity}, these bounds give
\[
  F(\Gamma_n(\Phi_\theta),\Psi_\theta)
  \ge
  F(T_n\cM_n(\rho_{n,\theta}),T_n(\sigma_{n,\theta}))
  -\sqrt{\delta_n}-\sqrt{\eta_n}.
\]
Finally, monotonicity of fidelity under the channel \(T_n\) gives
\[
  F(T_n\cM_n(\rho_{n,\theta}),T_n(\sigma_{n,\theta}))
  \ge F(\cM_n(\rho_{n,\theta}),\sigma_{n,\theta}).
\]
Taking the infimum over \(\theta\) proves \eqref{eq:LAN-transfer}.
\end{proof}

\paragraph{Conditional sector limits versus LAN.}
Although the same thermal state \(\tau_p\) appears in both arguments, it
represents different limiting operations.  The conditional-sector proof
fixes the input and output Young labels and, on fixed-height cutoffs,
identifies the ideal target with \(\tau_p\) and the Cartan-sector output with
\(\tau_p^{(\gamma)}\); see
\eqref{eq:target-fixed-height-convergence} and
\eqref{eq:output-fixed-height-convergence}.  This comparison directly
explains the achievable orbital factor \(\Forb(\gamma,p)\).  Full LAN instead
retains the Young-label fluctuation as
\(\mathsf N_{0,\Sigma_p}\), while orbital LAN marginalizes this classical
factor.  Thus the common occurrence of \(\tau_p\) reflects conditioning on
Young labels in the sector calculation and marginalizing their fluctuation
in orbital LAN.  We now combine \cref{lem:LAN-fidelity-transfer} with
\cref{thm:gaussian-amplification} to complete the minimax converses.

\subsection{Completion of the minimax proofs}
\label{sec:minimax-completion}

We now apply the same transfer principle in three increasingly broad settings.
The fixed-spectrum Gaussian bound completes \cref{thm:known-optimum}; the full
Gaussian bound yields the pointwise local minimax lemma below; and that lemma,
together with uniform achievability, yields the compact-set statement in
\cref{thm:unknown-optimum}.

\begin{proof}[Known-spectrum converse and completion of \cref{thm:known-optimum}]
For the known-spectrum problem, set
\(u_n(p):=\sup_{\cM_n}\inf_{g\in\U(d)}
F\!\bigl(\cM_n(\rho_{p,g}^{\ot n}),\rho_{p,g}^{\ot m_n}\bigr)\).
Choose \(\cM_n\) with worst-case fidelity at least \(u_n(p)-1/n\).  For
fixed \(L\), set
\[
  \Gamma_n=T_{m_n}^{\mathrm{orb}}\circ\cM_n\circ S_n^{\mathrm{orb}}.
\]
The local orbit \(\{g_{n,z}:z\in\mathsf Z_L\}\) is contained in the full
unitary orbit.  Hence \cref{lem:LAN-fidelity-transfer} gives
\[
  u_n(p)
  \le
  \inf_{z\in\mathsf Z_L}
  F(\Gamma_n(\Phi_z^{(p)}),\Psi_z^{(p,\gamma)})+o(1).
\]
The infimum is no larger than the average.  Taking \(\limsup_n\), then the
supremum over Gaussian-model channels, and finally \(L\to\infty\),
\eqref{eq:gaussian-orbital-optimum} gives
\(\limsup_n u_n(p)\le\Forb(\gamma,p)\).  Together with
\(\liminf_n u_n(p)\ge\Forb(\gamma,p)\), proved in
\cref{sec:label-kernels}, this proves \cref{thm:known-optimum}.
\end{proof}

The universal penalty cannot be expressed by optimizing over a singleton:
if \(p\) is fixed in advance, one is back in the known-spectrum problem.  The
useful pointwise reduction instead considers \(n^{-1/2}\)-neighborhoods of a
fixed interior spectrum and lets the local window grow after the large-\(n\)
limit.  For all sufficiently large \(n\), set
\[
  v_{n,L}(p)
  =\sup_{\cM_n}\inf_{(h,z)\in\Omega_L}
  F\!\left(
    \cM_n\bigl(\rho_{p+h/\sqrt n,g_{n,z}}^{\ot n}\bigr),
    \rho_{p+h/\sqrt n,g_{n,z}}^{\ot m_n}
  \right).
\]

\begin{lemma}[Pointwise local minimax limit]
\label{lem:unknown-local}
For every \(p\in\Delta_d^\circ\),
\[
  \lim_{L\to\infty}\limsup_{n\to\infty}v_{n,L}(p)
  =
  \lim_{L\to\infty}\liminf_{n\to\infty}v_{n,L}(p)
  =\Funiv(\gamma,p).
\]
\end{lemma}

\begin{proof}
Fix \(L\).  For all sufficiently large \(n\), the spectra
\(p+h/\sqrt n\), \((h,z)\in\Omega_L\), lie in one compact subset of
\(\Delta_d^\circ\).  Applying part~\textup{(a)} of
\cref{thm:unknown-optimum} uniformly on this compact set and using continuity
of \(\Funiv(\gamma,\cdot)\) yields
\[
  \liminf_{n\to\infty}v_{n,L}(p)\ge\Funiv(\gamma,p).
\]

For the converse, choose \(\cM_n\) with worst-case fidelity at least
\(v_{n,L}(p)-1/n\), and set
\[
  \Gamma_n=T_{m_n}^{\mathrm{full}}\circ\cM_n\circ S_n^{\mathrm{full}}.
\]
By \cref{lem:LAN-fidelity-transfer},
\[
  v_{n,L}(p)
  \le
  \inf_{(h,z)\in\Omega_L}
  F\!\left(
    \Gamma_n(\Theta_{h,z}^{\mathrm{in}}),
    \Theta_{h,z}^{\mathrm{tar}}
  \right)+o(1).
\]
Average over \(\Omega_L\), take \(\limsup_n\), and then use
\eqref{eq:gaussian-full-optimum} as \(L\to\infty\).  This
gives
\[
  \lim_{L\to\infty}\limsup_{n\to\infty}v_{n,L}(p)
  \le\Funiv(\gamma,p).
\]
Since both \(\limsup_n v_{n,L}(p)\) and
\(\liminf_n v_{n,L}(p)\) are nonincreasing in \(L\), the two outer limits
exist and equal \(\Funiv(\gamma,p)\).
\end{proof}

\begin{proof}[Proof of part~\textup{(b)} of \cref{thm:unknown-optimum}]
For the compact-set problem, set
\(w_n(K):=\sup_{\cM_n}\inf_{\substack{p\in K\\g\in\U(d)}}\)\allowbreak
\(\,F\!\bigl(\cM_n(\rho_{p,g}^{\ot n}),\rho_{p,g}^{\ot m_n}\bigr)\).
Part~\textup{(a)} of
\cref{thm:unknown-optimum} gives
\[
  \liminf_{n\to\infty}w_n(K)
  \ge\inf_{p\in K}\Funiv(\gamma,p).
\]
For the reverse inequality, fix
\(p_0\in\operatorname{int}_{\mathsf A_d}K\).  For every fixed \(L\), all local
spectra \(p_0+h/\sqrt n\), \((h,z)\in\Omega_L\), lie in \(K\) for
sufficiently large \(n\).  Thus \(w_n(K)\le v_{n,L}(p_0)\).  Letting first
\(n\to\infty\) and then \(L\to\infty\), \cref{lem:unknown-local} gives
\[
  \limsup_{n\to\infty}w_n(K)\le\Funiv(\gamma,p_0).
\]
Take the infimum over
\(\operatorname{int}_{\mathsf A_d}K\).  Its closure in \(\mathsf A_d\) is
\(K\), and \(\Funiv(\gamma,\cdot)\) is continuous, so
\[
  \limsup_{n\to\infty}w_n(K)
  \le\inf_{p\in K}\Funiv(\gamma,p).
\]
\end{proof}

\appendix
\clearpage
\section{Strict comparison with purify-and-clone}
\label[appendix]{app:qubit-rp-comparison}

This appendix compares our cloners with the actual mixed-system output of
random purification followed by Werner cloning.  We give the details for
qubits and then sketch the modewise extension to arbitrary dimension.

\begin{lemma}[Partial trace followed by qubit LAN]
\label{lem:qubit-partial-trace-lan}
Let \(p=(a,b)\), with \(a>b>0\), and set
\[
  \ket{\psi_p}=\sqrt a\,\ket{00}+\sqrt b\,\ket{11},
  \qquad
  e_s=\sqrt b\,\ket{00}-\sqrt a\,\ket{11},
  \quad e_+=\ket{01},
  \quad e_-=\ket{10}.
\]
Under the occupation--Fock embedding about \(\psi_p\), partial trace over
the second qubit followed by orbital qubit LAN converges strongly on
trace-class operators to the one-mode Gaussian channel \(\Lambda_p\) that
discards the \(e_s\)-mode and whose retained annihilation operator is
\begin{equation}\label{eq:qubit-partial-trace-mode}
  a_{\mathrm{orb}}
  =\frac{\sqrt a\,b_-+\sqrt b\,b_+^*}{\sqrt{a-b}}.
\end{equation}
Here \(b_+,b_-\) are the annihilation operators of the two off-diagonal
purification-tangent modes.
\end{lemma}

\begin{proof}
For precision, let \(J_M\) map the three-mode number vector
\(\ket{k_s,k_+,k_-}\), \(k_s+k_++k_-\le M\), to the corresponding
normalized symmetric occupation vector with the remaining particles in
\(\psi_p\).  If \(P_M\) projects onto these number vectors, define the
trace-preserving extension
\(
  \mathcal J_M(X)
  =J_MP_MXP_MJ_M^*
   +\Tr[(I-P_M)X]\proj{\psi_p}^{\ot M}
\).
Let \(\Gamma_M\) be the orbital marginal of the standard full qubit LAN
extraction channel
\cite{GutaKahn2006LANQubits,KahnGuta2009LANFiniteDimensional}, with its phase
chosen so that the output annihilator corresponds to the positive-CCR
fluctuation
\[
  A_M=\frac{1}{\sqrt{M(a-b)}}\sum_{\ell=1}^M E_{12}^{(\ell)}.
\]
Then
\(
  \Lambda_M
  :=\Gamma_M\circ\Tr_{E^M}\circ\mathcal J_M
  \to\Lambda_p
\)
strongly in trace norm.

It is enough first to test bounded coherent amplitudes
\(z=(z_s,z_+,z_-)\).  For the coherent projector \(\proj z\), the elementary
multinomial-to-Poisson limit gives
\[
  \sup_{\|z\|\le L}
  \|\mathcal J_M(\proj z)-\proj{\psi_{z,M}}^{\ot M}\|_1
  \longrightarrow0
\]
for every \(L<\infty\), where
\[
  \ket{\psi_{z,M}}
  \propto
  \ket{\psi_p}+M^{-1/2}(z_se_s+z_+e_++z_-e_-).
\]
After partial trace, the product on the right becomes exactly
\((\Tr_E\proj{\psi_{z,M}})^{\ot M}\).  Smooth spectral decomposition and
full qubit LAN, followed by discarding the classical coordinate, therefore
give a displaced copy of \(\tau_{b/a}\).  The first-order off-diagonal entry
of the reduced qubit is
\[
  M^{-1/2}
  (\sqrt b\,z_++\sqrt a\,\overline{z_-}),
\]
so, in the canonical orbital normalization, the limiting displacement is
\((\sqrt a\,z_-+\sqrt b\,\overline{z_+})/\sqrt{a-b}\).  This is precisely
the coherent-state action of \(\Lambda_p\) in
\eqref{eq:qubit-partial-trace-mode}; the \(e_s\)-mode is absent from the
retained orbital coordinate.

The relation is canonical because
\[
  [a_{\mathrm{orb}},a_{\mathrm{orb}}^*]
  =\frac{a-b}{a-b}=1.
\]
Together with the complementary mode
\[
  a_\perp
  =\frac{\sqrt a\,b_++\sqrt b\,b_-^*}{\sqrt{a-b}},
\]
it is a two-mode squeezing transformation; discarding \(a_\perp\) gives the
claimed Gaussian channel.  Finally, the linear span of coherent-state
projectors is trace-norm dense in the trace class: a bounded operator
annihilating this span has vanishing Husimi \(Q\)-symbol, hence is zero by
analyticity, and Hahn--Banach gives the density.  If \(Y\) belongs to this
span, contractivity gives
\[
  \limsup_M\|\Lambda_M(X)-\Lambda_p(X)\|_1
  \le2\|X-Y\|_1.
\]
Approximating \(X\) by such \(Y\) proves strong convergence; for the thermal
input below, the same last step is a number-tail truncation.
\end{proof}

\begin{proposition}[Purify-and-clone for a full-rank qubit]
\label{prop:qubit-rp-comparison}
Let \(p=(a,b)\), where \(a>b>0\), put \(q=b/a\), and suppose that
\(m_n/n\to\gamma>1\).  Denote by
\(\sigma_{n,m_n}^{\mathrm{RP}}(p)\) the mixed-system output obtained by
applying the random-purification channel to \(\rho_p^{\ot n}\), thereby
producing \(n\) identical copies of a random purification of \(\rho_p\),
applying Werner's pure-state \(n\to m_n\) cloner on \(\C^2\ot\C^2\), and
discarding all purifying registers.  Define
\begin{equation}\label{eq:qubit-rp-parameters}
  q_{\mathrm{RP}}
  =\frac{\gamma-1+\gamma q}{\gamma+(\gamma-1)q},
  \qquad
  R_{\mathrm{RP}}(\gamma,q)
  =\frac{\sqrt{(1-q)(1-q_{\mathrm{RP}})}}
         {1-\sqrt{q q_{\mathrm{RP}}}}.
\end{equation}
Then
\begin{equation}\label{eq:qubit-rp-upper-bound}
    \limsup_{n\to\infty}
    F\!\left(
      \sigma_{n,m_n}^{\mathrm{RP}}(p),
      \rho_p^{\ot m_n}
    \right)\le R_{\mathrm{RP}}(\gamma,q)<\Forb(\gamma,p).
\end{equation}
For \(\gamma=2\) and \(p=(2/3,1/3)\), moreover,
\begin{equation}\label{eq:qubit-rp-numerical-comparison}
  \begin{aligned}
    R_{\mathrm{RP}}(2,1/2)
    &=\frac{2+\sqrt{10}}{6}\approx0.8604\\
    &<\Funiv(2,p)\approx0.8856
     <\Forb(2,p)\approx0.9121.
  \end{aligned}
\end{equation}
Thus our known-spectrum cloner is strictly better for every full-rank
nondegenerate qubit spectrum, and our spectrum-independent cloner is also
strictly better in the displayed example.
\end{proposition}

\begin{proof}
Use the canonical purification \(\psi_p\) from
\cref{lem:qubit-partial-trace-lan}.
System-unitary covariance makes the fidelity independent of the chosen
representative of the orbit.
The random-purification output is the Haar mixture of
\(((I\ot U)\proj{\psi_p}(I\ot U)^*)^{\ot n}\), with \(U\in\U(2)\).
Write \(\mathcal W_{n,m}^{(4)}\) for Werner's pure-state cloner on
\(\C^4\).  Its unitary covariance and the unitary invariance of the partial
trace imply, for every \(U\),
\begin{align*}
 &\Tr_{E^m}\mathcal W_{n,m}^{(4)}\!\left(
   ((I\ot U)\proj{\psi_p}(I\ot U)^*)^{\ot n}
 \right)\\
 &\qquad=\Tr_{E^m}\!\left[
   (I\ot U)^{\ot m}
   \mathcal W_{n,m}^{(4)}(\proj{\psi_p}^{\ot n})
   (I\ot U^*)^{\ot m}
 \right]
 =\Tr_{E^m}\mathcal W_{n,m}^{(4)}(\proj{\psi_p}^{\ot n}).
\end{align*}
Consequently the Haar average drops out exactly after the purifying
registers are discarded:
\begin{equation}\label{eq:qubit-rp-haar-drops}
  \sigma_{n,m}^{\mathrm{RP}}(p)
  =\Tr_{E^m}\mathcal W_{n,m}^{(4)}(\proj{\psi_p}^{\ot n}).
\end{equation}

Use the tangent basis \(e_s,e_+,e_-\) from
\cref{lem:qubit-partial-trace-lan}.  In the symmetric occupation basis, let
\(\mathbf k=(k_s,k_+,k_-)\) record the numbers of particles in these
directions.  Werner's symmetric-projection formula shows that its output
before the partial trace in \eqref{eq:qubit-rp-haar-drops} is diagonal, with
eigenvalue
\begin{equation}\label{eq:qubit-werner-occupation-law}
  \pi_{n,m}(\mathbf k)
  =\frac{\binom{m-|\mathbf k|}{n}}
         {\binom{m+3}{n+3}},
  \qquad |\mathbf k|\le m-n,
\end{equation}
on the occupation vector indexed by \(\mathbf k\).  For fixed
\(\mathbf k\),
\[
  \pi_{n,m_n}(\mathbf k)
  \longrightarrow
  \gamma^{-3}
  \left(\frac{\gamma-1}{\gamma}\right)^{|\mathbf k|}.
\]
Both sides are normalized probability masses, so the discrete Scheff\'e
lemma gives trace-norm convergence, under the usual occupation--Fock
embedding, to three independent thermal modes with parameter
\(\vartheta=(\gamma-1)/\gamma\), each of mean photon number
\begin{equation}\label{eq:qubit-pure-cloner-noise}
  N=\frac{\vartheta}{1-\vartheta}=\gamma-1.
\end{equation}

Let \(\omega_n\) denote this output under the occupation--Fock embedding.
Then
\(\|\omega_n-\tau_\vartheta^{\ot3}\|_1\to0\), and contractivity together
with \cref{lem:qubit-partial-trace-lan} gives
\[
  \Gamma_{m_n}(\sigma_{n,m_n}^{\mathrm{RP}}(p))
  =\Lambda_{m_n}(\omega_n)
  \longrightarrow
  \Lambda_p(\tau_\vartheta^{\ot3})
\]
in trace norm.  For the tangent vacuum,
\eqref{eq:qubit-partial-trace-mode} has mean photon number \(b/(a-b)\),
and hence gives the ideal mixed-qubit LAN state \(\tau_q\).  If the two
off-diagonal input modes are instead independent thermal states of mean
\(N=\gamma-1\), their output is thermal with mean
\[
  N_{\mathrm{RP}}
  =\frac{aN+b(N+1)}{a-b}
  =\frac{\gamma-1+b}{a-b}.
\]
Its thermal parameter is
\[
  \frac{N_{\mathrm{RP}}}{N_{\mathrm{RP}}+1}
  =\frac{\gamma-1+b}{\gamma-1+a}
  =q_{\mathrm{RP}}.
\]
Thus the first limit above is \(\tau_{q_{\mathrm{RP}}}\), while ordinary
qubit LAN gives
\(\Gamma_{m_n}(\rho_p^{\ot m_n})\to\tau_q\) in trace norm.  Fidelity is
nondecreasing under channels and continuous under trace-norm convergence.
Therefore, using
\eqref{eq:thermal-root-fidelity},
\[
  \limsup_{n\to\infty}
  F\!\left(\sigma_{n,m_n}^{\mathrm{RP}}(p),\rho_p^{\ot m_n}\right)
  \le F(\tau_{q_{\mathrm{RP}}},\tau_q)
  =R_{\mathrm{RP}}(\gamma,q).
\]

Finally, the optimal mixed-state amplifier has thermal parameter
\[
  q_{\mathrm{opt}}=\frac{\gamma-1+q}{\gamma},
  \qquad
  q_{\mathrm{RP}}-q_{\mathrm{opt}}
  =\frac{q(\gamma-1)(1-q)}
         {\gamma[\gamma+(\gamma-1)q]}>0.
\]
For fixed \(q\), set
\[
  h_q(x)=\frac{\sqrt{(1-q)(1-x)}}{1-\sqrt{qx}}.
\]
For \(x>q\),
\[
  \frac{\dd}{\dd x}\log h_q(x)
  =\frac{\sqrt q-\sqrt x}
         {2\sqrt x(1-x)(1-\sqrt{qx})}<0.
\]
Since \(q_{\mathrm{RP}}>q_{\mathrm{opt}}>q\), it follows that
\[
  R_{\mathrm{RP}}(\gamma,q)
  =h_q(q_{\mathrm{RP}})
  <h_q(q_{\mathrm{opt}})
  =\Forb(\gamma,p).
\]
For \(\gamma=2\) and \(q=1/2\), one has
\(q_{\mathrm{RP}}=4/5\) and
\(R_{\mathrm{RP}}=(2+\sqrt{10})/6\); direct substitution into
\eqref{eq:main-values} gives the remaining comparisons in
\eqref{eq:qubit-rp-numerical-comparison}.
\end{proof}

\paragraph{Higher-dimensional sketch.}
The same occupation--Fock and partial-trace LAN mechanism is expected to
extend to every \(d\); here we only give the resulting modewise calculation.
Take
\(p_1>\cdots>p_d>0\), set \(q_{ij}=p_j/p_i\), and use the canonical
purification \(\sum_i\sqrt{p_i}\ket{ii}\).  After the Haar average drops out,
Werner cloning on \(\C^{d^2}\) has the Fock limit
\(\tau_\vartheta^{\ot(d^2-1)}\), where
\(\vartheta=(\gamma-1)/\gamma\), so every purification-tangent mode has mean
\(N=\gamma-1\).  Identifying a tangent vector with a matrix \(Z\) and writing
\(S=\diag(\sqrt{p_1},\ldots,\sqrt{p_d})\), the differential of partial trace
is \(Z\mapsto ZS+SZ^*\).  Assuming the corresponding general-dimensional
partial-trace LAN statement, the modewise calculation gives the
classical--quantum LAN output
\[
  \mathsf N_{0,(2\gamma-1)\Sigma_p}
  \ot\bigotimes_{i<j}\tau_{q_{ij}^{\mathrm{RP}}},
  \qquad
  q_{ij}^{\mathrm{RP}}
  =\frac{\gamma-1+\gamma q_{ij}}
         {\gamma+(\gamma-1)q_{ij}}.
\]
Indeed, the diagonal tangent block adds covariance
\(2(\gamma-1)\Sigma_p\), while the two ordered off-diagonal modes for each
\(i<j\) give the displayed thermal parameter; the complementary
purification directions are discarded.  The coherent-state and density
argument in the lemma applies independently to these blocks.

Writing \(\sigma_{n,m_n}^{\mathrm{RP}}(p)\) for the resulting mixed-system
output and using \(h_q\) from the proof above, full LAN and fidelity
monotonicity would yield
\begin{equation}\label{eq:general-rp-upper-bound}
  \begin{aligned}
  \limsup_{n\to\infty}
  F\!\left(\sigma_{n,m_n}^{\mathrm{RP}}(p),\rho_p^{\ot m_n}\right)
  &\le B_{\mathrm{RP}}(\gamma,p),\\
  B_{\mathrm{RP}}(\gamma,p)
  &:=
  \left(\frac{\sqrt{2\gamma-1}}{\gamma}\right)^{(d-1)/2}
  \prod_{i<j}h_{q_{ij}}\!\left(q_{ij}^{\mathrm{RP}}\right)\\
  &<\Fcl(\gamma,d)\prod_{i<j}f_\gamma(q_{ij})
   =\Funiv(\gamma,p)<\Forb(\gamma,p).
  \end{aligned}
\end{equation}
The strict middle inequality is modewise: the calculation following
\eqref{eq:qubit-rp-upper-bound} gives
\(q_{ij}^{\mathrm{RP}}>(\gamma-1+q_{ij})/\gamma\), and \(h_{q_{ij}}\)
is decreasing there, while
\(\sqrt{2\gamma-1}/\gamma<2\sqrt\gamma/(1+\gamma)\) for the classical
factor.  Thus the modewise calculation predicts a strict comparison with
both of our cloners in every dimension.  This is not a theorem proved here:
a complete proof would require the general-dimensional
partial-trace LAN statement and its convergence on the relevant thermal
inputs.

\section{Exact-weight PBW calculation of the sector fidelity}
\label[appendix]{app:sector-fidelity}

This appendix proves \cref{prop:sector-fidelity} by comparing the
target sector \(\rho_{p,\nu}\) with the Cartan-sector output
\(\cC_{\mu,\nu}(\rho_{p,\mu})\) on finitely many weight spaces near their
highest weights, then controlling the remaining target mass.
\Cref{lem:PBW-fixed-height} shows that normalized lowering-operator vectors
become asymptotically orthonormal on these spaces, while
\cref{lem:uniform-character-normalization} gives the uniform limit
\(p^\lambda/s_\lambda(p)\to\prod_{i<j}(1-p_j/p_i)\).  The Cartan inclusion
\(V_{\mu,\nu-\mu}\) then splits
each lowering operation between \(\cH_\mu\) and \(\cH_{\nu-\mu}\) with
asymptotically binomial weights, producing the thermal limits \(\tau_p\) and
\(\tau_p^{(\gamma)}\) and hence the fidelity \(\Forb(\gamma,p)\).

Because the calculations involve only weights of height bounded
in terms of \(L\), \(O_L(\cdot)\) may depend on \(L,d,c,C\) in
\(cN\le\Delta_\alpha^\lambda\le CN\) for \(\alpha\in\mathsf R_+\), but is
uniform in \(N\) and all dominant \(\lambda\) satisfying these inequalities.

\subsection{PBW vectors on complete weight spaces}

Identify the positive roots of \(\U(d)\) with
\[
  \mathsf R_+=\{e_i-e_j:1\le i<j\le d\}.
\]
For \(\alpha=e_i-e_j\), set
\[
  F_\alpha=E_{ji},
  \qquad E_\alpha=E_{ij},
  \qquad H_\alpha=[E_\alpha,F_\alpha]=E_{ii}-E_{jj},
  \qquad \Delta_\alpha^\lambda=\lambda_i-\lambda_j.
\]
Fix a total order on \(\mathsf R_+\).  For
\(\mathbf r\in\mathbb Z_{\ge0}^{\mathsf R_+}\), write
\[
  F^{\mathbf r}=\prod_{\alpha\in\mathsf R_+}^{\prec}F_\alpha^{r_\alpha},
  \qquad |\mathbf r|=\sum_\alpha r_\alpha.
\]
For a nonnegative simple-root combination \(\eta\), let
\(\operatorname{ht}(\eta)\) be its simple-root height and put
\[
  \mathsf P(\eta)
  =\left\{\mathbf r:\sum_{\alpha\in\mathsf R_+}r_\alpha\alpha=\eta\right\}.
\]
The set \(\mathsf P(\eta)\) is finite, and the PBW theorem implies that
\(\{F^{\mathbf r}v_\lambda:\mathbf r\in\mathsf P(\eta)\}\) spans
\(\cH_\lambda[\lambda-\eta]\) \cite{FultonHarris1991}.  Define
\begin{equation}\label{eq:PBW-height-subspace}
  \mathcal K_{\lambda,L}
  =\bigoplus_{\operatorname{ht}(\eta)\le L}
    \cH_\lambda[\lambda-\eta].
\end{equation}
Let \(P_{\lambda,L}\) denote the orthogonal projection onto
\(\mathcal K_{\lambda,L}\).

\begin{lemma}[PBW Gram matrices at fixed height]
\label{lem:PBW-fixed-height}
Let \(\lambda=\lambda^{(N)}\) range over dominant weights satisfying
\[
  cN\le\Delta_\alpha^\lambda\le CN
  \qquad(\alpha\in\mathsf R_+)
\]
for fixed \(0<c<C<\infty\).  Define
\begin{equation}\label{eq:normalized-PBW-vectors}
  \widetilde e_{\mathbf r}^\lambda
  =\prod_{\alpha\in\mathsf R_+}^{\prec}
   \frac{F_\alpha^{r_\alpha}}
        {\sqrt{r_\alpha!}(\Delta_\alpha^\lambda)^{r_\alpha/2}}
   v_\lambda.
\end{equation}
For every fixed \(L\), uniformly over
\(\operatorname{ht}(\eta)\le L\) and
\(\mathbf r,\mathbf s\in\mathsf P(\eta)\),
\begin{equation}\label{eq:PBW-Gram-fixed-height}
  \langle\widetilde e_{\mathbf r}^\lambda,
         \widetilde e_{\mathbf s}^\lambda\rangle
  =\delta_{\mathbf r,\mathbf s}+O_L(N^{-1/2}).
\end{equation}
Consequently, for all sufficiently large \(N\), the full Gram matrix
\(G_{\lambda,\eta}\) on every such exact-weight block is invertible.  If
\begin{equation}\label{eq:exact-block-orthonormalization}
  (e_{\mathbf r}^\lambda)_{\mathbf r\in\mathsf P(\eta)}
  =
  (\widetilde e_{\mathbf r}^\lambda)_{\mathbf r\in\mathsf P(\eta)}
  G_{\lambda,\eta}^{-1/2},
\end{equation}
then these vectors form an orthonormal basis of
\(\cH_\lambda[\lambda-\eta]\), and
\begin{equation}\label{eq:PBW-orthonormal-close}
  e_{\mathbf r}^\lambda
  =\widetilde e_{\mathbf r}^\lambda+O_L(N^{-1/2})
\end{equation}
uniformly for \(\operatorname{ht}(\eta)\le L\) and
\(\mathbf r\in\mathsf P(\eta)\).
\end{lemma}

\begin{proof}
For \(\mathbf r\in\mathsf P(\eta)\), every positive root has positive
simple-root height, so
\(|\mathbf r|\le\operatorname{ht}(\eta)\).  Put
\[
  \widehat F_\alpha=\frac{F_\alpha}{\sqrt{\Delta_\alpha^\lambda}},
  \qquad
  \widehat E_\alpha=\frac{E_\alpha}{\sqrt{\Delta_\alpha^\lambda}}.
\]
Each unnormalized Gram entry
\(\langle v_\lambda,(\widehat F^{\mathbf r})^*
\widehat F^{\mathbf s}v_\lambda\rangle\) is therefore a word of length at
most \(2L\).  Commuting root operators cannot increase the word length, and
each root has height at most \(d-1\).  Hence every intermediate vector lies
in
\[
  \mathcal K_{\lambda,R_L},
  \qquad R_L:=2L(d-1).
\]

For \(0\le R\le R_L\), the restrictions
\[
  \widehat E_\alpha:\mathcal K_{\lambda,R}\to\mathcal K_{\lambda,R},
  \qquad
  \widehat F_\alpha:\mathcal K_{\lambda,R}
  \to\mathcal K_{\lambda,R+d-1}
\]
satisfy
\(\|\widehat E_\alpha\|+\|\widehat F_\alpha\|=O_L(1)\).  To see this,
decompose \(\cH_\lambda\) orthogonally into irreducible
  \(\mathfrak{sl}_2(\alpha)\)-summands.  If a summand meets the weight space
  of \(\lambda-\eta\), where \(\operatorname{ht}(\eta)\le R\), let \(t\) be
  its level there and \(M\) its highest \(H_\alpha\)-eigenvalue.
  The top of this string, of weight
  \(\lambda-\eta+t\alpha\), is again a weight of \(\cH_\lambda\); hence
  \(\eta-t\alpha\) is a nonnegative simple-root combination and
  \(t\operatorname{ht}(\alpha)\le\operatorname{ht}(\eta)\le R\).  Moreover,
\[
  M=\Delta_\alpha^\lambda-\eta(H_\alpha)+2t
   =\Delta_\alpha^\lambda+O_L(1).
\]
At level \(t\), the \(F_\alpha\)- and \(E_\alpha\)-coefficients are
\(\sqrt{(t+1)(M-t)}\) and \(\sqrt{t(M-t+1)}\).  After division by
\(\sqrt{\Delta_\alpha^\lambda}\), both are \(O_L(1)\), proving the norm
bound.  All operator norms below are taken after restriction to
\(\mathcal K_{\lambda,R_L}\), with the codomain norm inherited from
\(\cH_\lambda\).

The matrix-unit commutator formula now gives
\begin{equation}\label{eq:normalized-root-commutator}
  [\widehat E_\alpha,\widehat F_\beta]
  =\delta_{\alpha\beta}I+O_L(N^{-1/2}).
\end{equation}
For \(\alpha=\beta\), the left side is
\(H_\alpha/\Delta_\alpha^\lambda=I+O_L(N^{-1})\) on the cutoff.
For \(\alpha\ne\beta\), the unnormalized commutator is zero or,
up to sign, a root operator \(E_\gamma\) or \(F_\gamma\).  In the nonzero
case, the preceding bound gives
\(\|[\widehat E_\alpha,\widehat F_\beta]\|
=O_L\bigl(\sqrt{\Delta_\gamma^\lambda}/
\sqrt{\Delta_\alpha^\lambda\Delta_\beta^\lambda}\bigr)
=O_L(\sqrt N/N)=O_L(N^{-1/2})\).

Commute every \(\widehat E_\alpha\) to the right in
\(\langle v_\lambda,(\widehat F^{\mathbf r})^*
\widehat F^{\mathbf s}v_\lambda\rangle\).  The scalar commutators in
\eqref{eq:normalized-root-commutator} contribute only when
\(\mathbf r=\mathbf s\), in which case they give
\(\prod_\alpha r_\alpha!\).  Every other term vanishes or contains an
\(O_L(N^{-1/2})\) remainder.  Only \(O_L(1)\) terms occur, and all remaining
operator factors have norm \(O_L(1)\); consequently,
\[
  \langle v_\lambda,(\widehat F^{\mathbf r})^*
  \widehat F^{\mathbf s}v_\lambda\rangle
  =\delta_{\mathbf r,\mathbf s}\prod_\alpha r_\alpha!
   +O_L(N^{-1/2}).
\]
Dividing by the factorials in \eqref{eq:normalized-PBW-vectors} proves
\eqref{eq:PBW-Gram-fixed-height}.

The number and sizes of the exact-weight blocks are bounded in terms of
\(L\), so the entrywise estimate gives
\(G_{\lambda,\eta}=I+O_L(N^{-1/2})\) in operator norm.  Thus
\(G_{\lambda,\eta}\) is invertible for large \(N\), and the PBW spanning
vectors form a basis.  Finally, functional calculus gives
\(G_{\lambda,\eta}^{-1/2}=I+O_L(N^{-1/2})\); hence
\eqref{eq:exact-block-orthonormalization} is an orthonormal basis and
satisfies \eqref{eq:PBW-orthonormal-close}.
\end{proof}

\subsection{Character normalization and fixed-height matrices}

For \(a=(a_1,\ldots,a_d)\), write \(p^a=\prod_i p_i^{a_i}\).

\begin{lemma}[Uniform character normalization]
\label{lem:uniform-character-normalization}
Let \(K\Subset\Delta_d^\circ\), let \(\delta_N\to0\), and suppose
\(p\in K\), \(\lambda\vdash N\), and
\(\|\lambda/N-p\|_\infty\le\delta_N\).  Then, uniformly over this family,
\begin{equation}\label{eq:uniform-character-normalization}
  \frac{p^\lambda}{s_\lambda(p)}
  =\prod_{i<j}\left(1-\frac{p_j}{p_i}\right)+o(1).
\end{equation}
\end{lemma}

\begin{proof}
Write the numerator in the Weyl character formula as
\[
  \sum_{\sigma\in\mathfrak S_d}\operatorname{sgn}(\sigma)A_\sigma,
  \qquad
  A_\sigma=\prod_i p_i^{\lambda_{\sigma(i)}+d-\sigma(i)}.
\]
For \(\sigma\ne\mathrm{id}\),
\[
  \frac1N\log\frac{A_{\mathrm{id}}}{A_\sigma}
  =\sum_i\frac{\lambda_i}{N}\log p_i
   -\sum_i\frac{\lambda_{\sigma(i)}}{N}\log p_i+O_K(N^{-1}).
\]
For \(x\in\mathbb R^d\), put
\[
  D_\sigma(x,p)
  =\sum_i x_i\log p_i-\sum_i x_{\sigma(i)}\log p_i.
\]
Strict rearrangement gives \(D_\sigma(p,p)>0\).  Since \(K\) is compact,
the permutations are finite, and the logarithms are uniformly bounded on
\(K\), the assumptions \(\|\lambda/N-p\|_\infty\le\delta_N\) and
\(\delta_N\to0\) give \(c_K>0\) such that, for every
\(\sigma\ne\mathrm{id}\) and all sufficiently large \(N\),
\[
  \frac1N\log\frac{A_{\mathrm{id}}}{A_\sigma}\ge c_K,
  \qquad
  A_\sigma\le e^{-c_KN}A_{\mathrm{id}}.
\]
Consequently,
\[
  \sum_{\sigma\in\mathfrak S_d}\operatorname{sgn}(\sigma)A_\sigma
  =A_{\mathrm{id}}\bigl(1+O_K(e^{-c_KN})\bigr),
  \qquad
  A_{\mathrm{id}}=p^\lambda\prod_i p_i^{d-i}.
\]
Since
\[
  \det(p_i^{d-j})_{i,j=1}^d
  =\prod_i p_i^{d-i}
   \prod_{i<j}\left(1-\frac{p_j}{p_i}\right),
\]
the Weyl character formula gives, uniformly,
\[
  s_\lambda(p)
  =p^\lambda
   \prod_{i<j}\left(1-\frac{p_j}{p_i}\right)^{-1}
   \bigl(1+O_K(e^{-c_KN})\bigr).
\]
This proves \eqref{eq:uniform-character-normalization}.
\end{proof}

For the remainder of the appendix, index the number basis of
\(\cF_{\mathrm{orb}}\) by
\(\mathbf r\in\mathbb Z_{\ge0}^{\mathsf R_+}\), and write it as
\(\{\ket{\mathbf r}\}\).  Define the finite-rank projection
\begin{equation}\label{eq:Fock-height-projection}
  P_L^{\cF}
  =\sum_{\operatorname{ht}(\sum_\alpha r_\alpha\alpha)\le L}
    \proj{\mathbf r}.
\end{equation}
For every fixed \(L\) and all sufficiently large \(N\), each highest
weight \(\lambda\) satisfying the gap bounds in
\cref{lem:PBW-fixed-height} has the orthonormal basis
\(e_{\mathbf r}^\lambda\) of \(\mathcal K_{\lambda,L}\) constructed there.
Define
\begin{equation}\label{eq:finite-coordinate-unitary}
  J_{\lambda,L}:\mathcal K_{\lambda,L}\longrightarrow
  P_L^{\cF}\cF_{\mathrm{orb}},
  \qquad
  J_{\lambda,L}e_{\mathbf r}^\lambda=\ket{\mathbf r}.
\end{equation}
By \cref{lem:PBW-fixed-height}, the two spaces have bases indexed
by the same \(\mathbf r\), so \(J_{\lambda,L}\) is unitary onto
\(P_L^{\cF}\cF_{\mathrm{orb}}\).  It is defined only on
\(\mathcal K_{\lambda,L}\), and it depends on \(\lambda\) and \(L\).

\begin{proof}[Proof of \cref{prop:sector-fidelity}]
Fix \(K\) and \(\delta_n\) as in the proposition, write \(m=m_n\), and put
\(\omega=\nu-\mu\).  Since \(m/n=\gamma_n\to\gamma>1\), we have
\(m/n-1\ge(\gamma-1)/2\) for all sufficiently large \(n\).
Uniformly over the allowed parameters,
\[
\begin{aligned}
  \left\|\frac\omega{m-n}-p\right\|_\infty
  &=\left\|
    \frac{m/n}{m/n-1}\left(\frac\nu m-p\right)
    -\frac1{m/n-1}\left(\frac\mu n-p\right)
  \right\|_\infty\\
  &\le\frac{m/n+1}{m/n-1}\,\delta_n=o(1).
\end{aligned}
\]
Every coordinate and every adjacent gap is uniformly positive on \(K\).
Hence \(\omega\vdash m-n\) for all sufficiently large \(n\).  For
\(\alpha=e_i-e_j\), set
\begin{equation}\label{eq:splitting-fraction}
\begin{aligned}
  s_{\alpha,n}
  &=\frac{\Delta_\alpha^\mu}{\Delta_\alpha^\nu}
    =\frac{n}{m}
      \frac{\mu_i/n-\mu_j/n}{\nu_i/m-\nu_j/m}
    =\frac1\gamma+o(1),\\
  1-s_{\alpha,n}
  &=\frac{\Delta_\alpha^\omega}{\Delta_\alpha^\nu}.
\end{aligned}
\end{equation}
Consequently, there are constants
\(0<c_K<C_K<\infty\), depending only on \(K\) and
\(\gamma\), such that, for all large \(n\),
\begin{equation}\label{eq:uniform-sector-gaps}
  c_Kn
  \le\Delta_\alpha^\mu,\Delta_\alpha^\nu,\Delta_\alpha^\omega
  \le C_Kn
  \qquad(\alpha\in\mathsf R_+).
\end{equation}

Let
\(V=V_{\mu,\omega}:\cH_\nu\to\cH_\mu\ot\cH_\omega\), and write
\(F_\alpha^{(\lambda)}=\dd\pi_\lambda(F_\alpha)\) for the action of the root
operator on \(\cH_\lambda\).  Intertwining and the highest-weight
normalization give
\[
  VF_\alpha^{(\nu)}
  =\bigl(F_\alpha^{(\mu)}\ot I+I\ot F_\alpha^{(\omega)}\bigr)V.
\]
Expanding each
\(\bigl(F_\alpha^{(\mu)}\ot I+I\ot F_\alpha^{(\omega)}\bigr)^{r_\alpha}\)
binomially gives the exact identity
\begin{equation}\label{eq:cartan-PBW-exact}
  V\widetilde e_{\mathbf r}^{\nu}
  =\sum_{\mathbf a\le\mathbf r}
    \beta_{\mathbf r,\mathbf a}^{(n)}
    \widetilde e_{\mathbf a}^{\mu}
    \ot\widetilde e_{\mathbf r-\mathbf a}^{\omega},
\end{equation}
where
\begin{equation}\label{eq:cartan-coefficients}
  \beta_{\mathbf r,\mathbf a}^{(n)}
  =\prod_{\alpha\in\mathsf R_+}
   \sqrt{\binom{r_\alpha}{a_\alpha}
   s_{\alpha,n}^{a_\alpha}
   (1-s_{\alpha,n})^{r_\alpha-a_\alpha}}.
\end{equation}
If \(\operatorname{ht}(\sum_\alpha r_\alpha\alpha)\le L\),
then \eqref{eq:cartan-PBW-exact} has \(O_L(1)\) terms, all of height at
most \(L\), with
\(0\le\beta_{\mathbf r,\mathbf a}^{(n)}\le1\).
Since \(m/n\to\gamma>1\), the sizes \(n,m,m-n\) are uniformly comparable, so
\eqref{eq:uniform-sector-gaps} and \cref{lem:PBW-fixed-height} give
\(e_{\mathbf b}^{\lambda}=\widetilde e_{\mathbf b}^{\lambda}
+O_L(n^{-1/2})\) for \(\lambda\in\{\mu,\nu,\omega\}\) and all relevant
\(\mathbf b\).  Substituting these relations into
\eqref{eq:cartan-PBW-exact} and using \(\|Vx\|=\|x\|\) gives
\begin{equation}\label{eq:cartan-orthonormal-fixed-height}
  Ve_{\mathbf r}^{\nu}
  =\sum_{\mathbf a\le\mathbf r}
    \beta_{\mathbf r,\mathbf a}^{(n)}
    e_{\mathbf a}^{\mu}\ot e_{\mathbf r-\mathbf a}^{\omega}
    +O_L(n^{-1/2})
\end{equation}
uniformly in norm for
\(\operatorname{ht}(\sum_\alpha r_\alpha\alpha)\le L\).

Put
\[
  \sigma_{\mu,\nu}=\cC_{\mu,\nu}(\rho_{p,\mu}),
  \qquad
  c_\lambda(p)=\frac{p^\lambda}{s_\lambda(p)},
  \qquad
  q_\alpha=\frac{p_j}{p_i}\quad(\alpha=e_i-e_j).
\]
For \(\lambda\in\{\mu,\nu\}\), the weight of
\(e_{\mathbf b}^\lambda\) is
\(\lambda-\sum_\alpha b_\alpha\alpha\), and hence
\[
  \rho_{p,\lambda}e_{\mathbf b}^\lambda
  =c_\lambda(p)\prod_\alpha q_\alpha^{b_\alpha}e_{\mathbf b}^\lambda.
\]
The sector channel gives the local formula
\[
  \sigma_{\mu,\nu}
  =\frac{\dim\cH_\mu}{\dim\cH_\nu}
   V^*(\rho_{p,\mu}\ot I_{\cH_\omega})V.
\]
All asymptotic estimates below are uniform over the allowed family.  The
character estimate, \cref{lem:uniform-character-normalization},
gives
\[
  c_\lambda(p)=\prod_\alpha(1-q_\alpha)+o(1)
  \quad(\lambda=\mu,\nu),
\]
while the Weyl dimension formula, \eqref{eq:uniform-sector-gaps}, and
\eqref{eq:splitting-fraction} give
\[
  \frac{\dim\cH_\mu}{\dim\cH_\nu}
  =\prod_{\alpha=e_i-e_j}
    \frac{\Delta_\alpha^\mu+j-i}{\Delta_\alpha^\nu+j-i}
  =\prod_{\alpha\in\mathsf R_+}
    \bigl(s_{\alpha,n}+O(n^{-1})\bigr)
  =\gamma^{-r_d}+o(1).
\]
Thus, substituting
\eqref{eq:cartan-orthonormal-fixed-height} and the weight action above into
this formula, and using
\(\|\rho_{p,\mu}\|\le1\), gives, for fixed \(L\),
\begin{align}
  \langle e_{\mathbf r}^\nu,
    \sigma_{\mu,\nu}e_{\mathbf s}^\nu\rangle
  &={}\frac{\dim\cH_\mu}{\dim\cH_\nu}
  \sum_{\substack{\mathbf a\le\mathbf r\\
                   \mathbf b\le\mathbf s}}
  \beta_{\mathbf r,\mathbf a}^{(n)}
  \beta_{\mathbf s,\mathbf b}^{(n)}
  \langle e_{\mathbf a}^\mu,
    \rho_{p,\mu}e_{\mathbf b}^\mu\rangle
  \langle e_{\mathbf r-\mathbf a}^\omega,
    e_{\mathbf s-\mathbf b}^\omega\rangle
  +O_L(n^{-1/2})\notag\\
  &={}
  \delta_{\mathbf r,\mathbf s}
  \frac{\dim\cH_\mu}{\dim\cH_\nu}c_\mu(p)
  \prod_{\alpha\in\mathsf R_+}
  \bigl(1-s_{\alpha,n}+s_{\alpha,n}q_\alpha\bigr)^{r_\alpha}
  +O_L(n^{-1/2}),
  \label{eq:sector-output-fixed-matrix}\\
  \langle e_{\mathbf r}^\nu,
    \rho_{p,\nu}e_{\mathbf s}^\nu\rangle
  &=\delta_{\mathbf r,\mathbf s}
    c_\nu(p)\prod_{\alpha\in\mathsf R_+}q_\alpha^{r_\alpha}.
  \label{eq:sector-target-fixed-matrix}
\end{align}

Recall the following notation from \cref{sec:gaussian-amplification}:
\[
  q_\alpha^{(\gamma)}=1-\frac{1-q_\alpha}{\gamma},
  \qquad
  \tau_q=(1-q)\sum_{k\ge0}q^k\proj{k},
  \qquad
  \tau_p=\bigotimes_{\alpha\in\mathsf R_+}\tau_{q_\alpha},
  \quad
  \tau_p^{(\gamma)}
  =\bigotimes_{\alpha\in\mathsf R_+}\tau_{q_\alpha^{(\gamma)}}.
\]
Since \(s_{\alpha,n}=\gamma^{-1}+o(1)\),
\eqref{eq:sector-output-fixed-matrix} and
\eqref{eq:sector-target-fixed-matrix} now give
\[
\begin{aligned}
  \langle e_{\mathbf r}^\nu,
    \sigma_{\mu,\nu}e_{\mathbf s}^\nu\rangle
  &=\delta_{\mathbf r,\mathbf s}
    \prod_\alpha(1-q_\alpha^{(\gamma)})
    \bigl(q_\alpha^{(\gamma)}\bigr)^{r_\alpha}+o(1),\\
  \langle e_{\mathbf r}^\nu,
    \rho_{p,\nu}e_{\mathbf s}^\nu\rangle
  &=\delta_{\mathbf r,\mathbf s}
    \prod_\alpha(1-q_\alpha)q_\alpha^{r_\alpha}+o(1).
\end{aligned}
\]
These are precisely the number-basis matrix elements of
\(\tau_p^{(\gamma)}\) and \(\tau_p\).
For fixed \(L\), the cutoff space has fixed finite dimension, so the uniform
entrywise limits imply convergence in trace norm:
\begin{align}
  \left\|
  J_{\nu,L}P_{\nu,L}\sigma_{\mu,\nu}P_{\nu,L}J_{\nu,L}^*
  -P_L^{\cF}\tau_p^{(\gamma)}P_L^{\cF}
  \right\|_1&\longrightarrow0,
  \label{eq:output-fixed-height-convergence}\\
  \left\|
  J_{\nu,L}P_{\nu,L}\rho_{p,\nu}P_{\nu,L}J_{\nu,L}^*
  -P_L^{\cF}\tau_pP_L^{\cF}
  \right\|_1&\longrightarrow0.
  \label{eq:target-fixed-height-convergence}
\end{align}

Both \(\sigma_{\mu,\nu}\) and \(\rho_{p,\nu}\) commute with
the diagonal torus.  The second assertion is immediate; the first follows
from torus invariance of \(\rho_{p,\mu}\) and covariance of
\(\cC_{\mu,\nu}\).
Consequently, \(P_{\nu,L}\) reduces both states, and the direct-sum identity
gives
\[
\begin{aligned}
  F(\sigma_{\mu,\nu},\rho_{p,\nu})
  ={}&F(P_{\nu,L}\sigma_{\mu,\nu}P_{\nu,L},
        P_{\nu,L}\rho_{p,\nu}P_{\nu,L})\\
  &+F((I-P_{\nu,L})\sigma_{\mu,\nu}(I-P_{\nu,L}),
       (I-P_{\nu,L})\rho_{p,\nu}(I-P_{\nu,L})).
\end{aligned}
\]
Schatten H\"older gives \(F(A,C)\le\sqrt{\Tr A\,\Tr C}\), so the second
term is at most
\begin{equation}\label{eq:sector-tail-fidelity-bound}
  \sqrt{\Tr[(I-P_{\nu,L})\rho_{p,\nu}]},
\end{equation}
because \(\Tr[(I-P_{\nu,L})\sigma_{\mu,\nu}]\le1\).

Taking traces in \eqref{eq:target-fixed-height-convergence} gives, for fixed
\(L\),
\[
  \Tr[(I-P_{\nu,L})\rho_{p,\nu}]
  =\Tr[(I-P_L^{\cF})\tau_p]+o(1)
\]
uniformly as \(n\to\infty\).  Since \(K\) is
compact in \(\Delta_d^\circ\), there is \(b_K<1\) such that
\(q_\alpha\le b_K\) for every \(p\in K\) and every positive root
\(\alpha\).  The occupation numbers under \(\tau_p\) are independent
geometric variables with parameters at most \(b_K\).  If their total
occupation is \(s\), then their total root height is at most \((d-1)s\).
Moreover, the probability of any fixed occupation vector of total occupation
\(s\) is at most \(b_K^s\), and there are
\(\binom{s+r_d-1}{r_d-1}\) such vectors.  Therefore
\begin{equation}\label{eq:uniform-thermal-height-tail}
  \sup_{p\in K}\Tr[(I-P_L^{\cF})\tau_p]
  \le
  \sum_{s>L/(d-1)}\binom{s+r_d-1}{r_d-1}b_K^s
  \longrightarrow0\qquad(L\to\infty).
\end{equation}
The same direct-sum decomposition for
\(F(\tau_p^{(\gamma)},\tau_p)\) bounds its complementary part by the square
root of the same target thermal tail.  Consequently,
\[
\begin{aligned}
 &\left|F(\sigma_{\mu,\nu},\rho_{p,\nu})
       -F(\tau_p^{(\gamma)},\tau_p)\right|\\
 &\quad\le
 \left|F(P_{\nu,L}\sigma_{\mu,\nu}P_{\nu,L},
         P_{\nu,L}\rho_{p,\nu}P_{\nu,L})
       -F(P_L^{\cF}\tau_p^{(\gamma)}P_L^{\cF},
          P_L^{\cF}\tau_pP_L^{\cF})\right|\\
 &\qquad+
 \sqrt{\Tr[(I-P_{\nu,L})\rho_{p,\nu}]}
 +\sqrt{\Tr[(I-P_L^{\cF})\tau_p]}.
\end{aligned}
\]

By \eqref{eq:fidelity-continuity} and
\eqref{eq:output-fixed-height-convergence}--\eqref{eq:target-fixed-height-convergence},
for every fixed \(L\),
\[
\begin{aligned}
  &F(P_{\nu,L}\sigma_{\mu,\nu}P_{\nu,L},
     P_{\nu,L}\rho_{p,\nu}P_{\nu,L})\\
  &\qquad\longrightarrow
  F(P_L^{\cF}\tau_p^{(\gamma)}P_L^{\cF},
    P_L^{\cF}\tau_pP_L^{\cF})
\end{aligned}
\]
uniformly.  First let \(n\to\infty\) at fixed \(L\), and then let
\(L\to\infty\) using
\eqref{eq:sector-tail-fidelity-bound}--\eqref{eq:uniform-thermal-height-tail}.
This proves
\[
  \left|F(\sigma_{\mu,\nu},\rho_{p,\nu})
  -F(\tau_p^{(\gamma)},\tau_p)\right|\longrightarrow0
\]
uniformly.  Finally,
\[
  F(\tau_p^{(\gamma)},\tau_p)
  =\prod_{i<j}F(\tau_{q_{ij}^{(\gamma)}},\tau_{q_{ij}})
  =\Forb(\gamma,p).
\]
Here multiplicativity of fidelity follows directly from multiplicativity of
the trace norm under tensor products, and the final equality uses
\eqref{eq:thermal-root-fidelity} and \eqref{eq:one-mode-f}.
\end{proof}

\section{Young local limit and randomized rounding}
\label[appendix]{app:young-rounding}

This appendix proves \cref{prop:young-rounding} by expressing the
randomized rounding law \(D_n(\widetilde\nu\mid\mu)\) as
the normalized intersection volume of the dilated source cell
around \(\gamma_n\mu\) with the target cell labeled by \(\widetilde\nu\).
\Cref{lem:young-local-limit} shows that the centered and rescaled target label
law \(P_{N,p}\) converges uniformly to \(\mathsf N_{0,\Sigma_p}\), while
\cref{lem:randomized-dilation} shows that sending \(P_{n,p}\) through \(D_n\)
changes this limit to \(\mathsf N_{0,\gamma\Sigma_p}\) and yields the label
affinity \(\Fcl(\gamma,d)\).  The final proof checks that typical rounded
labels satisfy \(\nu-\mu\vdash m_n-n\) and that replacement by
\(\nu_n^{\mathrm{fb}}\) moves only vanishing mass.

Define the half-open
fundamental cell
\[
  \mathsf C
  =\left\{\sum_{i=1}^{d-1}t_i(e_i-e_d):-\frac12\le t_i<\frac12\right\}
  \subset\mathsf H.
\]
Volumes are taken with the induced \((d-1)\)-dimensional Lebesgue measure.
The rounding rule \eqref{eq:randomized-rounding} is equivalent, outside a
null set of cell boundaries, to choosing the unique
\(\widetilde\nu\in\mathsf L_{m_n}\) for which
\[
  \gamma_n(\mu+U)\in\widetilde\nu+\mathsf C,
  \qquad
  U=\sum_{i=1}^{d-1}U_i(e_i-e_d).
\]
For \(i<d\), this is precisely the nearest-integer rule because
\(\lfloor x+\tfrac12\rfloor\) is the unique integer in
\((x-\tfrac12,x+\tfrac12]\); the \(d\)-th coordinate is then automatic
because both sides sum to \(m_n\).
Since \(U\) is uniform on \(\mathsf C\), the change of variables
\(x=\gamma_n(\mu+u)\) gives \(\dd x=\gamma_n^{d-1}\dd u\) and sends the
event domain to
\(\gamma_n(\mu+\mathsf C)\cap(\widetilde\nu+\mathsf C)\).
Consequently,
\begin{equation}\label{eq:cell-kernel}
  D_n(\widetilde\nu\mid\mu)
  =\frac{|\gamma_n(\mu+\mathsf C)\cap(\widetilde\nu+\mathsf C)|}
         {\gamma_n^{d-1}|\mathsf C|}.
\end{equation}
All label laws below are regarded as functions on
\(\mathsf L_N\): if a law \(R\) is originally defined on \(\mathsf Y_N\),
set \(R(\lambda)=0\) for every
\(\lambda\in\mathsf L_N\setminus\mathsf Y_N\).

For \(\lambda\in\mathsf L_N\), define the recentered and
rescaled lattice cell
\[
  \mathsf C_{N,p}(\lambda)=\frac{\lambda-Np+\mathsf C}{\sqrt N}.
\]
For any \(R:\mathsf L_N\to\R\) satisfying
\(\sum_{\lambda\in\mathsf L_N}|R(\lambda)|<\infty\), define
\[
  (\mathcal E_{N,p}R)(x)
  =\frac{N^{(d-1)/2}}{|\mathsf C|}R(\lambda)
  \quad(x\in\mathsf C_{N,p}(\lambda)).
\]
The cells \(\mathsf C_{N,p}(\lambda)\), \(\lambda\in\mathsf L_N\),
tile \(\mathsf H\), and their boundaries have zero induced Lebesgue measure.
The map \(\mathcal E_{N,p}\) preserves \(\ell^1\)-distance,
\(\|\mathcal E_{N,p}R-\mathcal E_{N,p}S\|_1=\|R-S\|_{\ell^1}\), and, for
nonnegative \(R,S\), Hellinger affinity,
\(\int_{\mathsf H}\sqrt{(\mathcal E_{N,p}R)(\mathcal E_{N,p}S)}\,\dd x
=\sum_{\lambda\in\mathsf L_N}\sqrt{R(\lambda)S(\lambda)}\).
Indeed, each \(\mathsf C_{N,p}(\lambda)\) has volume
\(|\mathsf C|N^{-(d-1)/2}\), which cancels the prefactor in both cellwise
integrals.

\begin{lemma}[Uniform Young local limit]\label{lem:young-local-limit}
For every compact \(K\Subset\Delta_d^\circ\),
\begin{equation}\label{eq:young-local-limit}
  \sup_{p\in K}
  \left\|\mathcal E_{N,p}P_{N,p}-\varphi_{0,\Sigma_p}\right\|_1
  \longrightarrow0.
\end{equation}
\end{lemma}

\begin{proof}
Let \(M_{N,p}\) be the multinomial law on nonnegative integer
compositions of \(N\), extended by zero to \(\mathsf L_N\).  We first show
\begin{equation}\label{eq:young-multinomial-l1}
  \sup_{p\in K}\|P_{N,p}-M_{N,p}\|_{\ell^1}\longrightarrow0.
\end{equation}
Set \(\mathsf W_{N,p}=\{\lambda\in\mathsf L_N:
\|\lambda-Np\|_\infty\le N^{2/3}\}\).  For all large \(N\),
\(\mathsf W_{N,p}\cap\operatorname{supp}M_{N,p}\subset\mathsf Y_N\), since
\(\lambda_i-\lambda_{i+1}\ge N(p_i-p_{i+1})-2N^{2/3}>0\) for \(i<d\),
uniformly over \(p\in K\).

For \(\lambda\in\mathsf Y_N\), the Weyl dimension and character formulas
give
\[
\begin{gathered}
  P_{N,p}(\lambda)=\dim\cS_\lambda\,s_\lambda(p),\qquad
  \dim\cS_\lambda
  =N!\frac{\prod_{i<j}(\lambda_i-\lambda_j+j-i)}
  {\prod_i(\lambda_i+d-i)!},\\
  s_\lambda(p)
  =\frac{\det[p_i^{\lambda_j+d-j}]_{i,j}}
  {\prod_{i<j}(p_i-p_j)}.
\end{gathered}
\]
Uniformly on \(\mathsf W_{N,p}\cap\operatorname{supp}M_{N,p}\),
\[
  s_\lambda(p)
  =\frac{\prod_i p_i^{\lambda_i+d-i}}
  {\prod_{i<j}(p_i-p_j)}\bigl(1+O_K(e^{-c_KN})\bigr)
\]
for some \(c_K>0\).  Indeed, both \((\lambda_i+d-i)/N\) and
\(\log p_i\) are uniformly strictly decreasing there, so strict
rearrangement makes every nonidentity determinant term exponentially
smaller than the identity term.

Comparison with
\(M_{N,p}(\lambda)=N!\prod_i p_i^{\lambda_i}/\prod_i\lambda_i!\) gives
\[
  \frac{P_{N,p}(\lambda)}{M_{N,p}(\lambda)}
  =\bigl(1+O_K(e^{-c_KN})\bigr)
  \prod_{i<j}\frac{\lambda_i-\lambda_j+j-i}{N(p_i-p_j)}
  \prod_i\frac{\lambda_i!(Np_i)^{d-i}}{(\lambda_i+d-i)!}.
\]
Writing the last product as
\(\prod_i\prod_{a=1}^{d-i}\frac{Np_i}{\lambda_i+a}\), each factor
\(\frac{\lambda_i-\lambda_j+j-i}{N(p_i-p_j)}\), \(i<j\), and
\(\frac{Np_i}{\lambda_i+a}\), \(1\le a\le d-i\), equals
\(1+O_K(N^{-1/3})\).  Therefore
\[
\begin{aligned}
  \|P_{N,p}-M_{N,p}\|_{\ell^1}
  \le{}&\sup_{\lambda\in\mathsf W_{N,p}\cap\operatorname{supp}M_{N,p}}
       \left|\frac{P_{N,p}(\lambda)}{M_{N,p}(\lambda)}-1\right|\\
  &+P_{N,p}(\mathsf W_{N,p}^c)+M_{N,p}(\mathsf W_{N,p}^c)
  \longrightarrow0
\end{aligned}
\]
uniformly for \(p\in K\), since the last two terms vanish by
\cref{lem:young-concentration} and Hoeffding's inequality.

By \eqref{eq:young-multinomial-l1} and the \(\ell^1\)-isometry
of \(\mathcal E_{N,p}\), it remains to show
\begin{equation}\label{eq:embedded-multinomial-limit}
  \sup_{p\in K}
  \|\mathcal E_{N,p}M_{N,p}-\varphi_{0,\Sigma_p}\|_1\longrightarrow0.
\end{equation}
Fix
\(\beta_0\in(1/2,2/3)\) and set
\[
  V_{N,p}=\{\lambda\in\mathsf L_N:
  \|\lambda-Np\|_\infty\le N^{\beta_0}\},
  \qquad x_\lambda=\frac{\lambda-Np}{\sqrt N}.
\]
Since \(\inf_{p\in K}p_d>0\),
\(V_{N,p}\subset\operatorname{supp}M_{N,p}\) for all sufficiently large
\(N\), uniformly in \(p\in K\).
Uniformly for \(p\in K\) and \(\lambda\in V_{N,p}\),
\begin{equation}\label{eq:multinomial-local-estimate}
  \frac{N^{(d-1)/2}}{|\mathsf C|}M_{N,p}(\lambda)
  =\varphi_{0,\Sigma_p}(x_\lambda)(1+o(1)).
\end{equation}
To see this, uniform Stirling and third-order Taylor expansions give
\[
  \log M_{N,p}(\lambda)
  =-\frac{d-1}{2}\log(2\pi N)-\frac12\sum_i\log p_i
   -\frac12\sum_i\frac{x_{\lambda,i}^2}{p_i}
   +O_K\!\left(N^{3\beta_0-2}+N^{\beta_0-1}+N^{-1}\right).
\]
This is the Gaussian exponent since
\(\langle x,\Sigma_p^{-1}x\rangle=\sum_i x_i^2/p_i\).
The normalizing constants agree since
\(|\mathsf C|^2=\det[\langle e_i-e_d,e_j-e_d\rangle]_{i,j<d}=d\), while
\(\ker\Sigma_p=\operatorname{span}\{(1,\ldots,1)\}\) and
\(\operatorname{cof}_{ii}(\Sigma_p)=\prod_jp_j\) give
\(\det_{\mathsf H}\Sigma_p=d\prod_jp_j\).
The remainder is
\(o(1)\) since \(1/2<\beta_0<2/3\).

Set
\[
  U_{N,p}=\bigcup_{\lambda\in V_{N,p}}\mathsf C_{N,p}(\lambda).
\]
The multinomial mass outside \(V_{N,p}\) and the Gaussian mass
outside \(U_{N,p}\) are exponentially small: for some
\(c_K,C_K>0\),
\begin{equation}\label{eq:multinomial-gaussian-tail}
  \sup_{p\in K}\left[
  M_{N,p}(V_{N,p}^c)
  +\int_{U_{N,p}^c}\varphi_{0,\Sigma_p}(x)\,\dd x
  \right]
  \le C_K e^{-c_K N^{2\beta_0-1}}.
\end{equation}
Indeed, if \(\mathsf C_{N,p}(\lambda)\) contains \(x\), then
\(\|x_\lambda\|_\infty\le\|x\|_\infty+(d-1)/(2\sqrt N)\).  Thus
\[
  \{x:\|x\|_\infty\le
  N^{\beta_0-1/2}-(d-1)/(2\sqrt N)\}\subset U_{N,p}.
\]
On the set at left \(\|x_\lambda\|_\infty\le N^{\beta_0-1/2}\),
so \(\lambda\in V_{N,p}\).  Hoeffding controls \(M_{N,p}(V_{N,p}^c)\);
the complementary inclusion and the uniform eigenvalue bounds for
\(\Sigma_p\) control the Gaussian tail on \(U_{N,p}^c\).

The Gaussian density has vanishing relative oscillation on
\(\mathsf C_{N,p}(\lambda)\), \(\lambda\in V_{N,p}\).
Indeed, uniformly for \(p\in K\), \(\lambda\in V_{N,p}\), and
\(x,y,z\in\mathsf C_{N,p}(\lambda)\),
\(\|x-y\|=O(N^{-1/2})\) and
\(\|z\|=O(N^{\beta_0-1/2})\).
Since \(\nabla\log\varphi_{0,\Sigma_p}(z)=-\Sigma_p^{-1}z\),
\begin{equation}\label{eq:gaussian-cell-oscillation}
\begin{aligned}
  \left|\log\frac{\varphi_{0,\Sigma_p}(x)}
                       {\varphi_{0,\Sigma_p}(y)}\right|
  &\le \sup_{z\in\mathsf C_{N,p}(\lambda)}
      \|\nabla\log\varphi_{0,\Sigma_p}(z)\|\,\|x-y\|\\
  &\le O_K(N^{\beta_0-1/2})\,O(N^{-1/2})
   =O_K(N^{\beta_0-1}),\\
  \left|\frac{\varphi_{0,\Sigma_p}(x)}
                   {\varphi_{0,\Sigma_p}(y)}-1\right|
  &=O_K(N^{\beta_0-1})=o(1).
\end{aligned}
\end{equation}
The second estimate follows by exponentiation.

The embedded multinomial density is consequently uniformly accurate on
\(U_{N,p}\).  Indeed, \(0\in\mathsf C\), so the point
\(x_\lambda=(\lambda-Np)/\sqrt N\) lies in
\(\mathsf C_{N,p}(\lambda)\).  Uniformly for \(p\in K\),
\(\lambda\in V_{N,p}\), and \(x\) in the same cell,
\begin{equation}\label{eq:embedded-multinomial-cellwise}
\begin{aligned}
  \frac{N^{(d-1)/2}M_{N,p}(\lambda)}
       {|\mathsf C|\varphi_{0,\Sigma_p}(x)}
  &=\frac{N^{(d-1)/2}M_{N,p}(\lambda)}
          {|\mathsf C|\varphi_{0,\Sigma_p}(x_\lambda)}
    \frac{\varphi_{0,\Sigma_p}(x_\lambda)}
         {\varphi_{0,\Sigma_p}(x)}\\
  &=(1+o(1))(1+o(1))=1+o(1).
\end{aligned}
\end{equation}
The two factors are \(1+o(1)\) by the multinomial local
estimate \eqref{eq:multinomial-local-estimate} and the Gaussian
cell-oscillation bound \eqref{eq:gaussian-cell-oscillation}, respectively.

For each \(p\in K\),
\[
\begin{aligned}
  \|\mathcal E_{N,p}M_{N,p}-\varphi_{0,\Sigma_p}\|_1
  ={}&\int_{U_{N,p}}
      |(\mathcal E_{N,p}M_{N,p})(x)-\varphi_{0,\Sigma_p}(x)|\,\dd x\\
     &+\int_{U_{N,p}^c}
      |(\mathcal E_{N,p}M_{N,p})(x)-\varphi_{0,\Sigma_p}(x)|\,\dd x
  \\[-2pt]
  \le{}&M_{N,p}(V_{N,p}^c)
  +\int_{U_{N,p}^c}\varphi_{0,\Sigma_p}(x)\,\dd x\\
  &+\sup_{\substack{\lambda\in V_{N,p}\\
                     x\in\mathsf C_{N,p}(\lambda)}}
  \left|
  \frac{N^{(d-1)/2}M_{N,p}(\lambda)}
       {|\mathsf C|\varphi_{0,\Sigma_p}(x)}-1
  \right|,
\end{aligned}
\]
Here the inequality uses
\(\int_{U_{N,p}}\varphi_{0,\Sigma_p}(x)\,\dd x\le1\) and the cell
normalization
\(\int_{U_{N,p}^c}(\mathcal E_{N,p}M_{N,p})(x)\,\dd x
=M_{N,p}(V_{N,p}^c)\).
Estimates \eqref{eq:multinomial-gaussian-tail} and
\eqref{eq:embedded-multinomial-cellwise} make all three terms on the right
vanish uniformly, proving
\eqref{eq:embedded-multinomial-limit}.

Finally, the \(\ell^1\)-isometry of \(\mathcal E_{N,p}\) gives
\[
\begin{aligned}
  \sup_{p\in K}
  \|\mathcal E_{N,p}P_{N,p}-\varphi_{0,\Sigma_p}\|_1
  \le{}&\sup_{p\in K}\|P_{N,p}-M_{N,p}\|_{\ell^1}\\
  &+\sup_{p\in K}
    \|\mathcal E_{N,p}M_{N,p}-\varphi_{0,\Sigma_p}\|_1
  \longrightarrow0
\end{aligned}
\]
by \eqref{eq:young-multinomial-l1} and
\eqref{eq:embedded-multinomial-limit}.
\end{proof}

\begin{lemma}[Continuous limit of randomized dilation]
\label{lem:randomized-dilation}
Let \(\widetilde Q_{m_n,p}=P_{n,p}D_n\) denote the
\(\widetilde\nu\)-marginal of the joint law obtained by composing the kernel
\(D_n\) with \(P_{n,p}\), before the replacement in
\eqref{eq:rounding-kernel}.  For every compact
\(K\Subset\Delta_d^\circ\),
\begin{align}
  \sup_{p\in K}
  \left\|\mathcal E_{m_n,p}\widetilde Q_{m_n,p}
  -\varphi_{0,\gamma\Sigma_p}\right\|_1&\longrightarrow0,
  \label{eq:rounded-density-limit}\\
  \sup_{p\in K}
  \left|F(\widetilde Q_{m_n,p},P_{m_n,p})-\Fcl(\gamma,d)\right|
  &\longrightarrow0.
  \label{eq:raw-rounded-affinity}
\end{align}
\end{lemma}

\begin{proof}
For \(a>0\) and \(f\in L^1(\mathsf H)\), set
\[
  (\operatorname{dil}_a f)(y)=a^{-(d-1)/2}f(y/\sqrt a),
\]
and define the target-cell averaging operator explicitly by
\[
  (\operatorname{Av}_{m_n,p}f)(x)
  :=\frac{m_n^{(d-1)/2}}{|\mathsf C|}
   \int_{\mathsf C_{m_n,p}(\nu)}f(y)\,\dd y,
  \qquad x\in\mathsf C_{m_n,p}(\nu).
\]
Since \(m_n=\gamma_n n\), the cell definitions and
\eqref{eq:cell-kernel} give
\[
\begin{aligned}
  \gamma_n^{-1/2}\mathsf C_{m_n,p}(\nu)
  &=\frac{\gamma_n^{-1}(\nu+\mathsf C)-np}{\sqrt n},\\
  \left|\gamma_n^{-1/2}\mathsf C_{m_n,p}(\nu)
       \cap\mathsf C_{n,p}(\mu)\right|
  &=\frac{|\mathsf C|}{n^{(d-1)/2}}D_n(\nu\mid\mu).
\end{aligned}
\]
Thus, for \(x\in\mathsf C_{m_n,p}(\nu)\), changing variables
by \(z=y/\sqrt{\gamma_n}\) and summing over the source cells gives
\begin{equation}\label{eq:rounding-intertwining}
\begin{aligned}
 &\bigl[\operatorname{Av}_{m_n,p}\operatorname{dil}_{\gamma_n}
   (\mathcal E_{n,p}P_{n,p})\bigr](x)\\
 &\quad=\frac{m_n^{(d-1)/2}}{|\mathsf C|}
   \sum_{\mu\in\mathsf L_n}\frac{n^{(d-1)/2}}{|\mathsf C|}
   P_{n,p}(\mu)
   \left|\gamma_n^{-1/2}\mathsf C_{m_n,p}(\nu)
        \cap\mathsf C_{n,p}(\mu)\right|\\
 &\quad=\frac{m_n^{(d-1)/2}}{|\mathsf C|}
   \sum_{\mu\in\mathsf L_n}P_{n,p}(\mu)D_n(\nu\mid\mu)
  =(\mathcal E_{m_n,p}\widetilde Q_{m_n,p})(x).
\end{aligned}
\end{equation}
For \(f,g\in L^1(\mathsf H)\),
\[
  \|\operatorname{Av}_{m_n,p}f-\operatorname{Av}_{m_n,p}g\|_1
  \le\|f-g\|_1,
  \qquad
  \|\operatorname{dil}_{\gamma_n}f
      -\operatorname{dil}_{\gamma_n}g\|_1
  =\|f-g\|_1.
\]
By \cref{lem:young-local-limit},
\(\|\mathcal E_{n,p}P_{n,p}-\varphi_{0,\Sigma_p}\|_1=o(1)\) uniformly on
\(K\).  Since
\(\operatorname{dil}_{\gamma_n}\varphi_{0,\Sigma_p}
=\varphi_{0,\gamma_n\Sigma_p}\), dilation isometry gives
\(\|\operatorname{dil}_{\gamma_n}(\mathcal E_{n,p}P_{n,p})
-\varphi_{0,\gamma_n\Sigma_p}\|_1=o(1)\).  Moreover,
\(\|\varphi_{0,\gamma_n\Sigma_p}-\varphi_{0,\gamma\Sigma_p}\|_1=o(1)\)
uniformly because \(\gamma_n\to\gamma\).  Cell averaging satisfies
\[
  \|\operatorname{Av}_{m_n,p}f-f\|_1
  \le 2^{d-1}\!\sup_{\substack{h\in\mathsf H\\
       \|h\|\le\operatorname{diam}(\mathsf C)/\sqrt{m_n}}}
       \|f(\,\cdot+h)-f\|_1.
\]
Since \(\gamma_n\to\gamma\) and \(K\) is compact,
\(\sup_{p\in K,\,n\ge n_0}
\|\nabla\varphi_{0,\gamma_n\Sigma_p}\|_1<\infty\) for some \(n_0\).
Together with
\(\|f(\,\cdot+h)-f\|_1\le\|h\|\,\|\nabla f\|_1\), this gives
\(\|\operatorname{Av}_{m_n,p}\varphi_{0,\gamma_n\Sigma_p}
-\varphi_{0,\gamma_n\Sigma_p}\|_1=o(1)\).  Using the contraction of
\(\operatorname{Av}_{m_n,p}\) and \eqref{eq:rounding-intertwining},
\[
\begin{aligned}
  \|\mathcal E_{m_n,p}\widetilde Q_{m_n,p}
       -\varphi_{0,\gamma\Sigma_p}\|_1
  ={}&\|\operatorname{Av}_{m_n,p}\operatorname{dil}_{\gamma_n}
       (\mathcal E_{n,p}P_{n,p})
       -\varphi_{0,\gamma\Sigma_p}\|_1\\
  \le{}&\|\operatorname{Av}_{m_n,p}
       [\operatorname{dil}_{\gamma_n}(\mathcal E_{n,p}P_{n,p})
       -\varphi_{0,\gamma_n\Sigma_p}]\|_1\\
  &+\|\operatorname{Av}_{m_n,p}\varphi_{0,\gamma_n\Sigma_p}
       -\varphi_{0,\gamma_n\Sigma_p}\|_1
   +\|\varphi_{0,\gamma_n\Sigma_p}
       -\varphi_{0,\gamma\Sigma_p}\|_1\\
  \le{}&\|\operatorname{dil}_{\gamma_n}(\mathcal E_{n,p}P_{n,p})
       -\varphi_{0,\gamma_n\Sigma_p}\|_1\\
  &+\|\operatorname{Av}_{m_n,p}\varphi_{0,\gamma_n\Sigma_p}
       -\varphi_{0,\gamma_n\Sigma_p}\|_1
   +\|\varphi_{0,\gamma_n\Sigma_p}
       -\varphi_{0,\gamma\Sigma_p}\|_1=o(1).
\end{aligned}
\]
The equality is \eqref{eq:rounding-intertwining}.
This proves \eqref{eq:rounded-density-limit}.

Applying \cref{lem:young-local-limit} at sample size \(m_n\) gives
\[
  \mathcal E_{m_n,p}P_{m_n,p}\longrightarrow\varphi_{0,\Sigma_p}
\]
uniformly in \(L^1\).  The embedding preserves Hellinger affinity, and
\eqref{eq:fidelity-continuity} also holds for nonnegative \(L^1\) densities.
Therefore
\[
  F(\widetilde Q_{m_n,p},P_{m_n,p})
  \longrightarrow
  F(\mathsf N_{0,\gamma\Sigma_p},\mathsf N_{0,\Sigma_p})
  =\Fcl(\gamma,d)
\]
uniformly on \(K\).
\end{proof}

\begin{proof}[Proof of \cref{prop:young-rounding}]
For \(i<d\), the nearest-integer rule gives
\[
  |\widetilde\nu_i-\gamma_n\mu_i|
  \le\frac{\gamma_n+1}{2}.
\]
Since the coordinate sums of \(\widetilde\nu\) and \(\gamma_n\mu\) are both
\(m_n\),
\[
  |\widetilde\nu_d-\gamma_n\mu_d|
  \le\frac{(d-1)(\gamma_n+1)}2.
\]
Thus one may take
\[
  c_{\mathrm{rnd}}
  =\frac{d-1}{2}\left(1+\sup_n\gamma_n\right)<\infty.
\]

Compactness of \(K\) gives \(c_0>0\) such that, for all large \(n\), every
\(p\in K\) and \(\mu\in\mathsf T_n(p)\) satisfy
\[
  \mu_d\ge c_0n,
  \qquad
  \mu_i-\mu_{i+1}\ge c_0n\quad(i<d).
\]
For every
\(\widetilde\nu\in\operatorname{supp}D_n(\cdot\mid\mu)\),
\[
  (\widetilde\nu-\mu)_d
  \ge(\gamma_n-1)\mu_d-c_{\mathrm{rnd}},
\]
and
\[
  (\widetilde\nu_i-\mu_i)
  -(\widetilde\nu_{i+1}-\mu_{i+1})
  \ge(\gamma_n-1)(\mu_i-\mu_{i+1})-2c_{\mathrm{rnd}}.
\]
These quantities are positive for all large \(n\), uniformly on \(K\).
Hence \(\widetilde\nu-\mu\) is a partition; since \(\mu\) is a partition,
so is \(\widetilde\nu\).
Moreover, \(\widetilde\nu_d>\mu_d>0\), so
\(\widetilde\nu\ne\nu_n^{\mathrm{fb}}\).  Hence, for
\(\mu\in\mathsf T_n(p)\), the law \(D_n(\cdot\mid\mu)\) is supported on
\(\mathsf Y_{m_n}\), assigns no mass to \(\nu_n^{\mathrm{fb}}\), and
undergoes no replacement in \eqref{eq:rounding-kernel}.  Consequently,
\[
  K_n^{\mathrm{rnd}}(\nu\mid\mu)=D_n(\nu\mid\mu),
  \qquad \nu\in\mathsf Y_{m_n},
\]
and every \(\nu\) in this support satisfies
\(\nu-\mu\vdash m_n-n\) and
\(\|\nu-\gamma_n\mu\|_\infty\le c_{\mathrm{rnd}}\).

Let \(\eps_{n,p}\) be the probability, under \(\mu\sim P_{n,p}\) and
\(\widetilde\nu\sim D_n(\cdot\mid\mu)\), that \(\widetilde\nu\) is
replaced in \eqref{eq:rounding-kernel}.  The preceding argument gives
\[
  \eps_{n,p}\le P_{n,p}(\mathsf T_n(p)^c),
\]
which tends uniformly to zero by \cref{lem:young-concentration}.

By \eqref{eq:raw-rounded-affinity}, the unreplaced output law has affinity
\(\Fcl(\gamma,d)+o(1)\) with \(P_{m_n,p}\), uniformly on \(K\).  Moving the
rejected mass to \(\nu_n^{\mathrm{fb}}\) changes the law by at most
\(2\eps_{n,p}\) in \(\ell^1\).  Applying \eqref{eq:fidelity-continuity} to
classical distributions therefore gives
\eqref{eq:rounding-label-fidelity}.
\end{proof}

\section{Unified Gaussian flat-prior converse}
\label[appendix]{app:gaussian-converse}

This appendix proves the two converse inequalities in
\cref{thm:gaussian-amplification} by averaging any candidate channel
\(\Lambda\) over
larger parameter boxes and testing the unavoidable noise in its output.
\Cref{lem:weighted-fidelity} turns fidelity into such a test expectation.
\Cref{lem:compact-cp-limit} shows that any bounded sequence of
nearly displacement-covariant maps has a subsequence converging to an exactly
covariant map with the same trace bound.
\Cref{prop:least-noise-moment} shows that every covariant map with
the required gain is at least as noisy as the amplifier, as measured by the
number test, even allowing for lost trace.  \Cref{lem:scalar-convolution}
governs the classical-distribution part of a classical--quantum channel:
translation covariance forces it to scale \(x\) to \(\sqrt\gamma x\), add
independent noise, and possibly lose mass.
\Cref{prop:flat-prior-converse} combines these bounds, giving
\(\Forb(\gamma,p)\) for
\(\Phi_z^{(p)}\mapsto\Psi_z^{(p,\gamma)}\) and
\(\Fcl(\gamma,d)\Forb(\gamma,p)\) for
\(\Theta_{h,z}^{\mathrm{in}}\mapsto\Theta_{h,z}^{\mathrm{tar}}\).

The classical parameter space has
dimension \(k=0\) for the orbital problem and \(k=d-1\) for the full problem.
All trace-class-valued integrals below are Bochner integrals in trace norm.

\subsection{Weighted fidelity and compact limits}

\begin{lemma}[Weighted root-fidelity bound]\label{lem:weighted-fidelity}
Let \(R,T\) be positive trace-class operators, and let \(W\) be bounded,
positive, and injective.  Define
\[
  \Tr(TW^{-1})
  :=\lim_{\varepsilon\downarrow0}
  \Tr\bigl(T(W+\varepsilon I)^{-1}\bigr)\in[0,\infty].
\]
If this quantity is finite, then
\begin{equation}\label{eq:weighted-quantum-fidelity}
  F(R,T)^2\le\Tr(RW)\Tr(TW^{-1}).
\end{equation}

Let \(\mathsf E\) be a finite-dimensional Euclidean space,
let \(R\in L^1(\mathsf E;\cT_1(\cF_{\mathrm{orb}}))\) be positive, let
\(g\) be a probability density on \(\mathsf E\), and let \(\chi>0\) be
measurable.  For positive \(R,S\) in this \(L^1\)-space, set
\(F(R,S)=\int_{\mathsf E}F(R(y),S(y))\,\dd y\), extending the hybrid-state
convention in \cref{sec:gaussian-LAN}.
The notation \(gT\) means the density \(y\mapsto g(y)T\).
Whenever the right-hand side below is finite,
\begin{equation}\label{eq:weighted-hybrid-fidelity}
  F(R,gT)^2
  \le
  \left(\int_{\mathsf E}\chi(y)\Tr[R(y)W] \,\dd y\right)
  \left(\int_{\mathsf E}\frac{g(y)}{\chi(y)}\,\dd y\right)
  \Tr(TW^{-1}).
\end{equation}
\end{lemma}

\begin{proof}
For \(W_\varepsilon=W+\varepsilon I\), the Schatten H\"older inequality gives
\[
  F(R,T)
  =\|(R^{1/2}W_\varepsilon^{1/2})(W_\varepsilon^{-1/2}T^{1/2})\|_1
  \le\sqrt{\Tr(RW_\varepsilon)}\sqrt{\Tr(TW_\varepsilon^{-1})}.
\]
Squaring and letting \(\varepsilon\downarrow0\), using
\(\Tr(RW_\varepsilon)=\Tr(RW)+\varepsilon\Tr R\) and monotone convergence
for the second trace, proves \eqref{eq:weighted-quantum-fidelity}.  For the
hybrid state,
\[
  F(R,gT)
  =\int_{\mathsf E}\sqrt{g(y)}\,F(R(y),T)\,\dd y.
\]
Apply the first inequality pointwise and then Cauchy--Schwarz to the factors
\(\sqrt{\chi(y)\Tr[R(y)W]}\) and
\(\sqrt{g(y)/\chi(y)}\).  This gives
\eqref{eq:weighted-hybrid-fidelity}.
\end{proof}

\begin{lemma}[Compact-observable limits of completely positive maps]
\label{lem:compact-cp-limit}
Let \(\cH_0,\cH_1\) be separable Hilbert spaces and let
\(\Gamma_j:\cT_1(\cH_0)\to\cT_1(\cH_1)\) be completely positive maps with
\(\sup_j\|\Gamma_j\|_{1\to1}<\infty\).  There are a subsequence and a
completely positive map
\(\Gamma:\cT_1(\cH_0)\to\cT_1(\cH_1)\) such that
\begin{equation}\label{eq:compact-map-convergence}
  \Tr[\Gamma_j(X)K]\longrightarrow\Tr[\Gamma(X)K]
\end{equation}
for every \(X\in\cT_1(\cH_0)\) and every compact
\(K\in\mathcal K(\cH_1)\).  If
\(\Tr\Gamma_j(X)\le c\Tr X\) for every \(X\ge0\), then
\(\Tr\Gamma(X)\le c\Tr X\).

Suppose in addition that
\(U_g\in\U(\cH_0)\) and \(V_g\in\U(\cH_1)\) are families of
unitaries indexed by a set \(G\), and that, for every fixed \(g\in G\) and
\(X\in\cT_1(\cH_0)\),
\[
  \|\Gamma_j(U_g X U_g^*)-V_g\Gamma_j(X)V_g^*\|_1
  \longrightarrow0.
\]
Then
\[
  \Gamma(U_g X U_g^*)=V_g\Gamma(X)V_g^*
\]
for every \(g\in G\) and \(X\in\cT_1(\cH_0)\).
\end{lemma}

\begin{proof}
Let \(M=\sup_j\|\Gamma_j\|_{1\to1}\), and let
\(\mathcal D\subset\mathcal K(\cH_1)\) be a countable norm-dense
\(\mathbb Q(\mathrm i)\)-linear \(*\)-algebra of finite-rank operators.
The predual \(\cT_1(\cH_0)\) is separable, so a diagonal weak-* compactness
argument gives a subsequence for which
\(\Gamma_j^*(K)\) converges weak-* in \(\mathcal B(\cH_0)\) for every
\(K\in\mathcal D\).  The bound
\(\|\Gamma_j^*(K)\|\le M\|K\|\) extends both the limit and the convergence
to every compact \(K\), producing a
bounded complex-linear map
\(\Gamma^*:\mathcal K(\cH_1)\to\mathcal B(\cH_0)\).

To check complete positivity, fix \(r\) and \(A\ge0\) in
\(M_r(\mathcal K(\cH_1))\).  Since \(A^{1/2}\) is compact, choose
\(B_\ell\in M_r(\mathcal D)\) with \(B_\ell\to A^{1/2}\).  Since
\(\Gamma_j^{*(r)}(B_\ell^*B_\ell)\ge0\), first letting \(j\to\infty\) and
then \(\ell\to\infty\) gives \(\Gamma^{*(r)}(A)\ge0\).

Since \(\mathcal K(\cH_1)^*=\cT_1(\cH_1)\), define \(\Gamma(X)\) by
\[
  \Tr[\Gamma(X)K]=\Tr[X\Gamma^*(K)]
  \qquad(K\in\mathcal K(\cH_1)).
\]
Then \(\Gamma\) is completely positive and
\eqref{eq:compact-map-convergence} holds.  If \(P_N\uparrow I\) are
finite-rank projections and \(X\ge0\), then
\[
  \Tr\Gamma(X)
  =\sup_N\lim_j\Tr[\Gamma_j(X)P_N]
  \le c\Tr X.
\]
Finally, for fixed \(g,X\) and compact \(K\), the operator
\(V_g^*KV_g\) is compact, and hence
\[
\begin{aligned}
  &\Tr\!\left[
    \bigl(\Gamma(U_gXU_g^*)-V_g\Gamma(X)V_g^*\bigr)K
  \right]\\
  &\qquad=\lim_j\Tr\!\left[
    \bigl(\Gamma_j(U_gXU_g^*)-V_g\Gamma_j(X)V_g^*\bigr)K
  \right]=0.
\end{aligned}
\]
Compact operators separate trace-class operators, so the limiting map is
covariant.
\end{proof}

\subsection{The least-noise oscillator moment}

Fix \(\gamma>1\).  For \(s\ge1\) and
\(q_1,\ldots,q_s\in(0,1)\), put
\[
  q_\ell^{(\gamma)}=1-\frac{1-q_\ell}{\gamma},
  \qquad
  \tau=\bigotimes_{\ell=1}^s\tau_{q_\ell},
  \qquad
  \tau^{(\gamma)}
  =\bigotimes_{\ell=1}^s\tau_{q_\ell^{(\gamma)}}.
\]

The proof below uses the following two external inputs, isolated here to make
the dependence on \cite{GutaMatsumoto2006MixedGaussianCloning} explicit.

\begin{fact}[Covariant-amplifier representation]
\label{fact:covariant-amplifier-representation}
Every displacement-covariant one-mode channel of amplitude gain
\(\sqrt\gamma\) admits an idler realization
\[
  a_{\mathrm{out}}
  =\sqrt\gamma\,a_{\mathrm{in}}+\sqrt{\gamma-1}\,b^*.
\]
The idler state need not be Gaussian; if the channel is also phase covariant,
it may be chosen number diagonal
\cite[Sec.~II.C and Appendix; Sec.~III for general gain]
{GutaMatsumoto2006MixedGaussianCloning}.
\end{fact}

\begin{fact}[Vacuum-idler stochastic order]
\label{fact:vacuum-idler-stochastic-order}
For a thermal input \(\tau_q\) and a number-diagonal idler, let \(p_k\) be
the output number distribution of the amplifier in
\cref{fact:covariant-amplifier-representation}.  Then
\[
  \sum_{k=0}^m(1-q^{(\gamma)})(q^{(\gamma)})^k
  \ge \sum_{k=0}^m p_k
  \qquad(m\ge0).
\]
Thus the vacuum-idler output \(\tau_{q^{(\gamma)}}\) is stochastically
smallest among these outputs
\cite[Lem.~2.3 and Sec.~III]{GutaMatsumoto2006MixedGaussianCloning}.
\end{fact}

\begin{proposition}[Least-noise moment]
\label{prop:least-noise-moment}
Let \(\Gamma\) be a completely positive trace-nonincreasing map on the
\(s\)-mode trace class satisfying
\[
  \Gamma(D(z)XD(z)^*)
  =D(\sqrt\gamma\,z)\Gamma(X)D(\sqrt\gamma\,z)^*
  \qquad(z\in\mathbb C^s).
\]
If
\[
  W(\widehat{\mathbf N})
  =\prod_{\ell=1}^s w_\ell(\widehat N_\ell),
\]
where every \(w_\ell\) is bounded, nonnegative, and nonincreasing on
\(\mathbb Z_{\ge0}\), then
\begin{equation}\label{eq:least-noise-moment}
  \Tr[\Gamma(\tau)W(\widehat{\mathbf N})]
  \le\Tr[\Gamma(\tau)]
       \Tr[\tau^{(\gamma)}W(\widehat{\mathbf N})].
\end{equation}
\end{proposition}

\begin{proof}
Averaging \(\Gamma\) over independent phase rotations preserves its
trace bound and displacement covariance, and it does not change either side
of \eqref{eq:least-noise-moment}.  We may therefore assume that \(\Gamma\)
is phase covariant in every mode.

Consider one mode.  Write \(\Tr\Gamma(X)=\Tr(AX)\) with \(0\le A\le I\).
Displacement covariance implies \(D(z)^*AD(z)=A\); hence irreducibility of
the Weyl representation gives \(A=cI\), where
  \(c=\Tr\Gamma(\tau_q)\).  If \(c>0\), then
  \(\Lambda=c^{-1}\Gamma\) is a phase-covariant amplifier of amplitude gain
  \(\sqrt\gamma\).  By
  \cref{fact:covariant-amplifier-representation}, it has a number-diagonal
  idler realization.  \Cref{fact:vacuum-idler-stochastic-order} therefore
  gives
\begin{equation}\label{eq:imported-stochastic-order}
  \sum_{n=0}^m(1-q^{(\gamma)})(q^{(\gamma)})^n
  \ge
  \sum_{n=0}^m\langle n|\Lambda(\tau_q)|n\rangle
  \qquad(m\ge0).
\end{equation}
The distribution on the left is the vacuum-idler output, namely that of
\(\tau_{q^{(\gamma)}}\), and both output states are number diagonal.  For a
bounded nonincreasing \(w\), summation by parts writes its expectation as a
nonnegative linear combination of these cumulative probabilities.  Both
distributions in \eqref{eq:imported-stochastic-order} have total mass one,
so the common limiting constant cancels.  Hence
\[
  \Tr[\Lambda(\tau_q)w(\widehat N)]
  \le\Tr[\tau_{q^{(\gamma)}}w(\widehat N)].
\]
Multiplication by \(c\) proves the one-mode case; the case \(c=0\) is
immediate.

For the multimode case, the claim is immediate if some \(w_\ell\)
vanishes identically; otherwise, divide each \(w_\ell\) by
\(w_\ell(0)\) and assume \(0\le w_\ell\le1\).  Put
\[
  W_{\ge\ell}
  =I_1\ot\cdots\ot I_{\ell-1}\ot
   \bigotimes_{j=\ell}^s w_j(\widehat N_j),
  \qquad W_{\ge s+1}=I,
\]
and, for an operator \(X\) on mode \(\ell\), define
\[
  \Gamma_\ell(X)
  =\Tr_{\ne\ell}\!\left[
    W_{\ge\ell+1}^{1/2}
    \Gamma(\tau_{q_1}\ot\cdots\ot\tau_{q_{\ell-1}}\ot X\ot
           \tau_{q_{\ell+1}}\ot\cdots\ot\tau_{q_s})
    W_{\ge\ell+1}^{1/2}
  \right].
\]
Because \(0\le W_{\ge\ell+1}\le I\) and this operator is the identity on
mode \(\ell\), each \(\Gamma_\ell\) is completely positive, trace
nonincreasing, and displacement covariant in that mode.  The phase-covariant
reduction following \eqref{eq:least-noise-moment} also makes \(\Gamma_\ell\)
phase covariant because the fixed thermal inputs and \(W_{\ge\ell+1}\) are
phase invariant.  Moreover,
\[
  \Tr[\Gamma_\ell(\tau_{q_\ell})w_\ell(\widehat N_\ell)]
  =\Tr[\Gamma(\tau)W_{\ge\ell}],
  \qquad
  \Tr\Gamma_\ell(\tau_{q_\ell})
  =\Tr[\Gamma(\tau)W_{\ge\ell+1}].
\]
Applying the one-mode inequality to \(\Gamma_\ell\) and iterating over
\(\ell=1,\ldots,s\) gives \eqref{eq:least-noise-moment}.
\end{proof}

Specialize to \(s=r_d\), \(q_\ell=q_{ij}\),
\(\tau=\tau_p\), and \(\tau^{(\gamma)}=\tau_p^{(\gamma)}\).
For
\(\mathbf k=(k_{ij})_{i<j}\in\mathbb Z_{\ge0}^{r_d}\), let
\(t_{\mathbf k}\) and \(t_{\mathbf k}^{(\gamma)}\) be the corresponding
number-basis eigenvalues, and define
\begin{equation}\label{eq:quantum-witness}
  W_\gamma
  =\sum_{\mathbf k}
  \sqrt{\frac{t_{\mathbf k}}{t_{\mathbf k}^{(\gamma)}}}
  \proj{\mathbf k}.
\end{equation}
Equivalently,
\[
  W_\gamma=\prod_{i<j}w_{ij}(\widehat N_{ij}),
  \qquad
  w_{ij}(n)=
  \sqrt{\frac{1-q_{ij}}{1-q_{ij}^{(\gamma)}}}
  \left(\frac{q_{ij}}{q_{ij}^{(\gamma)}}\right)^{n/2}.
\]
Here \(\widehat N_{ij}\) is the number operator of the mode indexed by
\(i<j\), and \(\lvert\mathbf k\rvert:=\sum_{i<j}k_{ij}\) is the total
occupation number.
Since \(q_{ij}<q_{ij}^{(\gamma)}\), each \(w_{ij}(n)\) is positive, bounded,
and strictly decreasing in \(n\), and the eigenvalues of \(W_\gamma\) tend
to zero as \(\lvert\mathbf k\rvert\to\infty\).
Thus \(W_\gamma\) is bounded,
injective, and compact.  Its definition gives
\begin{equation}\label{eq:quantum-witness-identities}
  \Tr(\tau_p^{(\gamma)}W_\gamma)
  =\Tr(\tau_p W_\gamma^{-1})
  =\sum_{\mathbf k}\sqrt{t_{\mathbf k}t_{\mathbf k}^{(\gamma)}}
  =\Forb(\gamma,p),
\end{equation}
where the inverse is understood as in \cref{lem:weighted-fidelity}.

\subsection{Classical trace kernels}

\begin{lemma}[Translation-covariant scalar kernels]
\label{lem:scalar-convolution}
Let \(\mathsf E\) be a finite-dimensional Euclidean space and let
\(\mathcal R:L^1(\mathsf E)\to C_0(\mathsf E)^*\) be a positive
contraction.  For a finite Radon measure \(\lambda\), define
\(T_a\lambda(A)=\lambda(A-a)\) for Borel \(A\).  Suppose
\begin{equation}\label{eq:scalar-measure-covariance}
  \mathcal R(u(\,\cdot-h))=T_{\sqrt\gamma h}\mathcal R(u)
  \qquad(u\in L^1(\mathsf E),\ h\in\mathsf E)
\end{equation}
as an equality of finite Radon measures.  Then there is a finite positive
measure \(\mu\), with \(\mu(\mathsf E)\le1\), such that
\begin{equation}\label{eq:scalar-convolution}
  \mathcal R(u)(A)
  =\int_{\mathsf E}u(x)\,\mu(A-\sqrt\gamma x)\,\dd x
\end{equation}
for every \(u\in L^1(\mathsf E)\) and every Borel set \(A\).
\end{lemma}

\begin{proof}
If \(\mathsf E=\{0\}\), take \(\mu=\mathcal R(1)\).  Otherwise choose a
nonnegative approximate identity \(u_j\in C_c(\mathsf E)\) with \(\int u_j=1\)
and \(\operatorname{supp}u_j\subset\{x:\|x\|\le j^{-1}\}\).  The positive
measures \(\mathcal R(u_j)\) lie in the unit ball of \(C_0(\mathsf E)^*\), so
weak-* compactness and separability yield a positive subsequential weak-* limit
\(\mu\) with \(\mu(\mathsf E)\le1\).

For \(u\in L^1(\mathsf E)\), the identity
\(u*u_j=\int_{\mathsf E}u(x)u_j(\,\cdot-x)\,\dd x\) holds as a Bochner
integral in \(L^1(\mathsf E)\), so bounded linearity and covariance give
\[
  \mathcal R(u*u_j)=\int_{\mathsf E}u(x)T_{\sqrt\gamma x}\mathcal R(u_j)\,\dd x.
\]
Pair with \(\phi\in C_0(\mathsf E)\).  The approximate-identity property and
contractivity give \(\mathcal R(u*u_j)\to\mathcal R(u)\) in total variation.
On the right, weak-* convergence applies pointwise because
\(y\mapsto\phi(y+\sqrt\gamma x)\) lies in \(C_0(\mathsf E)\), and the
integrand is bounded by \(|u(x)|\|\phi\|_\infty\).  Dominated convergence gives
\[
  \langle\phi,\mathcal R(u)\rangle=\int_{\mathsf E}u(x)\langle\phi,T_{\sqrt\gamma x}\mu\rangle\,\dd x.
\]
By definition of \(T_a\), the right-hand side pairs \(\phi\) with the finite
Radon measure in \eqref{eq:scalar-convolution}, whose total variation is at
most \(\|u\|_1\mu(\mathsf E)\).  Riesz uniqueness proves the claim.
\end{proof}

\subsection{One flat-prior converse for both Gaussian models}

\begin{proposition}[Flat-prior converse inequality]
\label{prop:flat-prior-converse}
Let \(k\in\{0,d-1\}\).  For \(k=0\), let
\(\mathsf E=\{0\}\), omit the classical factor, and interpret \(h\)-integration
as evaluation at \(0\).  For \(k=d-1\), let \(\mathsf E=\mathsf H\) with its
induced Euclidean measure.  Let
\[
  \mathsf E_L=\{h\in\mathsf E:\|h\|_\infty\le L\},
  \qquad
  \Xi_L=\mathsf E_L\times\mathsf Z_L,
\]
with \(|\mathsf E_L|=1\) when \(k=0\).  Define
\[
  R_{0,z}^{(0)}=\Phi_z^{(p)},
  \qquad S_{0,z}^{(0)}=\Psi_z^{(p,\gamma)},
\]
\[
  R_{h,z}^{(d-1)}=\Theta_{h,z}^{\mathrm{in}},
  \qquad S_{h,z}^{(d-1)}=\Theta_{h,z}^{\mathrm{tar}},
  \qquad
  c_k=\left(\frac{2\sqrt\gamma}{1+\gamma}\right)^{k/2}.
\]
For \(k=0\), the supremum below is over quantum channels on
\(\cF_{\mathrm{orb}}\); for \(k=d-1\), it is over classical--quantum
channels.  Then
\begin{equation}\label{eq:unified-flat-prior-converse}
  \limsup_{L\to\infty}\sup_\Lambda
  \frac1{|\Xi_L|}\int_{\Xi_L}
  F\!\left(\Lambda(R_{h,z}^{(k)}),S_{h,z}^{(k)}\right)
  \,\dd h\,\dd z
  \le c_k\Forb(\gamma,p).
\end{equation}
\end{proposition}

\begin{proof}
For \(k=0\), set \(f_1=f_\gamma=\chi=1\) on the one-point space.  For
\(k=d-1\), let \(f_1\) and \(f_\gamma\) be the centered Gaussian densities
with covariances \(\Sigma_p\) and \(\gamma\Sigma_p\), and set
\[
  \chi(y)=\sqrt{\frac{f_1(y)}{f_\gamma(y)}}
  =\gamma^{k/4}
   \exp\!\left[-\frac{\gamma-1}{4\gamma}
     \langle y,\Sigma_p^{-1}y\rangle\right].
\]
This function is bounded and strictly positive; for \(k=d-1\), it also
belongs to \(C_0(\mathsf E)\).  In both cases,
\begin{equation}\label{eq:classical-witness-identity}
  \int_{\mathsf E}\frac{f_1(y)}{\chi(y)}\,\dd y
  =\int_{\mathsf E}\sqrt{f_1(y)f_\gamma(y)}\,\dd y
  =c_k.
\end{equation}

For \(\xi=(h,z)\), define
\[
  (\mathcal U_\xi^{\mathrm{in}}R)(x)
  =D(z)R(x-h)D(z)^*,
  \qquad
  (\mathcal U_\xi^{\mathrm{out}}R)(y)
  =D(\sqrt\gamma z)R(y-\sqrt\gamma h)D(\sqrt\gamma z)^*.
\]
Then
\(R_\xi^{(k)}=\mathcal U_\xi^{\mathrm{in}}(f_1\tau_p)\) and
\(S_\xi^{(k)}=\mathcal U_\xi^{\mathrm{out}}(f_1\tau_p)\).  For a channel
\(\Lambda\), let
\[
  \overline\Lambda_L
  =\frac1{|\Xi_L|}\int_{\Xi_L}
   \mathcal U_{-\eta}^{\mathrm{out}}\circ\Lambda\circ
   \mathcal U_\eta^{\mathrm{in}}\,\dd\eta.
\]
Invariance of \(F\) under the isometries
\(\mathcal U_\xi^{\mathrm{out}}\) and joint concavity, applied to the Bochner
average, give
\begin{equation}\label{eq:unified-folner-concavity}
  F_L(\Lambda)
  =\frac1{|\Xi_L|}\int_{\Xi_L}
  F\!\left(\Lambda(R_\xi^{(k)}),S_\xi^{(k)}\right)\,\dd\xi
  \le F(\overline\Lambda_L(f_1\tau_p),f_1\tau_p).
\end{equation}

Let \(F_*=\limsup_{L\to\infty}\sup_\Lambda F_L(\Lambda)\).  Choose
\(L_j\to\infty\) with \(\sup_\Lambda F_{L_j}(\Lambda)\to F_*\), and channels
\(\widetilde\Lambda_j\) satisfying
\(F_{L_j}(\widetilde\Lambda_j)\ge\sup_\Lambda F_{L_j}(\Lambda)-j^{-1}\).
Set
\[
  \Lambda_j=\overline{\widetilde\Lambda_j}_{L_j},
  \qquad R_j=\Lambda_j(f_1\tau_p).
\]
Then
\begin{equation}\label{eq:unified-maximizing-sequence}
  F_{L_j}(\widetilde\Lambda_j)\longrightarrow F_*,
  \qquad
  F_{L_j}(\widetilde\Lambda_j)\le F(R_j,f_1\tau_p).
\end{equation}
For \(\eta\in\mathsf E\times\mathbb C^{r_d}\), write
\[
  A_{j,\eta}
  =\mathcal U_{-\eta}^{\mathrm{out}}\circ
   \widetilde\Lambda_j\circ\mathcal U_\eta^{\mathrm{in}},
  \qquad
  \Lambda_j=\frac1{|\Xi_{L_j}|}\int_{\Xi_{L_j}}A_{j,\eta}\,\dd\eta.
\]
The group law and the change of variables \(u=\eta+\xi\) give
\[
\begin{aligned}
  \Lambda_j(\mathcal U_\xi^{\mathrm{in}}R)
  &=\mathcal U_\xi^{\mathrm{out}}
    \frac1{|\Xi_{L_j}|}\int_{\Xi_{L_j}+\xi}A_{j,u}(R)\,\dd u,\\
  \mathcal U_\xi^{\mathrm{out}}\Lambda_j(R)
  &=\mathcal U_\xi^{\mathrm{out}}
    \frac1{|\Xi_{L_j}|}\int_{\Xi_{L_j}}A_{j,u}(R)\,\dd u.
\end{aligned}
\]
The two integrals cancel on the intersection of the two averaging sets,
and \(\|A_{j,u}\|_{1\to1}\le1\).  Therefore, for fixed \(\xi\),
\begin{equation}\label{eq:unified-approx-covariance}
  \|\Lambda_j(\mathcal U_\xi^{\mathrm{in}}R)
    -\mathcal U_\xi^{\mathrm{out}}\Lambda_j(R)\|_1
  \le
  \frac{|(\Xi_{L_j}+\xi)\mathbin\triangle\Xi_{L_j}|}{|\Xi_{L_j}|}
  \|R\|_1
  \longrightarrow0.
\end{equation}
Indeed, \(\Xi_1\) is a bounded polytope and \(\Xi_L=L\Xi_1\), so the
ratio is \(O_\xi(L^{-1})\).

Define
\[
  \Gamma_j^\chi(X)
  =\int_{\mathsf E}\chi(y)[\Lambda_j(f_1X)](y)\,\dd y.
\]
These maps are completely positive and satisfy
\(\|\Gamma_j^\chi\|_{1\to1}\le\|\chi\|_\infty\).  Taking
\(\xi=(0,z)\) in \eqref{eq:unified-approx-covariance}, multiplying by
\(\chi\), and integrating gives
\[
\begin{aligned}
 &\|\Gamma_j^\chi(D(z)XD(z)^*)
   -D(\sqrt\gamma z)\Gamma_j^\chi(X)D(\sqrt\gamma z)^*\|_1\\
 &\qquad\le\|\chi\|_\infty
   \frac{|(\Xi_{L_j}+(0,z))\mathbin\triangle\Xi_{L_j}|}{|\Xi_{L_j}|}
   \|X\|_1\longrightarrow0.
\end{aligned}
\]
Hence \cref{lem:compact-cp-limit} gives, after passing to a subsequence, a
completely positive displacement-covariant limit \(\Gamma^\chi\) such that
\begin{equation}\label{eq:weighted-map-limit}
  \Tr[\Gamma_j^\chi(X)K]
  \longrightarrow\Tr[\Gamma^\chi(X)K]
\end{equation}
for trace-class \(X\) and compact \(K\).  With
\[
  M_j
  =\int_{\mathsf E}\chi(y)\Tr[R_j(y)W_\gamma] \,\dd y
  =\Tr[\Gamma_j^\chi(\tau_p)W_\gamma],
\]
compactness of \(W_\gamma\) gives
\begin{equation}\label{eq:moment-limit}
  M_j\longrightarrow\Tr[\Gamma^\chi(\tau_p)W_\gamma].
\end{equation}
Moreover, \cref{lem:weighted-fidelity},
\eqref{eq:classical-witness-identity}, and
\eqref{eq:quantum-witness-identities} give
\begin{equation}\label{eq:unified-fidelity-to-moment}
  F(R_j,f_1\tau_p)^2
  \le M_j\,c_k\Forb(\gamma,p).
\end{equation}

If \(k=0\), then \(\Gamma_j^\chi=\Lambda_j\) are channels.  The trace-bound
clause of \cref{lem:compact-cp-limit} with \(c=1\) makes the limit trace
nonincreasing.  By
\cref{prop:least-noise-moment},
\[
\begin{aligned}
  \lim_j M_j
  &=\Tr[\Gamma^\chi(\tau_p)W_\gamma]\\
  &\le\Tr\Gamma^\chi(\tau_p)\,
       \Tr(\tau_p^{(\gamma)}W_\gamma)
  \le\Forb(\gamma,p).
\end{aligned}
\]
Combining \eqref{eq:unified-maximizing-sequence} with
\eqref{eq:unified-fidelity-to-moment} yields
\(F_*\le\Forb(\gamma,p)\).

Assume henceforth that \(k=d-1\).  For \(u\in L^1(\mathsf E)\), define
\[
  \mathcal R_j(u)(A)
  =\int_A\Tr\bigl([\Lambda_j(u\tau_p)](y)\bigr)\,\dd y.
\]
The maps \(\mathcal R_j:L^1(\mathsf E)\to C_0(\mathsf E)^*\) are positive
contractions.  Using countable dense subsets of \(L^1(\mathsf E)\) and
\(C_0(\mathsf E)\), pass to a further subsequence converging weak-* on every \(u\)
to a positive contraction \(\mathcal R\).  Taking \(z=0\) in
\eqref{eq:unified-approx-covariance}, pairing with a function in
\(C_0(\mathsf E)\), and passing to the limit gives
\[
  \mathcal R(u(\,\cdot-h))
  =T_{\sqrt\gamma h}\mathcal R(u).
\]
Let \(\mu\) be the measure from \cref{lem:scalar-convolution}; then
\(\mu(\mathsf E)\le1\).

Finite-rank projections, \eqref{eq:weighted-map-limit}, the weak-*
convergence of \(\mathcal R_j(f_1)\), and
\(\chi\in C_0(\mathsf E)\) give
\[
\begin{aligned}
  \Tr\Gamma^\chi(\tau_p)
  &\le\liminf_j\Tr\Gamma_j^\chi(\tau_p)\\
  &=\lim_j\int_{\mathsf E}\chi(y)\,\mathcal R_j(f_1)(\dd y)
   =\int_{\mathsf E}\chi(y)\,\mathcal R(f_1)(\dd y).
\end{aligned}
\]
By \eqref{eq:scalar-convolution}, the last integral equals
\[
  \int_{\mathsf E}\mu(\dd r)
  \int_{\mathsf E}f_1(x)\chi(r+\sqrt\gamma x)\,\dd x.
\]
Completing the square gives
\[
  \int_{\mathsf E}f_1(x)\chi(r+\sqrt\gamma x)\,\dd x
  =c_k\exp\!\left[-\frac{\gamma-1}{2\gamma(\gamma+1)}
    \langle r,\Sigma_p^{-1}r\rangle\right].
\]
Since \(\mu(\mathsf E)\le1\), it follows that
\(\Tr\Gamma^\chi(\tau_p)\le c_k\).

Finally write \(\Tr\Gamma^\chi(X)=\Tr(AX)\).  Covariance implies that \(A\)
commutes with every Weyl displacement, so \(A=cI\), with
\[
  c=\Tr\Gamma^\chi(\tau_p)\le c_k\le1.
\]
Thus \(\Gamma^\chi\) is trace nonincreasing, and
\cref{prop:least-noise-moment} gives
\[
\begin{aligned}
  \lim_j M_j
  &=\Tr[\Gamma^\chi(\tau_p)W_\gamma]\\
  &\le\Tr\Gamma^\chi(\tau_p)\,
       \Tr(\tau_p^{(\gamma)}W_\gamma)
  \le c_k\Forb(\gamma,p).
\end{aligned}
\]
Combining this with \eqref{eq:unified-maximizing-sequence} and
\eqref{eq:unified-fidelity-to-moment} yields
\[
  F_*\le c_k\Forb(\gamma,p),
\]
which proves the full classical--quantum case.
\end{proof}


\begin{thebibliography}{BCFJS96}
\bibitem[WZ82]{WoottersZurek1982}
W. K. Wootters and W. H. Zurek,
A single quantum cannot be cloned,
\doi{10.1038/299802a0}{\emph{Nature} \textbf{299}, 802--803 (1982)}.

\bibitem[Die82]{Dieks1982}
D. Dieks,
Communication by EPR devices,
\doi{10.1016/0375-9601(82)90084-6}{\emph{Physics Letters A} \textbf{92}, 271--272 (1982)}.

\bibitem[BH96]{BuzekHillery1996}
V. Buzek and M. Hillery,
Quantum copying: Beyond the no-cloning theorem,
\doi{10.1103/PhysRevA.54.1844}{\emph{Phys. Rev. A} \textbf{54}, 1844--1852 (1996)}.

\bibitem[GM97]{GisinMassar1997}
N. Gisin and S. Massar,
Optimal quantum cloning machines,
\doi{10.1103/PhysRevLett.79.2153}{\emph{Phys. Rev. Lett.} \textbf{79}, 2153--2156 (1997)}.

\bibitem[BDE+98]{BrussDiVincenzoEkertFuchsMacchiavelloSmolin1998}
D. Bruss, D. P. DiVincenzo, A. Ekert, C. A. Fuchs, C. Macchiavello, and J. A. Smolin,
Optimal universal and state-dependent quantum cloning,
\doi{10.1103/PhysRevA.57.2368}{\emph{Phys. Rev. A} \textbf{57}, 2368--2378 (1998)},
\arxiv{quant-ph/9705038}.

\bibitem[CB99]{CheflesBarnett1999}
A. Chefles and S. M. Barnett,
Strategies and networks for state-dependent quantum cloning,
\doi{10.1103/PhysRevA.60.136}{\emph{Phys. Rev. A} \textbf{60}, 136--144 (1999)},
\arxiv{quant-ph/9812035}.

\bibitem[Wer98]{Werner1998}
R. F. Werner,
Optimal cloning of pure states,
\doi{10.1103/PhysRevA.58.1827}{\emph{Phys. Rev. A} \textbf{58}, 1827--1832 (1998)},
\arxiv{quant-ph/9804001}.

\bibitem[KW99]{KeylWerner1999}
M. Keyl and R. F. Werner,
Optimal cloning of pure states, testing single clones,
\doi{10.1063/1.532887}{\emph{J. Math. Phys.} \textbf{40}, 3283--3299 (1999)},
\arxiv{quant-ph/9807010}.

\bibitem[MP95]{MassarPopescu1995}
S. Massar and S. Popescu,
Optimal extraction of information from finite quantum ensembles,
\doi{10.1103/PhysRevLett.74.1259}{\emph{Phys. Rev. Lett.} \textbf{74}, 1259--1263 (1995)}.

\bibitem[BEM98]{BrussEkertMacchiavello1998}
D. Bruss, A. Ekert, and C. Macchiavello,
Optimal universal quantum cloning and state estimation,
\doi{10.1103/PhysRevLett.81.2598}{\emph{Phys. Rev. Lett.} \textbf{81}, 2598--2601 (1998)}.

\bibitem[BA06]{BaeAcin2006}
J. Bae and A. Acin,
Asymptotic quantum cloning is state estimation,
\doi{10.1103/PhysRevLett.97.030402}{\emph{Phys. Rev. Lett.} \textbf{97}, 030402 (2006)},
\arxiv{quant-ph/0603078}.

\bibitem[CD06]{ChiribellaDAriano2006}
G. Chiribella and G. M. D'Ariano,
Quantum information becomes classical when distributed to many users,
\doi{10.1103/PhysRevLett.97.250503}{\emph{Phys. Rev. Lett.} \textbf{97}, 250503 (2006)}.

\bibitem[CY14]{ChiribellaYang2014Optimal}
G. Chiribella and Y. Yang,
Optimal asymptotic cloning machines,
\doi{10.1088/1367-2630/16/6/063005}{\emph{New J. Phys.} \textbf{16}, 063005 (2014)},
\arxiv{1404.0990}.

\bibitem[KC22]{KimChitambar2022}
C. Kim and E. Chitambar,
Process-optimized phase-covariant quantum cloning,
\doi{10.1103/PhysRevA.106.022405}{\emph{Phys. Rev. A} \textbf{106}, 022405 (2022)},
\arxiv{2107.03042}.

\bibitem[CKNWW05]{CerfKrugerNavezWernerWolf2005}
N. J. Cerf, O. Kr\"uger, P. Navez, R. F. Werner, and M. M. Wolf,
Non-Gaussian cloning of quantum coherent states is optimal,
\doi{10.1103/PhysRevLett.95.070501}{\emph{Phys. Rev. Lett.} \textbf{95}, 070501 (2005)},
\arxiv{quant-ph/0410058}.

\bibitem[DG98]{DuanGuo1998}
L.-M. Duan and G.-C. Guo,
Probabilistic cloning and identification of linearly independent quantum states,
\doi{10.1103/PhysRevLett.80.4999}{\emph{Phys. Rev. Lett.} \textbf{80}, 4999--5002 (1998)}.

\bibitem[NG98]{NiuGriffiths1998}
C.-S. Niu and R. B. Griffiths,
Optimal copying of one quantum bit,
\doi{10.1103/PhysRevA.58.4377}{\emph{Phys. Rev. A} \textbf{58}, 4377--4393 (1998)},
\arxiv{quant-ph/9805063}.

\bibitem[BCDM00]{BrussCinchettiDArianoMacchiavello2000}
D. Bruss, M. Cinchetti, G. M. D'Ariano, and C. Macchiavello,
Phase-covariant quantum cloning,
\doi{10.1103/PhysRevA.62.012302}{\emph{Phys. Rev. A} \textbf{62}, 012302 (2000)},
\arxiv{quant-ph/9909046}.

\bibitem[Cer00]{Cerf2000Asymmetric}
N. J. Cerf,
Asymmetric quantum cloning in any dimension,
\doi{10.1080/09500340008244036}{\emph{J. Mod. Opt.} \textbf{47}, 187--209 (2000)},
\arxiv{quant-ph/9805024}.

\bibitem[MJPV99]{MuraoJonathanPlenioVedral1999}
M. Murao, D. Jonathan, M. B. Plenio, and V. Vedral,
Quantum telecloning and multiparticle entanglement,
\doi{10.1103/PhysRevA.59.156}{\emph{Phys. Rev. A} \textbf{59}, 156--161 (1999)},
\arxiv{quant-ph/9806082}.

\bibitem[CDP08]{ChiribellaDArianoPerinotti2008}
G. Chiribella, G. M. D'Ariano, and P. Perinotti,
Optimal cloning of unitary transformation,
\doi{10.1103/PhysRevLett.101.180504}{\emph{Phys. Rev. Lett.} \textbf{101}, 180504 (2008)},
\arxiv{0804.0129}.

\bibitem[CYY13]{ChiribellaYangYao2013}
G. Chiribella, Y. Yang, and A. C.-C. Yao,
Quantum replication at the Heisenberg limit,
\doi{10.1038/ncomms3915}{\emph{Nat. Commun.} \textbf{4}, 2915 (2013)},
\arxiv{1304.2910}.

\bibitem[GCMBC14]{GendraCalsamigliaMunozTapiaBaganChiribella2014}
B. Gendra, J. Calsamiglia, R. Mu\~noz-Tapia, E. Bagan, and G. Chiribella,
Probabilistic metrology attains macroscopic cloning of quantum clocks,
\doi{10.1103/PhysRevLett.113.260402}{\emph{Phys. Rev. Lett.} \textbf{113}, 260402 (2014)},
\arxiv{1406.2218}.

\bibitem[DSS15]{DurSekatskiSkotiniotis2015}
W. D\"ur, P. Sekatski, and M. Skotiniotis,
Deterministic superreplication of one-parameter unitary transformations,
\doi{10.1103/PhysRevLett.114.120503}{\emph{Phys. Rev. Lett.} \textbf{114}, 120503 (2015)},
\arxiv{1410.6008}.

\bibitem[CYH15]{ChiribellaYangHuang2015}
G. Chiribella, Y. Yang, and C. Huang,
Universal superreplication of unitary gates,
\doi{10.1103/PhysRevLett.114.120504}{\emph{Phys. Rev. Lett.} \textbf{114}, 120504 (2015)}.

\bibitem[CY16]{ChiribellaYang2016}
G. Chiribella and Y. Yang,
Quantum superreplication of states and gates,
\doi{10.1007/s11467-016-0556-7}{\emph{Front. Phys.} \textbf{11}, 110304 (2016)},
\arxiv{1512.05839}.

\bibitem[SSD15]{SekatskiSkotiniotisDur2015NoSignaling}
P. Sekatski, M. Skotiniotis, and W. D\"ur,
No-signaling bounds for quantum cloning and metrology,
\doi{10.1103/PhysRevA.92.022355}{\emph{Phys. Rev. A} \textbf{92}, 022355 (2015)},
\arxiv{1412.3361}.

\bibitem[HT24]{HuTomamichel2024}
Y. Hu and M. Tomamichel,
Fundamental limits on quantum cloning from the no-signaling principle,
\doi{10.1103/PhysRevA.109.022221}{\emph{Phys. Rev. A} \textbf{109}, 022221 (2024)},
\arxiv{2305.02002}.

\bibitem[SGKS25]{SekatskiGuryanovaKothakondaSkotiniotis2025}
P. Sekatski, Y. Guryanova, N. B. T. Kothakonda, and M. Skotiniotis,
Cloning quantum channels,
\arxiv{2509.08059} (2025).

\bibitem[BGSQ26]{BrzicGrinkoStudzinskiQuintino2026}
V. Brzi\'c, D. Grinko, M. Studzi\'nski, and M. T. Quintino,
Optimal pure state cloning and transposition are complementary channels,
\arxiv{2603.23628} (2026).

\bibitem[SIGA05]{ScaraniIblisdirGisinAcin2005}
V. Scarani, S. Iblisdir, N. Gisin, and A. Acin,
Quantum cloning,
\doi{10.1103/RevModPhys.77.1225}{\emph{Rev. Mod. Phys.} \textbf{77}, 1225--1256 (2005)},
\arxiv{quant-ph/0511088}.

\bibitem[FWJ+14]{FanWangJingYueShiZhangMu2014}
H. Fan, Y.-N. Wang, L. Jing, J.-D. Yue, H.-D. Shi, Y.-L. Zhang, and L.-Z. Mu,
Quantum cloning machines and the applications,
\doi{10.1016/j.physrep.2014.06.004}{\emph{Phys. Rep.} \textbf{544}, 241--322 (2014)},
\arxiv{1301.2956}.

\bibitem[BCFJS96]{BarnumCavesFuchsJozsaSchumacher1996}
H. Barnum, C. M. Caves, C. A. Fuchs, R. Jozsa, and B. Schumacher,
Noncommuting mixed states cannot be broadcast,
\doi{10.1103/PhysRevLett.76.2818}{\emph{Phys. Rev. Lett.} \textbf{76}, 2818--2821 (1996)},
\arxiv{quant-ph/9511010}.

\bibitem[KH08]{KalevHen2008}
A. Kalev and I. Hen,
No-broadcasting theorem and its classical counterpart,
\doi{10.1103/PhysRevLett.100.210502}{\emph{Phys. Rev. Lett.} \textbf{100},
210502 (2008)}.

\bibitem[CEM99]{CiracEkertMacchiavello1999}
J. I. Cirac, A. K. Ekert, and C. Macchiavello,
Optimal purification of single qubits,
\doi{10.1103/PhysRevLett.82.4344}{\emph{Phys. Rev. Lett.} \textbf{82}, 4344--4347 (1999)},
\arxiv{quant-ph/9812075}.

\bibitem[Ras03a]{Rastegin2003Upper}
A. E. Rastegin,
Upper bound on the global fidelity for mixed-state cloning,
\doi{10.1103/PhysRevA.67.012305}{\emph{Phys. Rev. A} \textbf{67}, 012305 (2003)}.

\bibitem[Ras03b]{Rastegin2002Relative}
A. E. Rastegin,
A lower bound on the relative error of mixed-state cloning and related operations,
\emph{J. Opt. B: Quantum Semiclass. Opt.} \textbf{5}, S647--S650 (2003),
\arxiv{quant-ph/0208159}.

\bibitem[Ras03c]{Rastegin2003StateDependent}
A. E. Rastegin,
Global-fidelity limits of state-dependent cloning of mixed states,
\doi{10.1103/PhysRevA.68.032303}{\emph{Phys. Rev. A} \textbf{68}, 032303 (2003)},
\arxiv{quant-ph/0301132}.

\bibitem[Fan03]{Fan2003}
H. Fan,
Quantum cloning of mixed states in symmetric subspace,
\doi{10.1103/PhysRevA.68.054301}{\emph{Phys. Rev. A} \textbf{68}, 054301 (2003)},
\arxiv{quant-ph/0308058}.

\bibitem[FLS07]{FanLiuShi2007}
H. Fan, B. Y. Liu, and K. J. Shi,
Quantum cloning of identical mixed qubits,
\doi{10.26421/QIC7.5-6-8}{\emph{Quantum Inf. Comput.} \textbf{7}, 551--558 (2007)},
\arxiv{quant-ph/0601017}.

\bibitem[Kaz10]{Kazakov2007}
A. Ya. Kazakov,
``Partial'' quantum cloning and quantum cloning of mixed states,
\doi{10.1142/S0219749910006186}{\emph{Int. J. Quantum Inf.}
\textbf{8}, 435--442 (2010)},
\arxiv{0711.2860}.

\bibitem[CC07]{ChenChen2007MixedBroadcast}
L. Chen and Y.-X. Chen,
Mixed qubit cannot be universally broadcast,
\doi{10.1103/PhysRevA.75.062322}{\emph{Phys. Rev. A} \textbf{75}, 062322 (2007)},
\arxiv{quant-ph/0612205}.

\bibitem[LW17]{LemmWilde2017}
M. Lemm and M. M. Wilde,
Information-theoretic limitations on approximate quantum cloning and broadcasting,
\doi{10.1103/PhysRevA.96.012304}{\emph{Phys. Rev. A} \textbf{96}, 012304 (2017)},
\arxiv{1608.07569}.

\bibitem[DMP05]{DArianoMacchiavelloPerinotti2005}
G. M. D'Ariano, C. Macchiavello, and P. Perinotti,
Superbroadcasting of mixed states,
\doi{10.1103/PhysRevLett.95.060503}{\emph{Phys. Rev. Lett.} \textbf{95}, 060503 (2005)},
\arxiv{quant-ph/0506251}.

\bibitem[BDMP05]{BuscemiDArianoMacchiavelloPerinotti2005Optimal}
F. Buscemi, G. M. D'Ariano, C. Macchiavello, and P. Perinotti,
Optimal superbroadcasting of mixed qubit states,
\arxiv{quant-ph/0510155} (2005).

\bibitem[BDMP06]{BuscemiDArianoMacchiavelloPerinotti2006}
F. Buscemi, G. M. D'Ariano, C. Macchiavello, and P. Perinotti,
Universal and phase-covariant superbroadcasting for mixed qubit states,
\doi{10.1103/PhysRevA.74.042309}{\emph{Phys. Rev. A} \textbf{74}, 042309 (2006)},
\arxiv{quant-ph/0602125}.

\bibitem[DF07]{DangFan2007}
G.-F. Dang and H. Fan,
Optimal broadcasting of mixed states,
\doi{10.1103/PhysRevA.76.022323}{\emph{Phys. Rev. A} \textbf{76}, 022323 (2007)},
\arxiv{quant-ph/0609182}.

\bibitem[DPS06]{DArianoPerinottiSacchi2006Continuous}
G. M. D'Ariano, P. Perinotti, and M. F. Sacchi,
Superbroadcasting of continuous variable mixed states,
\doi{10.1088/1367-2630/8/6/099}{\emph{New J. Phys.} \textbf{8}, 99 (2006)},
\arxiv{quant-ph/0601114}.

\bibitem[PFBC24]{ParzygnatFullwoodBuscemiChiribella2024}
A. J. Parzygnat, J. Fullwood, F. Buscemi, and G. Chiribella,
Virtual quantum broadcasting,
\doi{10.1103/PhysRevLett.132.110203}{\emph{Phys. Rev. Lett.} \textbf{132}, 110203 (2024)},
\arxiv{2310.13049}.

\bibitem[BGA11]{BowlesGutaAdesso2011Purification}
P. Bowles, M. Gu{\c t}{\u a}, and G. Adesso,
Asymptotically optimal purification and dilution of mixed qubit and Gaussian states,
\doi{10.1103/PhysRevA.84.022320}{\emph{Phys. Rev. A} \textbf{84}, 022320 (2011)},
\arxiv{1102.2341}.

\bibitem[BGA12]{BowlesGutaAdesso2012Reversal}
P. Bowles, M. Gu{\c t}{\u a}, and G. Adesso,
Asymptotically optimal quantum channel reversal for qudit ensembles and multimode Gaussian states,
\doi{10.1088/1367-2630/14/11/113041}{\emph{New J. Phys.} \textbf{14}, 113041 (2012)},
\arxiv{1208.3569}.

\bibitem[Cav82]{Caves1982LinearAmplifiers}
C. M. Caves,
Quantum limits on noise in linear amplifiers,
\doi{10.1103/PhysRevD.26.1817}{\emph{Phys. Rev. D} \textbf{26}, 1817--1839 (1982)}.

\bibitem[CIR00]{CerfIpeRottenberg2000}
N. J. Cerf, A. Ipe, and X. Rottenberg,
Cloning of continuous quantum variables,
\doi{10.1103/PhysRevLett.85.1754}{\emph{Phys. Rev. Lett.} \textbf{85}, 1754--1757 (2000)},
\arxiv{quant-ph/9909037}.

\bibitem[BCI+01]{BraunsteinCerfIblisdirVanLoockMassar2001}
S. L. Braunstein, N. J. Cerf, S. Iblisdir, P. van Loock, and S. Massar,
Optimal cloning of coherent states with a linear amplifier and beam splitters,
\doi{10.1103/PhysRevLett.86.4938}{\emph{Phys. Rev. Lett.} \textbf{86}, 4938--4941 (2001)},
\arxiv{quant-ph/0012046}.

\bibitem[GJ07]{GutaJencova2007LAN}
M. Gu{\c t}{\u a} and A. Jen{\v c}ov\'a,
Local asymptotic normality in quantum statistics,
\doi{10.1007/s00220-007-0340-1}{\emph{Commun. Math. Phys.} \textbf{276}, 341--379 (2007)},
\arxiv{quant-ph/0606213}.

\bibitem[GK06]{GutaKahn2006LANQubits}
M. Gu{\c t}{\u a} and J. Kahn,
Local asymptotic normality for qubit states,
\doi{10.1103/PhysRevA.73.052108}{\emph{Phys. Rev. A} \textbf{73}, 052108 (2006)},
\arxiv{quant-ph/0512075}.

\bibitem[KG09]{KahnGuta2009LANFiniteDimensional}
J. Kahn and M. Gu{\c t}{\u a},
Local asymptotic normality for finite dimensional quantum systems,
\doi{10.1007/s00220-009-0787-3}{\emph{Commun. Math. Phys.} \textbf{289}, 597--652 (2009)},
\arxiv{0804.3876}.

\bibitem[YFG13]{YamagataFujiwaraGill2013LAN}
K. Yamagata, A. Fujiwara, and R. D. Gill,
Quantum local asymptotic normality based on a new quantum likelihood ratio,
\doi{10.1214/13-AOS1147}{\emph{Ann. Statist.} \textbf{41}, 2197--2217 (2013)},
\arxiv{1210.3749}.

\bibitem[GM06]{GutaMatsumoto2006MixedGaussianCloning}
M. Gu{\c t}{\u a} and K. Matsumoto,
Optimal cloning of mixed Gaussian states,
\doi{10.1103/PhysRevA.74.032305}{\emph{Phys. Rev. A} \textbf{74}, 032305 (2006)},
\arxiv{quant-ph/0605161}.

\bibitem[ZC17]{ZhaoChiribella2017}
X. Zhao and G. Chiribella,
Quantum amplification and purification of noisy coherent states,
\doi{10.1103/PhysRevA.95.042303}{\emph{Phys. Rev. A} \textbf{95}, 042303 (2017)},
\arxiv{1701.04602}.

\bibitem[LN24]{LahiryNussbaum2024}
S. Lahiry and M. Nussbaum,
Minimax estimation of low-rank quantum states and their linear functionals,
\doi{10.3150/23-BEJ1610}{\emph{Bernoulli} \textbf{30}, 610--635 (2024)},
\arxiv{2111.03279}.

\bibitem[KW01]{KeylWerner2001Spectrum}
M. Keyl and R. F. Werner,
Estimating the spectrum of a density operator,
\doi{10.1103/PhysRevA.64.052311}{\emph{Phys. Rev. A} \textbf{64}, 052311 (2001)},
\arxiv{quant-ph/0102027}.

\bibitem[PRV67]{ParthasarathyRangaRaoVaradarajan1967}
K. R. Parthasarathy, R. Ranga Rao, and V. S. Varadarajan,
Representations of complex semi-simple Lie groups and Lie algebras,
\doi{10.2307/1970351}{\emph{Annals of Mathematics} \textbf{85}, 383--429 (1967)}.

\bibitem[FH91]{FultonHarris1991}
W. Fulton and J. Harris,
\emph{Representation Theory: A First Course},
Graduate Texts in Mathematics 129, Springer (1991).

\bibitem[TWZ25]{TangWrightZhandry2025}
E. Tang, J. Wright, and M. Zhandry,
Conjugate queries can help,
\arxiv{2510.07622} (2025).

\bibitem[LTHC26]{LiTheilHarrowChuang2026}
Z. Li, E. Theil, A. W. Harrow, and I. Chuang,
An exponential sample-complexity advantage for coherent quantum inference,
\arxiv{2605.21457} (2026).

\bibitem[LTHC26b]{LiTheilHarrowChuangQPA2026}
Z. Li, E. Theil, A. W. Harrow, and I. Chuang,
Quantum purity amplification for arbitrary eigenstates and multiple outputs,
\arxiv{2605.21570} (2026).

\end{thebibliography}

\end{document}
